%%%%%%%%%%%%%%%%%%%%%%%%%%%%%%%%%%%%%%%%%%%%%%%%%%%%%%%%%%%%%%%%%%%%%%
% Lớp 11 | Dãy số. Cấp số cộng. Cấp số nhân | Cấp số cộng
% Dạng 1D2?2-7: Toán thực tế áp dụng cấp số cộng
% Số câu: 91  (sinh tự động bởi tach_cau.py; từ nay thầy sửa trực tiếp file này)
%%%%%%%%%%%%%%%%%%%%%%%%%%%%%%%%%%%%%%%%%%%%%%%%%%%%%%%%%%%%%%%%%%%%%%

%%%%%%%%%%%%%%%% TRẮC NGHIỆM 4 PHƯƠNG ÁN %%%%%%%%%%%%%%%%

% nguồn: Ok_Toan11_TracNghiem_Lan2
\begin{ex}%[1D2H2-7]%[ma:a7cddca7]
	Người  ta trồng $820$ cây theo một hình tam giác như sau: Hàng thứ nhất trồng $1$ cây, kể từ hàng thứ hai trở đi số cây trồng mỗi hàng nhiều hơn $1$ cây so với hàng liền trước nó. Hỏi có tất cả bao nhiêu hàng cây?
	\choice 
		{$42$} 
		{$41$} 
		{\True $40$} 
		{$39$} 
	\loigiai{Giả sử trồng được  $n$ hàng cây $(n\ge 1$, $n\in \mathbb{N}$).\\
		Số cây ở mỗi  hàng lập thành cấp số cộng có $u_1=1$ và công sai $d=1$.\\
		 Theo giả thiết
		 \allowdisplaybreaks
		\begin{eqnarray*}
			 &&S_n=820\\
			&\Leftrightarrow&\dfrac {n}{2}\left[2u_1+(n-1)d \right]=820\\ 
			&\Leftrightarrow&n(n+1)=1\,640\\ 
			&\Leftrightarrow&n^2+n-1\,640=0\\ 
			 &\Leftrightarrow&\hoac{&n=40\\&n=-41.}
		\end{eqnarray*}
		So với điều kiện, suy ra $n=40$.  Vậy có tất cả $40$ hàng cây.
	 }
\end{ex}

% nguồn: Ok_Toan11_TracNghiem_Lan2
\begin{ex}%[1D2H2-7]%[ma:f1c44571]
	Hùng đang tiết kiệm để mua một cây đàn piano có giá $142$ triệu đồng. Trong tháng đầu tiên, anh ta để dành được $20$ triệu đồng. Mỗi tháng tiếp theo anh ta để dành được $3$ triệu đồng và đưa số tiền tiết kiệm của mình. Hỏi ít nhất vào tháng thứ bao nhiêu  thì Hùng mới có đủ tiền để mua cây đàn piano đó?
	 \choice
		 {$43$}
		 {$41$}
		 {$40$}
		{\True  $42$}
	\loigiai{Tổng số tiền  Hùng tiết kiệm được vào mỗi tháng lập thành một cấp số cộng $(u_n)$ có số hạng đầu $u_1=20$ và công sai $d=3$.\\
		Tổng số tiền Hùng tiết kiệm được  vào tháng thứ $n$ bằng
		\[u_n=u_1+(n-1)d=20+(n-1)\cdot  3=3n+17.\]
		Hùng có đủ tiền mua cây đàn  $3n+17\ge 142$ $\Leftrightarrow n\ge \dfrac {125}{3}\approx 41{,}67$.\\
		Vậy ít nhất vào tháng thứ $42$ thì  Hùng mới có đủ tiền để mua cây đàn piano đó.
	} 
\end{ex}

% nguồn: Ok_Toan11_TracNghiem_Lan2
\begin{ex}%[1D2H2-7]%[ma:49d78c8f]
	Lan đang tiết kiệm để mua laptop. Trong tuần đầu tiên, cô ta  để dành $200$ đô la, và trong mỗi tuần tiếp theo, cô ta đã thêm $16$ đô la vào tài khoản tiết kiệm của mình. Chiếc laptop Lan cần mua có giá $1\,000$ đô la. Hỏi vào tuần thứ bao nhiêu thì cô ấy có đủ tiền để mua chiếc laptop đó?
	 \choice
		 {$49$}
		 {$50$}
		 {\True $51$}
		{$52$} 
	\loigiai{Gọi $n$ là số tuần cô ta đã thêm $16$ đô la vào tài khoản  tiết kiệm của mình.\\
		 Số tiền cô ta tiết kiệm được sau $n$ tuần đó là $T=200+16n$.\\
		Theo đề bài,  ta có $T=200+16n=1\,000\Leftrightarrow n=50$.\\
		Vậy kể cả tuần đầu thì tuần thứ $51$ cô ta có  đủ tiền để mua chiếc laptop đó.
	} 
\end{ex}

% nguồn: Ok_Toan11_TracNghiem_Lan2
\begin{ex}%[1D2H2-7]%[ma:0e0c75ce]
	Trong tháng $12$, lớp 11A dự kiến quyên góp tiền để đi làm từ thiện như sau: Ngày đầu quyên góp được mỗi bạn bỏ $2\,000$ đồng vào lợn, từ ngày thứ hai trở đi mỗi  bạn bỏ vào lợn hơn ngày liền trước là $500$ đồng. Hỏi sau $28$ ngày lớp 11A quyên góp được bao nhiêu tiền? Biết lớp có $40$ bạn.
	 \choice
		{$8\,800\,000$ đồng} 
		{\True  $9\,800\,000$ đồng}
		{$10\,800\,000$  đồng}
		{$7\,800\,000$  đồng}
	\loigiai{Số tiền mỗi bạn lớp 11A quyên góp để làm từ thiện lập thành một cấp số cộng có  số hạng đầu $u_1=2\,000$, công sai $d=500$.\\
		Vậy sau $28$ ngày số tiền mỗi bạn quyên góp  được là $28\cdot 2\,000+\dfrac {27\cdot 28\cdot 500}{2}=245\,000$ đồng.\\
		Vậy sau $28$ ngày tổng số tiền quyên góp được của  lớp 11A là $245\,000\cdot 40=9\,800\,000$ đồng.
	 }
\end{ex}

% nguồn: Ok_Toan11_TracNghiem_Lan2
\begin{ex}%[1D2H2-7]%[ma:171b4d8d]
	Một người làm việc cho một công ty. Theo hợp đồng trong năm đầu tiên, tháng lương thứ nhất là $6$ triệu đồng và lương tháng sau cao hơn tháng trước là $200$ ngàn đồng. Hỏi theo hợp đồng, tháng thứ $7$ người đó nhận được lương là bao  nhiêu?
	\choice 
		{$7$ triệu} 
		{$7{,}3$ triệu} 
		{\True $7{,}2$ triệu} 
		{$7{,}4$  triệu}
	 \loigiai{Gọi lương tháng thứ $n$ của người đó là $A_n$.\\
		Ta có $A_1=6$.\\ 
		Lương tháng sau hơn tháng trước $0{,}2$ triệu nên ta có $\left(A_n\right)$ là một cấp số cộng với số  hạng đầu $A_1=6$ và công sai $d=0{,}2$.\\
		Số hạng tổng quát của dãy số là $A_n=A_1+(n-1)d$, với  $n>1$.\\
		Vậy tới tháng thứ  $7$, người đó nhận được lương là $A_7=A_1+6d=6+6\cdot 0{,}2=7{,}2$.
	 }
\end{ex}

% nguồn: Ok_TracNghiemLop11_Lan1
\begin{ex}%[1D2H2-7]%[ma:82f71f6f]
	Một thửa  ruộng bậc thang có thửa thấp nhất (bậc thấp nhất) nằm ở độ cao $900$ mét so với mực nước biển và độ chênh lệch giữa thửa trên và thửa dưới (hai thửa liên tiếp) trung bình là $1{,}5$ mét. Hỏi bậc thứ $19$ của thửa ruộng đó có độ cao là bao nhiêu so với mực nước biển?
	 \choice
		{$930$ mét} 
		{$928{,}5$  mét}
		{$925{,}5$  mét}
		 {\True $927$ mét}
	 \loigiai{
		Gọi $u_n$ là chiều cao (mét) so với mực nước biển của thửa ruộng bậc thang ở  bậc thứ $n$.\\
		Khi đó, $(u_n)$ là  một cấp số cộng với $u_1=900 \text{m}$ và $d=1{,}5 \text{m}$.\\
		 Ta có: $u_{19}=u_1+18 \cdot d=900+18 \cdot 1{,}5=927$ mét.\\
		Vậy bậc thứ  $19$ của thửa ruộng có độ cao là $927$ mét so với mực nước biển.
	 }
\end{ex}

% nguồn: Ok_TracNghiemLop11_Lan1
\begin{ex}%[1D2H2-7]%[ma:9fd79bf3]
	Trong dịp tết trồng cây, một trường THPT muốn  trồng $50$ hàng cây theo quy cách như sau:
	\begin{itemize}
		\item Hàng 1: Trồng  $2$ cây.
		 \item Hàng 2: Trồng $4$ cây.
		\item Hàng 3:  Trồng $6$ cây.\\
		$\cdots \cdots$ 
	\end{itemize}
	 Hỏi nhà trường trồng được bao nhiêu cây?
	\choice 
		{\True $ 2550$} 
		{$ 2650$} 
		{$ 5300$} 
		{$  2500$}
	 \loigiai{
	Ta có  $u_1=2$, $d=2$ nên $S_{50}=50\cdot2+\dfrac{50\cdot \left(50-1\right).2}{2}=2550$.}
\end{ex}

% nguồn: Ok_Toan11_TracNghiem_Lan2
\begin{ex}%[1D2V2-7]%[ma:255cecbe]
	Công ty A muốn thuê hai mảnh đất để làm $2$ nhà kho, một mảnh trong vòng $10$ năm và $1$ mảnh trong vòng  $15$ năm ở hai chỗ khác nhau. Công ty bất động sản C, công ty bất động sản B đều muốn cho thuê. Hai công ty đưa ra phương án cho thuê như sau:
	 Công ty C: Năm đầu tiên tiền thuê đất là $60$ triệu và kể từ năm thứ hai trở đi mỗi năm tăng thêm $3$ triệu đồng.
	Công ty B: Trả tiền theo quí, quý đầu tiên là $8$ triệu đồng và  từ quý thứ hai trở đi mỗi quý tăng thêm $500\,000$ đồng.
	Hỏi công ty A nên lựa chọn thuê đất của công ty bất động sản nào để chi phí là thấp nhất biết rằng các mảnh đất cho thuê về diện tích,  độ tiện lợi đều như nhau?
	\choice 
		{Chọn công  ty B để thuê cả hai mảnh đất}
		 {\True Chọn công ty C để thuê cả hai mảnh đất}
		{Chọn công ty C để thuê đất $10$ năm, công ty B thuê đất $15$ năm} 
		{Chọn công ty B để thuê đất $10$ năm, công  ty C thuê đất $15$ năm}
	\loigiai{Gọi $B_n$, $C_n$ lần lượt là số tiền công  ty A cần trả theo các tính của hai công ty B và C.\\
		 Theo bài ra ta có\\
		$B_n$ là tổng $n$ số hạng đầu tiên của một cấp số cộng với $u_1=8$ triệu  đồng $d=0{,}5$ triệu đồng.\\
		$C_n$ là tổng $n$ số hạng đầu tiên của một  cấp số cộng với $u_1=60$ triệu đồng $d=3$ triệu đồng.\\
		Do đó: Nếu thuê đất của công ty B trong vòng $15$ năm $=60$ quý số tiền công ty A phải  trả là\\
		$B_{60}=(2\cdot  8+59\cdot 0{,}5)\cdot 30=1\,365$ triệu đồng.\\
		Nếu thuê đất của công ty C  trong vòng $15$ năm số tiền công ty A phải trả là\\
		$C_{15}=(2\cdot 60+14\cdot 3)\cdot 7{,}5=1\,215$ triệu đồng.\\ 
		Vậy thuê mảnh đất trong vòng $15$ năm của công  ty C.\\
		Nếu thuê đất của công ty $B$ trong vòng $10$ năm $=40$ quý số tiền công ty A phải trả là  $B_{40}=(2\cdot 8+39\cdot 0{,}5)\cdot 20=710$ triệu đồng.\\
		Nếu thuê đất của công ty C trong vòng $10$ năm số tiền công ty A phải trả là $C_{10}=(2\cdot 60+9\cdot 3)\cdot 4{,}5=661{,}5$  triệu đồng.\\
		Vậy thuê mảnh đất trong vòng $10$ năm của công  ty C.
	} 
\end{ex}

% nguồn: Ok_Toan11_TracNghiem_Lan2
\begin{ex}%[1D2V2-7]%[ma:4bb2540a]
	Một cầu thang đường lên cổng trời của một điểm giải trí ở công viên tỉnh X được hàn bằng sắt có  hình dáng các bậc thang đều là hình chữ nhật với cùng chiều rộng là $35$ cm và chiều dài của nó theo thứ tự mỗi bậc đều giảm dần đi $7$ cm. Biết rằng bậc đầu tiên của cầu thang là hình chữ nhật có chiều dài $189$ cm và bậc cuối cùng cầu thang là hình chữ nhật có chiều dài $63$ cm. Hỏi giá thành làm cầu thang đó gần với số nào dưới đây nếu giá thành làm một mét vuông cầu thang đó là $1\,250\,000$ đồng trên một mét vuông?
	\choice 
		 {$9\,500\,000$ đồng}
		{$11\,000\,000$  đồng}
		{$10\,000\,000$  đồng}
		{\True $10\,500\,000$ đồng} 
	\loigiai{Ta có chiều dài của mỗi mặt cầu thang theo thứ tự lập thành một cấp số cộng với số hạng đầu tiên là $u_1=189$, công sai $d=-7$ và số hạng cuối  cùng là $u_n=63$.\\
		Khi đó áp dụng công thức tính số hạng tổng quát  ta có\\
		$u_n=u_1+(n-1)d\Leftrightarrow 63=189-7(n-1)\Leftrightarrow  n=19$.\\
		Tổng chiều dài của $19$  hình chữ nhật đó là $S_{19}=19\cdot \dfrac {u_1+u_{19}}{2}=2\,394$.\\
		Diện tích của $19$ bậc thang là $S=35\cdot 2\,394=83\,790$ (cm$^2$) $=8{,}379$  (m$^2$).\\
		Tổng số tiền để  làm cầu thang đó là $T=8{,}379\cdot 1\,250\,000=10\,473\,750$ đồng.
	 }
\end{ex}


%%%%%%%%%%%%%%%% ĐÚNG/SAI %%%%%%%%%%%%%%%%

% nguồn: Ok_dungsai_lop11_Lan2
\begin{ex}%[1D2N2-7]%[ma:2fc30629]
	Hùng đang tiết kiệm để mua một cây guitar. Trong tuần đầu tiên, anh ta để dành $42$ đô la, và trong mỗi tuần tiết theo, anh ta đã thêm $8$ đô  la vào tài khoản tiết kiệm của mình. Cây guitar Hùng cần mua có giá $400$ đô la. Hỏi anh ấy cần tối thiểu bao nhiêu tuần để có đủ tiền để mua cây guitar đó?
	 \choiceTF
		 {$47$}
		 {$45$}
		 {$44$}
		{\True  $46$}
	 \loigiai{
		Gọi $n$ là số tuần anh ta đã thêm $8$ đô la vào tài khoản tiết kiệm của  mình.\\
		Số tiền anh ta tiết kiệm được sau $n$ tuần đó  là $S=42+8n$.\\
		Theo bài ra $S=42+8n\ge 400\Leftrightarrow n\ge 44{,}75\Rightarrow n=45$.\\ 
		Vậy kể cả tuần đầu thì tuần thứ $46$ anh ta có đủ tiền để mua cây guitar  đó. Do đó
		\begin{itemchoice}
			\itemch  {\bf Sai}.\\
			\itemch {\bf Sai}.\\ 
			 \itemch {\bf Sai}.\\
			\itemch {\bf Đúng}.\\ 
		\end{itemchoice}
	} 
\end{ex}

% nguồn: Ok_dungsai_lop11_Lan2
\begin{ex}%[1D2H2-7]%[ma:be82c5f1]
	Trong hội chợ tết, một  công ty sữa muốn xếp $10\,000$ hộp sữa theo thứ tự từ trên xuống dưới như sau Hàng thứ nhất có $1$ hộp sữa, hàng thứ hai có $3$ hộp sữa, hàng thứ ba có $5$ hộp sữa, \ldots\, cứ như thế, số lượng hộp sữa của hàng sau lớn hơn số lượng hộp sữa của hàng trước nó là $2$ hộp sữa.
	\choiceTF 
		{\True Gọi $u_n$ là số hộp sữa ở hàng thứ $n$ thì $(u_n)$ là một cấp số cộng có số hạng đầu $u_1=1$  và công sai $d=2$}
		{Số hộp sữa của hàng  thứ $10$ là $20$ hộp sữa}
		{\True Để xếp được $20$  hàng thì cần $400$ hộp sữa}
		{Hàng cuối cùng có $198$ hộp sữa} 
	\loigiai{ 
		Dãy số hộp sữa ở mỗi hàng $(u_n)$ là một cấp số cộng có số hạng đầu $u_1=1$ và công sai $d=2$.\\ 
		Số hộp  sữa của hàng thứ $10$ là $u_{10} = u_1+9d = 1+9\cdot 2=19$ hộp sữa.\\
		Để xếp được $20$ hàng thì số hộp sữa cần có là $S_{20} = \dfrac{20}{2}(2u_1+19\cdot d) =  \dfrac{20}{2}(2\cdot 1+19\cdot 2) = 400$ hộp sữa.\\
		Ta  có
		\begin{eqnarray*}
			 && S_n=\dfrac{n}{2}\left[2u_1+(n-1)d\right]\\
			&\Leftrightarrow&  10\,000=\dfrac{n}{2}\left[2\cdot 1+(n-1)\cdot2\right]\\
			&\Leftrightarrow& 10\,000=n^2\\ 
			 &\Leftrightarrow& \hoac{&n=100~ \text{(nhận)}\\&n=-100~\text{(loại).}}
		\end{eqnarray*}
		Vậy  xếp được tất cả $100$ hàng. Nên $u_{100}=1+(100-1)\cdot 2=199$.
		\begin{itemchoice}
			\itemch  Dãy số hộp sữa ở mỗi hàng là một cấp số cộng có số hạng đầu $u_1=1$ và công sai $d=2$.
			\itemch Số hộp sữa của hàng thứ $10$ là $19$ hộp  sữa.
			\itemch Để xếp  được $20$ hàng thì cần $400$ hộp sữa.
			\itemch  Hàng cuối cùng có $199$ hộp sữa.
		\end{itemchoice}
	} 
\end{ex}

% nguồn: Ok_dungsai_lop11_Lan2
\begin{ex}%[1D2H2-7]%[ma:d8b1637e]
	Trên một bàn cờ có nhiều ô vuông, người ta đặt $8$ hạt dẻ vào ô đầu tiên, sau đó đặt tiếp vào ô thứ hai số hạt nhiều hơn ô thứ nhất là $3$, tiếp tục đặt vào ô thứ ba số hạt nhiều hơn ô thứ hai là $3$, \ldots và cứ thế  tiếp tục đến ô thứ $n$. Biết rằng đặt hết số ô trên bàn cờ người ta phải sử dụng $18\,865$ hạt ?
	\choiceTF 
		{\True Số  hạt dẻ ở mỗi ô trên bàn cờ lập thành một cấp số cộng với $u_1=8$ và $d=3$}
		{\True Số hạt dẻ ở ô thứ  $5$ là $20$}
		{Bàn cờ đó  có $111$ ô vuông}
		 {Tổng số hạt dẻ ở $10$ ô cuối là $215$}
	\loigiai{ 
		\begin{itemchoice}
			\itemch {\bf Đúng}.\\ Dễ thấy số hạt dẻ ở mỗi ô trên  bàn cờ lập thành một cấp số cộng với $u_1=8$ và $d=3$.
			\itemch {\bf Đúng}.\\ Số hạt dẻ ở ô thứ $5$ chính là số hạng thứ $5$  của dãy mà $u_5=u_1+4d=8+4\cdot 3=20$.
			\itemch {\bf Sai}.\\  Vì theo đề bài ta có $S_n=18\,865$, với $n\in{\mathbb{N}^{*}}$.\\
			Mặt  khác $S_n=\dfrac{n}{2}\big[2u_1+(n-1)d\big]=\dfrac{n}{2}\big[2\cdot 8+(n-1)\cdot 3\big]=\dfrac{n}{2}(13+3n)$.\\
			Như vậy 
			\begin{eqnarray*}
				&&S_n=18\,865\\&\Leftrightarrow& \dfrac{n}{2}(13+3n)=18\,865\\ 
				 &\Leftrightarrow& 3n^2+13n-37\,730=0\\
				&\Leftrightarrow& \hoac{&n=110\\&n=-\dfrac{343}{3}.} 
			\end{eqnarray*}
			 \itemch {\bf Sai}.\\ Do bàn cờ có $110$ ô vuông (theo câu c) nên tổng $10$ ô cuối sẽ là tổng của $10$ số hạng từ $u_{101}$ đến $u_{110}$. Mà $u_{101}=u_1+100d=308$, $u_{110}=u_1+109d=335$. \\
			Vậy tổng số hạt  dẻ của $10$ ô cuối là $S=\dfrac{\left(u_{101}+u_{120}\right)\cdot 10}{2}=3\,215$.
		\end{itemchoice}
	} 
\end{ex}

% nguồn: Ok_dungsai_lop11_Lan2
\begin{ex}%[1D2H2-7]%[ma:76dc3254]
	Một sinh viên sau khi ra trường và xin vào làm cho một trung tâm với mức lương khởi điểm là $100$  triệu đồng một năm. Cứ sau mỗi năm, trung tâm trả thêm cho sinh viên $20$ triệu đồng. Gọi $u_n$ (triệu đồng) là số tiền lương mà sinh viên đó nhận được ở năm thứ $n$. Hỏi các mệnh đề sau đúng hay sai?
	 \choiceTF
		{\True Số tiền lương sinh  viên nhận được ở năm thứ hai là $120$ triệu đồng}
		{Số tiền lương sinh viên nhận được  ở năm thứ $10$ là $300$ triệu đồng}
		 {Dãy số $\left(u_n \right)$ là cấp số cộng có $u_1=120$ và công sai $d=20$}
		{Giả sử, mỗi năm bạn sinh viên chi tiêu tiết kiệm  hết $70$ triệu đồng. Vậy sau ít nhất $12$ năm thì sinh viên đó mua được căn chung cư $2$ tỉ đồng}
	\loigiai{ 
		Ta thấy, số tiền lương năm sau hơn năm trước $20$ triệu đồng  nên $\left(u_n \right)$ là cấp số cộng có $u_1=100$ và công sai $d=20$. Do đó $u_n=u_1+\left(n-1\right)d=100+\left(n-1\right)\cdot 20=20n+80$.
		\begin{itemchoice}
			\itemch \textbf{Đúng.} \\ 
			Số tiền lương sinh viên nhận được ở năm thứ hai là $u_2=120$ (triệu  đồng).
			\itemch \textbf{Sai.} \\ 
			Số tiền lương  sinh viên nhận được ở năm thứ $10$ là $u_{10}=20\cdot 10+80=280$ (triệu đồng).
			\itemch \textbf{Sai.} \\ 
			Vì $u_1=100$. 
			\itemch \textbf{Sai.}  \\
			Tổng số tiền bạn sinh viên tiết  kiệm được sau $n$ năm là \\
			$S=\dfrac{n}{2} \left[2u_1+\left(n-1\right)d\right]-70n=\dfrac{n}{2} \left[2\cdot 100+\left(n-1\right)\cdot 20\right]-70n=10n^2+20n$ (triệu đồng).\\ 
			Ta có $S\ge 2\,000$ $\Leftrightarrow 10n^2+20n-2\,000\ge 0\Leftrightarrow \hoac{&{n\ge  13{,}1} \\
			&{n\le-15{,}1}. 
			}$\\ 
			Do đó sau ít nhất sau $14$ năm thì sinh viên đó có thể mua  được chung cư $2$ tỉ đồng.
		\end{itemchoice}
	} 
\end{ex}

% nguồn: Ok_dungsai_lop11_Lan2
\begin{ex}%[1D2H2-7]%[ma:e8a2f555]
	Người ta trồng $3\,003$ cây theo dạng một hình tam giác như sau: hàng  thứ nhất trồng $1$ cây, hàng thứ hai trồng $2$ cây, hàng thứ ba trồng $3$ cây,\ldots, cứ tiếp tục trồng như thế cho đến khi hết số cây. Gọi số cây ở hàng thứ $n$ là $u_n$.
	\choiceTF 
		{\True Dãy số $u_n$ là $1$ cấp số cộng với số hạng đầu  $u_1=1$}
		{\True Số hạng thứ $3$ của  dãy này là $3$}
		{\True Tổng số  cây ở $10$ hàng đầu là $55$ cây}
		{Nếu  trồng theo quy luật trên ta sẽ trồng được $77$ hàng}
	\loigiai{\begin{itemchoice}
		\itemch {\bf Đúng}.\\ Vì $u_1=1$, $u_2=2$, $u_3=3$, \ldots. Dễ nhận thấy dãy số $(u_n)$ là cấp số cộng  có $u_1=1$, công sai $d=1$.
		\itemch {\bf Đúng}.\\ Ta có $u_3$ chính là số  số cây ở hàng thứ $3$. Vậy $u_3=3$.
		\itemch {\bf Đúng}.\\  Vì $S_{10}=\dfrac{\left(u_1+u_{10}\right)\cdot 10}{2}=\dfrac{(1+10)10}{2}=55$.
		\itemch  {\bf Sai}.\\ Vì $S=\dfrac{n\big[2u_1+(n-1)d\big]}{2}=3\,003$.\\
		Suy ra \begin{eqnarray*}
		&&\dfrac{n\big[2\cdot 1+(n-1)1\big]}{2}=3\,003\\ 
		 &\Leftrightarrow& n(n+1)=6\,006\\
		&\Leftrightarrow& n^2+n-6\,006=0\\ 
		&\Leftrightarrow&  \hoac{&n=77\\&n=-78.}
	\end{eqnarray*}
	Vì $n\in \mathbb{N^*}$ nên  $n=77$.\\ Vậy số hàng cây được trồng là $77$.
	\end{itemchoice}
} 
\end{ex}

% nguồn: Ok_dungsai_lop11_Lan2
\begin{ex}%[1D2H2-7]%[ma:3f0ff78a]
	Do nhu cầu đi lại của gia đình, anh Bình  quyết định thực hiện tích góp tiền để mua một chiếc ôtô HONDA CRV trị giá $1{,}259$ tỉ đồng.\\
	Đợt thứ nhất: anh Bình đã tích góp theo nguyên tắc tháng sau tích góp nhiều hơn tháng ngay trước đó số tiền là $2$ triệu đồng và cứ như thế đến tháng thứ $10$ anh phải góp  $21$ triệu đồng. Đến hết đợt thứ nhất anh Bình có tất cả $624$ triệu đồng.\\
	Đợt thứ hai kế tiếp: do muốn rút ngắn thời gian mua xe  thì số tiền còn lại anh tiếp tục tích góp với tháng đầu là $5$ triệu đồng và mỗi tháng tiếp theo số tiền gấp đôi tháng kề trước nó.
	 \choiceTF
		{\True Đợt thứ nhất anh Bình tích lũy tiền theo  dãy số với cấp số cộng có công sai là $d=2$ triệu và $u_1=3$ triệu}
		{Anh Bình tích lũy tiền hết  đợt thứ nhất trong $25$ tháng}
		{\True Đợt thứ hai anh Bình tích lũy  tiền theo dãy số với cấp số nhân có công bội là $q=2$ triệu và $u_1=5$ triệu}
		 {\True Để đủ tiền mua ôtô thì anh Bình thì anh Bình tích góp ít nhất $31$ tháng}
	 \loigiai{
		\begin{itemchoice}
			\itemch \textbf{Đúng.} \\ 
			Đợt thứ nhất anh Bình tích lũy theo cấp số cộng với công sai $d=2$ triệu.\\ 
			Theo đề bài ta có $21=u_1+9\cdot2\Rightarrow u_1=3$  triệu.
			\itemch \textbf{Sai.} \\ 
			Hết đợt thứ nhất anh Bình có  tất cả $624$ triệu đồng nên ta có\\
			$S_n=624\Rightarrow 624=\dfrac{n\left[2u_1+(n-1)\cdot3\right]}{2}  \Rightarrow n=24$ tháng.
			\itemch  \textbf{Đúng.} \\
			Theo đề  số tiền còn lại anh Bình tích góp theo cấp số nhân có công bội là $q=2$ triệu và $u_1=5$ triệu.
			\itemch  \textbf{Đúng.} \\
			Số tiền cần tích lũy ở đợt hai là $1259-624=635$ triệu  đồng.\\
			Từ đó ta có $635=5\cdot\dfrac{1-2^n}{1-2}$ $\Rightarrow n=7$  tháng.\\
			Tổng cộng hai đợt  cần có ít nhất $24+7=31$ tháng.
		\end{itemchoice}
	} 
\end{ex}

% nguồn: Ok_dungsai_lop11_Lan2
\begin{ex}%[1D2H2-7]%[ma:3fccef57]
	Trong  hội chợ tết, một công ty sữa muốn xếp $1\,000$ hộp sữa theo thứ tự từ trên xuống dưới như sau: Hàng thứ nhất có $1$ hộp sữa, hàng thứ hai có $3$ hộp sữa, hàng thứ ba có $5$ hộp sữa, $\dots$ cứ như thế, số lượng hộp sữa của hàng sau lớn hơn số lượng hộp sữa của hàng trước nó là $2$ hộp sữa (mô hình như hình dưới).
	\begin{center}
		\begin{tikzpicture}[line join=round,line cap=round,scale=0.2,transform shape]
				\definecolor{antiquebrass}{rgb}{0.8,0.58,0.46} 
				\definecolor{apricot}{rgb}{0.98,0.81,0.69} 
				\definecolor{babypink}{rgb}{0.96,0.76,0.76} 
				 \definecolor{bittersweet}{rgb}{1.0,0.44,0.37}
				 \definecolor{burntsienna}{rgb}{0.91,0.45,0.32}
				 \definecolor{citrine}{rgb}{0.89,0.82,0.04}
				 \definecolor{candypink}{rgb}{0.89,0.44,0.48}
				 \definecolor{inchworm}{rgb}{0.7,0.93,0.36}
				 \definecolor{pinksherbet}{rgb}{0.97,0.56,0.65}
				\definecolor{alizarin}{rgb}{0.82,0.1,0.26} 
				\definecolor{deepcarminepink}{rgb}{0.94,0.19,0.22} 
				 \definecolor{melon}{rgb}{0.99,0.74,0.71}
				 \definecolor{azure (colorwheel)}{rgb}{0.0,0.5,1.0}
				\definecolor{gainsboro}{rgb}{0.86,0.86,0.86} 
				\clip  (-2,0) rectangle (25,20);
				\tikzset{ong_tho/.pic={ 
				\def\G{ 
					(-1.02,-2.4)--(-1.06,-2.5) 
					 ..controls +(40:.02) and +(-140:.02)..(-1,-2.45)
					..controls +(-40:.2) and  +(-140:.25)..(-.52,-2.42)
					..controls  +(140:.2) and +(40:.2)..cycle
					(.05,-2.35) 
					 ..controls +(60:.15) and +(-140:.02)..(.5,-2.2)
					..controls +(-100:.25) and +(-40:.2)..cycle; 
				 }
				\draw  \G;
				\fill[citrine]  \G;
				 \def\Vg{
					 (-1.06,-2.5)
					 ..controls +(40:.02) and +(-140:.02)..(-1,-2.45)
					 ..controls +(-40:.2) and +(-140:.25)..(-.52,-2.42)--(-.43,-2.42)
					..controls +(-110:.2) and +(-100:.25)..cycle 
					(.5,-2.2) 
					..controls +(-100:.25) and +(-40:.2)..(.05,-2.35)--(-.05,-2.35) 
					..controls  +(-40:.3) and +(-120:.3)..(.6,-2.3)
					..controls +(130:.05)  and +(-30:.05)..cycle;
				 }
				 \fill[white] \Vg;
				\draw \Vg; 
				\def\A{ 
					 (-1.3,-.2)
					..controls +(-140:1)  and +(60:1.2)..(-2,-1.7)
					..controls  +(-70:.4) and +(-140:.35)..(-1.3,-1.85)
					--(-1.35,-2.45) 
					..controls +(25:.6) and +(170:.4)..(-.5,-2.35) 
					..controls  +(-10:.6) and +(170:.4)..(.45,-2.15)
					 ..controls +(-10:.2) and +(-120:.1)..(.8,-2)--(.75,-1.9)
					..controls +(35:.6)  and +(-25:.8)..(0,-.097)
					..controls +(150:.25) and +(35:.2)..cycle; 
				 }
				 \draw \A;
				\fill[citrine] \A; 
				\def\At{ 
					(-1.95,-1.7) 
					..controls  +(-70:.4) and +(-100:.6)..(-1.25,-1.45)--(-1.2,-1.45)
					..controls +(-95:.55) and  +(-70:.5)..(-2,-1.7)--cycle
					 (-1.35,-2.45)
					 ..controls +(25:.6) and +(170:.4)..(-.5,-2.35)
					..controls +(-10:.6)  and +(170:.4)..(.45,-2.15)
					..controls  +(-10:.2) and +(-120:.1)..(.8,-2)
					..controls +(-70:.4) and  +(20:.2)..(.3,-2.25)
					..controls +(-160:.2) and  +(-20:.3)..(-.5,-2.42)
					 ..controls +(160:.2) and +(30:.4)..(-1.3,-2.52)--cycle
					(.53,-.7)--(.55,-.8) 
					..controls +(-30:.2) and +(120:.4)..(.65,-1.5) 
					..controls +(-60:.2) and  +(-40:.4)..(.43,-1.6)
					..controls +(140:.2)  and +(-100:.1)..(.4,-1)
					..controls +(-160:.4) and +(150:.1)..(.55,-1.8) 
					..controls +(-30:.2)  and +(-65:.35)..(.7,-1.4)
					 ..controls +(125:.2) and +(-35:.35)..cycle;
				} 
				\fill[white]  \At;
				\draw \At; 
				 \def\T{
					 (-1.4,.5)
					..controls  +(110:.1) and +(140:.3)..(-1.5,.3)
					 ..controls +(-80:.1) and +(160:.1)..(-1.35,.1)--cycle
					(-.1,.8) 
					..controls +(-10:.2) and  +(70:.1)..(-.1,.47)--cycle;
				} 
				\fill[apricot]  \T;
				\draw  \T;
				\def\D{ 
					(-1.5,.6) 
					..controls +(103:1.5)  and +(83:1.1)..(-.1,.8)
					 ..controls +(-70:.85) and +(-83:1.05)..cycle;
				 }
				 \fill[apricot] \D;
				\draw  \D;
				 \def\dai{
					(-.6,-1.65)--(-.5,-2.4)--(-.25,-2.5)--(-.05,-2.3)--(-.2,-1.59) 
					..controls  +(-170:.2) and +(10:.2)..cycle;
				} 
				\fill[citrine] \dai; 
				\draw  \dai;
				 \draw (-1.2,-1.55)
					..controls +(-40:.4)  and +(-140:.3)..(.4,-1.35)
					 (-.9,-1.65)
					..controls  +(120:1) and +(70:.9)..(.1,-1.49);
				 \def\R{
					 (-1,.75)
					..controls +(170:.5) and +(110:.6)..(-1.3,-.15) 
					..controls +(-110:.1) and +(120:.1)..(-1.25,-.55) 
					..controls +(100:.1) and +(-100:.1)..(-1.22,-.2) 
					..controls +(-35:.2) and +(60:.4)..(-.6,-1.3) 
					..controls +(30:.35) and  +(-100:.4)..(-.2,-.05)
					..controls +(-50:.2) and  +(100:.2)..(0,-.45)
					..controls +(70:.2)  and +(-30:.1)..(-.05,0)
					..controls  +(150:.1) and +(-70:.2)..(-.1,.5)
					 ..controls +(100:.3) and +(30:.3)..(-.65,.8)--(-.63,.78)
					..controls +(30:.35) and +(40:.4)..(-.25,.25) 
					..controls +(130:.2)  and +(10:.1)..(-.65,.35)--(-.64,.33)
					..controls  +(10:.2) and +(120:.1)..(-.3,.2)
					 ..controls +(-140:.1) and +(20:.02)..(-.5,.1)
					..controls +(-160:.05)  and +(20:.05)..(-.75,.05)
					..controls +(-160:.05) and +(-20:.05)..(-1.05,.08) 
					..controls  +(75:.15) and +(160:.1)..(-.8,.3)--(-.82,.32)
					..controls +(160:.15)  and +(60:.1)..(-1.1,.1)
					..controls +(170:.4)  and +(170:.4)..(-1.05,.72)--cycle;
				 }
				\fill[white] \R; 
				 \draw \R;
				 \def\gay{
					(.35,-2.8) 
					..controls  +(-70:.15) and +(-92:1)..(.47,-2)
					 ..controls +(88:.15) and +(-100:.3)..(.4,.5)--(.42,.9)
					..controls +(100:.3) and  +(-100:.2)..(.55,1.3)
					..controls +(80:.7)  and +(110:.6)..(-.15,1.2)
					..controls  +(-40:.4) and +(85:.6)..(.2,.4)
					..controls  +(-80:.4) and +(95:.2)..cycle;
				} 
				\fill[antiquebrass] \gay; 
				\draw \gay; 
				 \def\dao{
					 (-1.9,-1.7)
					..controls +(140:1)  and +(-70:.4)..(-1.9,-.68)
					 ..controls +(-5:1) and +(10:1)..cycle;;
				} 
				\fill[pinksherbet]  \dao;
				\draw \dao; 
				\draw (-1.8,-1.7) 
					..controls +(160:.5) and  +(-70:.6)..(-1.9,-.7);
				\def\la{ 
					 (-1.7,-1.7)
					..controls +(110:.65)  and +(150:.5)..(-.8,-1.15)
					 ..controls +(-160:.5) and +(10:.8)..cycle
					(-1.8,-1.7) 
					..controls +(110:.65) and  +(50:.5)..(-2.55,-1.3)
					 ..controls +(-20:.4) and +(-175:.5)..cycle;
				} 
				 \fill[inchworm] \la;
				\draw \la; 
				 \def\tay{
					(.55,-.8) 
					..controls  +(110:.05) and +(-30:.05)..(.52,-.7)
					..controls  +(110:.05) and +(-30:.05)..(.5,-.56)
					..controls  +(110:.05) and +(-30:.05)..(.46,-.46)
					..controls +(150:.05)  and +(30:.05)..(.2,-.5)
					 ..controls +(-150:.05) and +(110:.05)..(.21,-.6)
					..controls +(-150:.02) and +(110:.05)..(.22,-.7) 
					..controls +(-150:.02)  and +(110:.05)..(.25,-.8)
					..controls +(-150:.02) and  +(110:.05)..(.28,-.9)
					..controls  +(20:.02) and +(110:.05)..(.45,-.95)
					..controls  +(-70:.05) and +(170:.3)..(.58,-1.24)
					..controls +(70:.2) and +(-40:.15)..cycle 
					 (-1.9,-1.76)
					 ..controls +(-40:.02) and +(-30:.02)..(-1.8,-1.79)
					..controls +(-40:.02)  and +(-30:.02)..(-1.7,-1.79)
					..controls +(-40:.02) and  +(-30:.02)..(-1.6,-1.8)
					 ..controls +(-5:.06) and +(-120:.4)..(-1.25,-1.55)
					..controls +(-110:.06) and  +(-170:.1)..(-1.5,-1.65)
					..controls +(95:.06) and +(10:.06)..(-1.6,-1.55) 
					..controls  +(170:.04) and +(10:.04)..(-1.8,-1.52)
					 ..controls +(170:.04) and +(10:.04)..(-1.9,-1.56)
					..controls +(-110:.04)  and +(160:.04)..cycle;
				 }
				 \fill[apricot] \tay;
				\draw  \tay;
				\draw  (-1.6,-1.75)
					 ..controls +(150:.04) and +(-140:.04)..(-1.6,-1.55)
					 (-1.7,-1.74)
					..controls +(150:.04)  and +(-140:.04)..(-1.7,-1.53)
					 (-1.8,-1.73)
					 ..controls +(150:.04) and +(-140:.04)..(-1.8,-1.52);
				\draw  (.52,-.69)
					..controls +(110:.05)  and +(30:.05)..(.22,-.7)
					(.5,-.56) 
					 ..controls +(110:.05) and +(30:.05)..(.21,-.6)
					(.55,-.8) 
					..controls +(110:.05) and +(30:.05)..(.25,-.8); 
				 \def\B{
					 (-.1,.52)--(-.05,.44)
					 ..controls +(-100:.15) and +(150:.2)..(.08,.1)
					..controls +(-88:1.05) and  +(-15:1.1)..(.43,.48)
					..controls +(150:.15) and  +(15:.2)..(.1,.55)--(0.02,.63)--cycle;
				 }
				\fill[burntsienna]  \B;
				\draw  \B;
				\draw  (.08,.1)
					 ..controls +(-15:.2) and +(-150:.1)..(.4,.15)
					(.43,.48) 
					..controls  +(-65:.15) and +(80:.15)..(.45,.2);
				\draw[fill=burntsienna,rotate=35] (.3,.5) ellipse (.08  and 0.04);
				 \draw (-.98,.6)
					..controls +(-170:.1) and  +(50:.12)..(-1.17,.52)
					(-.58,.65) 
					..controls +(40:.1) and  +(160:.12)..(-.35,.65);
				\draw (-1.3,.87) 
					 ..controls +(-30:.3) and +(-130:.3)..(-.4,1)
					 (-.93,.8)
					..controls +(-30:.07) and +(-130:.09)..(-.8,.81) 
					 ..controls +(50:.07) and +(-110:.12)..(-.65,.85);
				\draw (-1.3,-1.85) 
					..controls +(40:.2)  and +(-70:.12)..(-1.02,-1.3);
				 \fill[rotate=10,color=candypink] (-.67,.57) ellipse (.11 and 0.08);
				\fill[rotate=10,color=candypink!50]  (-.2,.57) ellipse (.12 and 0.05);
				 \fill[rotate=5,color=candypink!50] (-1.15,.5) ellipse (.11 and 0.04);
				 \draw[bittersweet] (-.92,.22)
					..controls +(-40:.1)  and +(-120:.2)..(-.5,.3)
					 ..controls +(-110:.2) and +(-60:.1)..cycle;
				\draw (.75,-1.9)..controls +(170:.2) and +(-60:.1)..(.5,-1.75); 
				 \draw (.46,1.15)
					..controls +(-170:.8) and +(110:.35)..(.4,1.4); 
					 }}
				 \tikzset{hop_sua/.pic={
				 \def\V{
					 (-2.15,1.9)
					..controls +(30:.35)  and +(150:.35)..(2.15,1.9)
					..controls  +(-30:.05) and +(30:.05)..(2.15,1.8)
					 ..controls +(150:.35) and +(30:.35)..(-2.15,1.8)
					..controls  +(150:.05) and +(-150:.05)..(-2.15,1.9)
					 (-2.15,-3.15)
					..controls +(-30:.5) and +(-150:.5)..(2.15,-3.15) 
					..controls +(-30:.05)  and +(30:.05)..(2.15,-3.25)
					..controls +(-150:.5) and  +(-30:.5)..(-2.15,-3.25)
					..controls +(150:.05)  and +(-150:.05)..cycle;
				 }
				 \fill[gainsboro] \V;
				\draw  \V;
				 \def\H{
					 (2.15,1.8)
					..controls +(150:.35) and +(30:.35)..(-2.15,1.8)--(-2.15,-3.15) 
					..controls +(-30:.5) and  +(-150:.5)..(2.15,-3.15)--cycle;
				 }
				 \fill[deepcarminepink!80] \H;
				\draw  \H;
				 \def\X{
					 (-1.8,1)
					 ..controls +(30:.8) and +(150:.8)..(1.8,1)
					..controls +(-30:.3) and +(30:.4)..(1.6,-.1) 
					..controls +(150:.35) and +(30:.35)..(-1.6,-.1) 
					..controls  +(150:.2) and +(-150:.2)..cycle;
				} 
				 \fill[melon] \X;
				\draw  \X;
				 \node[rotate=10,right,alizarin,scale=2.7,inner sep=0,align=left,
					font=\fontfamily{phv}\selectfont] at  (-1.7,.4)
					 {Ông};
				\node[rotate=-5,right,alizarin,scale=2.7,inner sep=0,align=left, 
					font=\fontfamily{phv}\selectfont] at  (0,.6)
					{Thọ}; 
				\node[right,white,scale=.8,inner sep=0,align=left, 
					font=\fontfamily{qag}\selectfont] at  (-1.2,-2.8)
					{Sữa  đặc có đường};
				 \node[,right,white,scale=.8,inner sep=0,align=left,
					font=\fontfamily{qag}\selectfont] at  (1.3,-2.25)
					 {$380$ g};
				 \node[right,azure (colorwheel),scale=.7,inner sep=0,align=left,
					font=\fontfamily{qag}\selectfont]  at (-.5,1.5)
					{VINAMILK}; 
					}} 
				\draw  (0,0)
					foreach \m in  {1,2,3,4,5}{
					(0+4.4*\m,5.4)  pic[scale=1]{hop_sua}
					 (.4+4.4*\m,4.4) pic[scale=.58]{ong_tho}
					 };
				\draw  (0,0)
					foreach \m in {2,3,4}{ 
					(0+4.4*\m,10.8)  pic[scale=1]{hop_sua}
					(.4+4.4*\m,9.8) pic[scale=.58]{ong_tho} 
					 };
				 \draw (0,0)
					 foreach \m in {3}{
					(0+4.4*\m,16.2) pic[scale=1]{hop_sua} 
					(.4+4.4*\m,15.2)  pic[scale=.58]{ong_tho}
					 };
				\node at (0,5)  [above] {\Large Hàng thứ ba};
				\node at  (0,10) [above] {\Large Hàng thứ hai};
				\node  at (0,16) [above] {\Large Hàng thứ nhất};
				\draw[->]  (-1.5,15.7)--(1.5,15.7);
				\draw[->] (-1.5,9.7)--(1.5,9.7); 
				\draw[->]  (-1.5,4.7)--(1.5,4.7);
			\end{tikzpicture}
	\end{center}
	\choiceTF 
		{\True Gọi $u_n$ là số hộp sữa ở hàng thứ $n$ thì $\left(u_n  \right)$ là một cấp số cộng có số hạng đầu $u_1=1$ và công sai $d=2$}
		{Số hộp sữa của hàng  thứ $10$ là $20$ hộp sữa}
		{\True Để xếp được $20$ hàng thì cần $400$ hộp  sữa}
		 {Hàng cuối cùng có $900$ hộp sữa}
	\loigiai{ 
		\begin{itemchoice}
			\itemch \textbf{Đúng.} \\ 
			 Gọi $u_n$ là số hộp sữa ở hàng thứ $n$ thì $\left(u_n \right)$ là một cấp số cộng có số hạng đầu $u_1=1$ và công sai $d=2$.
			\itemch  \textbf{Sai.} \\
			Số hộp sữa của hàng thứ $10$  là $u_{10}=u_1+9d=1+9\cdot2=19$ hộp sữa.
			\itemch  \textbf{Đúng.} \\
			Để xếp được $20$ hàng thì số  hộp sữa cần có là
			\[S_{20}=\dfrac{20}{2} \left(2u_1+19\cdot d\right)=\dfrac{20}{2} \left(2\cdot 1+19\cdot 2\right)=400\,\,\text{hộp  sữa.}\]
			\itemch  \textbf{Sai.} \\
			 Ta có $S_n=\dfrac{n}{2} \left[2\cdot1+\left(n-1\right)\cdot2\right]=n^2$.
			Giả sử với $1\,000$ hộp sữa ta chỉ xếp được nhiều  nhất $n$ hàng thì $n$ là số nguyên lớn nhất thoả mãn bất phương trình \[S_n \le 1\,000\Leftrightarrow n^2\le 1\,000\Leftrightarrow-10\sqrt{10} \le n\le 10\sqrt{10}.\] Suy ra $n=31$. \\
			Vậy, hàng  cuối cùng có $u_{31}=\dfrac{31}{2} \left(2u_1+30\cdot d\right)=\dfrac{31}{2} \left(2\cdot 1+30\cdot2\right)=961$ hộp sữa.
		\end{itemchoice}
	 }
\end{ex}

% nguồn: Ok_dungsai_lop11_Lan2
\begin{ex}%[1D2H2-7]%[ma:ef7515ea]
	Người ta trồng $3\,240$ cây theo một hình tam giác như sau: hàng thứ  nhất trồng $1$ cây, hàng thứ hai trồng $2$ cây, hàng thứ ba trồng $3$ cây, \ldots
	\choiceTF 
		{\True  Số cây mỗi hàng lập thành một cấp số cộng $(u_n)$ có số hạng đầu là $u_1=1$}
		{Số  cây mỗi hàng lập thành một cấp số cộng $(u_n)$ có công sai là $d=2$}
		{\True Có  tất cả $80$ hàng cây}
		{Hàng thứ $20$ trồng  được $40$ cây}
	\loigiai{ 
		 Số cây mỗi hàng (bắt đầu từ hàng thứ nhất) lập thành một cấp số cộng $(u_n)$ có $u_1=1$, $d=1$.\\
		Giả  sử có $n$ hàng cây thì $u_1+u_2+\cdots+u_n=3\,240=S_n$.\\
		 Ta có
		\begin{eqnarray*}
			&  & S_n=3\,240\\
			&\Leftrightarrow &  nu_1+\dfrac{n(n-1)}{2}d=3\,240\\
			 &\Leftrightarrow & n^2+n-6\,480=0\\
			 &\Leftrightarrow & n=80.
		\end{eqnarray*}
		Số cây hàng thứ $20$  trồng được là $u_{20}=u_1+19d=20$.\\
		\begin{itemchoice}
			 \itemch {\bf Đúng}.\\
			\itemch  {\bf Sai}.\\
			\itemch  {\bf Đúng}.\\
			 \itemch {\bf Sai}.\\
		\end{itemchoice}
	 }
\end{ex}

% nguồn: Ok_dungsai_lop11_Lan1
\begin{ex}%[1D2V2-7]%[ma:b409b325]
 Hãy cho  biết tính đúng sai của mỗi mệnh đề dưới đây?
 \choiceTF 
  {Giá của một chiếc xe ô tô lúc mới mua là $680$ triệu đồng. Cứ sau mỗi năm sử dụng, giá của chiếc xe ô tô giảm $55$ triệu đồng. Giá còn lại của chiếc xe sau  $5$ năm sử dụng là $450$ triệu đồng}
  {Mai làm việc ở công ty X với mức lương năm đầu  tiên là $60$ triệu đồng. Sau mỗi năm lương của Mai được tăng thêm $5$ triệu đồng. Tổng lương của Mai sau $20$ năm làm việc ở công ty X là $2200$ triệu đồng }
  {\True Hùng đang tiết kiệm để mua một cây guitar. Trong tuần đầu tiên, anh ta để dành $42$ đô la, và trong mỗi tuần tiếp  theo, anh ta đã thêm $8$ đô la vào tài khoản tiết kiệm của mình. Cây guitar Hùng cần mua có giá $400$ đô la. Tuần thứ $46$ thì anh ấy có đủ tiền để mua guitar đó}
  {\True Người ta trồng $3003$ cây theo hình một tam giác như sau: hàng thứ nhất có $1$ cây, hàng thứ hai có $2$ cây, hàng thứ ba có $3$ cây,$\ldots$ Có tất cả $77$  hàng cây được trồng}
  \loigiai{
  \begin{itemchoice}
    \itemch Giá của chiếc xe sau $n$ năm là $u_n=680-55(n-1)$.\\
   Vậy sau 5 năm sử dụng giá của  chiếc xe là
   $u_5=680-55(5-1)=460$ (triệu  đồng).
   \itemch Vì mỗi năm lương của Mai được tăng thêm $5$ triệu đồng nên lương các năm của Mai là  cấp số cộng với $ u_1=60$ và $d=5$.\\
   Vậy $S_{20}=20u_1+5\times  \dfrac{20\cdot 19}{2} = 2150 $ triệu đồng.
   \itemch Gọi $u_1=8$ là số tiền mà Hùng để dành trong  tuần thứ $2$.\\
   $u_2=8$ là số tiền mà  Hùng để dành trong tuần thứ $3$.\\
   $\cdots$  \\
   $u_n$ là  là số tiền mà Hùng để dành trong tuần thứ $n+1$.\\
   Ta có $(u_n)$ là cấp số cộng có số hạng $u_1=8$ và công sai $d=0$.\\ 
   Ta  có $S_n=42+\dfrac{n\left(2u_1+(n-1)d\right)}{2}=400 \Leftrightarrow n=44{,}75$.\\
   Suy ra vào tuần thứ $46$ thì anh ấy có đủ  tiền để mua guitar đó.
   \itemch Theo đề ta  có\\
   Hàng $1$ có $1$ cây.\\ 
    Hàng $2$ có $2$ cây.\\
   $\ldots$\\ 
    Hàng $n$ có $n$ cây.\\
   Với cách trồng theo hình thức này thì số cây của các hàng tuân  theo cấp số cộng có $u_1=1$ và công sai $d=1$.\\
   Giả sử $3003$ cây trồng được $n$ hàng, ta  có
    $S_n=nu_1+\dfrac{n(n-1)d}{2}
    \Leftrightarrow 3003=n+\dfrac{n(n-1)}{2}
   \Leftrightarrow  n^2+n-6006=0
   \Leftrightarrow  \hoac{&n=77\\&n=-78 \, \text{(loại)}.}$\\
   Vậy $n=77$. 
  \end{itemchoice}
 } 
\end{ex}

% nguồn: Ok_dungsai_lop11_Lan2
\begin{ex}%[1D2V2-7]%[ma:6f3f6747]
	Khi kí kết  hợp đồng lao động với người lao động, một doanh nghiệp đề xuất hai phương án trả lương như sau\\
	\textbf{Phương án 1}: Năm thứ nhất, tiền lương  là $120$ triệu. Kể từ năm thứ hai trở đi, mỗi năm tiền lương được tăng $18$ triệu.\\
	\textbf{Phương án 2}: Quý thứ nhất, tiền lương là $24$ triệu.  Kể từ quý thứ hai trở đi, mỗi quý tiền lương được tăng $1{,}8$ triệu. Xét tính đúng sai của các khẳng định sau
	 \choiceTF
		{\True Trong phương án 1: dãy số tiền lương là cấp số cộng có số hạng đầu tiên là $u_1=120$, công  sai $d_1=18$}
		{Trong phương án 1: tiền lương người lao động nhận được trong năm thứ ba là $174$ triệu} 
		{\True Trong phương án  1: tổng tiền lương người lao động nhận được trong ba năm là $414$ triệu}
		{\True Nếu kí hợp đồng lao động trong  ba năm, với mong muốn nhận được tổng số tiền lương cao nhất thì người lao động nên chọn phương án 1}
	\loigiai{ 
		\textbf{Phương  án 1:}\\
		Dãy số tiền lương theo năm là một cấp số cộng có  số hạng đầu $u_1=120$ và công sai $d_1=18$.\\
		Tiền lương năm thứ ba là $u_3 = 120 + 2\cdot 18 = 156$ triệu.\\ 
		Tổng tiền lương nhận được trong ba năm là $S_3  = \dfrac{(u_1+u_3).3}{2} = \dfrac{(120+156)\cdot 3}{2} = 414$ triệu.\\
		 \textbf{Phương án 2:}\\
		Dãy số tiền lương theo quý là  cấp số cộng có số hạng đầu tiên là $v_1 = 24$, công sai $d_2=1{,}8$. Ba năm tương ứng với $12$ quý.\\
		Tiền lương quý thứ $12$ là $v_{12} = 24 + 11\cdot 1{,}8  = 43{,}8$ triệu.\\
		Tổng số tiền lương nhận được  sau $3$ năm tương ứng với $12$ quý là\\
		$S_{12} =  \dfrac{(v_1+v_{12})\cdot 12}{2} = (24+43{,}8)\cdot 6 = 406{,}8$ triệu.\\
		\textbf{So  sánh:}\\
		Ta nhận thấy  $414 > 406{,}8$.
		\begin{itemchoice}
			\itemch Trong phương án 1, dãy số tiền lương là cấp số cộng có  số hạng đầu tiên là $u_1=120$, công sai $d_1=18$.
			\itemch  Trong phương án 1, tiền lương năm thứ ba là $u_3 = 156$ triệu.
			\itemch Trong phương án 1, tổng tiền lương trong  ba năm là $414$ triệu.
			\itemch Vì tổng lương phương án 1 ($414$ triệu) lớn hơn  tổng lương phương án 2 ($406,8$ triệu), người lao động nên chọn phương án 1.
		\end{itemchoice}
	 }
\end{ex}

% nguồn: Ok_dungsai_lop11_Lan2
\begin{ex}%[1D2V2-7]%[ma:f68da018]
	Một nhà hát có $25$ hàng ghế với $16$ ghế ở hàng thứ nhất, $18$ ghế ở hàng thứ hai, $20$ ghế ở hàng thứ ba và cứ tiếp tục theo quy luật đó, tức là hàng sau nhiều hơn hàng liền trước nó $2$ ghế. Gọi $u_n$ (ghế) là tổng số ghế ở hàng thứ $n$. Xét  tính đúng sai của các khẳng định sau
	 \choiceTF
		{\True $u_2=18$} 
		{\True Dãy số $(u_n)$ là cấp số cộng có công sai $d=2$} 
		{Số ghế ở hàng thứ $20$ nhỏ hơn  $54$}
		{Tổng số ghế trong nhà hát  hơn $1\,000$}
	 \loigiai{
		Số ghế ở mỗi hàng của nhà hát lập thành một cấp số cộng $(u_n)$ gồm $25$ số hạng, với số hạng đầu $u_1=16$ và công sai $d=2$.\\ 
		Ta có $u_2 = u_1+d =  16+2=18$.\\
		Số ghế ở hàng thứ $20$ là $u_{20} = u_1+19d = 16+19\cdot  2 = 54$.\\
		 Tổng số ghế trong nhà hát là
		\begin{eqnarray*}
			S_{25} &=&  u_1+u_2+\ldots +u_{25}\\
			& =&  \dfrac{25}{2}[2u_1+(25-1)d]\\
			& = &\dfrac{25}{2}(2\cdot 16+24\cdot 2)\\ 
			 &=&1\,000.
		\end{eqnarray*}
		\begin{itemchoice}
			\itemch $u_2=18$. 
			\itemch Dãy số  $(u_n)$ là cấp số cộng có công sai $d=2$.
			\itemch Số ghế ở hàng  thứ $20$ bằng $54$.
			\itemch  Tổng số ghế trong nhà hát bằng $1\,000$.
		\end{itemchoice}
	 }
\end{ex}

% nguồn: Ok_dungsai_lop11_Lan2
\begin{ex}%[1D2V2-7]%[ma:d0316f1e]
	\immini[thm] 
	{Do nhu cầu đi lại của gia đình,  anh Bình quyết định thực hiện tích góp tiền để mua một chiếc ô tô HONDA CRV trị giá $1{,}259$ tỉ đồng.\\
	 \textbf{Đợt thứ nhất}: anh Bình đã tích góp theo nguyên tắc tháng sau tích góp nhiều hơn tháng ngay trước đó số tiền là $2$ triệu đồng và cứ như thế đến tháng thứ $10$ anh phải góp $21$ triệu}
	{\begin{tikzpicture}[line cap=round,line join=round,scale=0.2]
		\path[fill=black] (12.03,3.71)..  controls +(179:2.5) and +(-12:0.9)..
			(7.15,4.02).. controls +(168:0.7)  and +(0:0.1)..
			(6.17,4.2).. controls +(180:0.1)  and +(-90:0)..
			(6.03,4.23).. controls +(90:0)  and +(-23:0.1)..
			(5.9,4.33).. controls +(153:0.4) and +(-68:0.6).. 
			(4.88,5.59).. controls +(112:0.3)  and +(-10:0.1)..
			(4.64,5.89).. controls  +(169:0.2) and +(-75:0.2)..
			(4.44,6.1) --  (4.39,6.26) -- (4.16,6.29).. controls +(173:0.6) and +(-48:0.3)..
			(3.36,6.58).. controls +(134:0.3) and +(-90:0.1).. 
			 (3,7.21).. controls +(90:0.1) and +(166:0.1)..
			(3.12,7.3).. controls  +(0:0) and +(-27:0)..
			(3.13,7.32).. controls +(117:0)  and +(-135:0.1)..
			(3.21,7.46).. controls  +(45:0.1) and +(-78:0.1)..
			(3.28,7.72).. controls +(100:0.3) and +(148:0.2).. 
			(3.36,7.93)..  controls +(-30:0.2) and +(-104:0.5)..
			(3.69,8.26).. controls +(72:0) and  +(-135:0.1)..
			(3.78,8.38)..  controls +(41:0.1) and +(179:3.9)..
			(7.88,8.39).. controls +(-1:3.5) and +(159:0.1).. 
			 (12.05,8.27).. controls +(-18:0.1) and +(133:0.2)..
			(12.44,7.95) -- (12.69,7.68)  -- (13.37,7.72).. controls +(3:0.4) and +(-178:0.3)..
			 (14.62,7.79).. controls +(4:0.3) and +(180:0.1)..
			(15.35,7.83).. controls +(0:0.3)  and +(-71:0.5)..
			(15.48,8.28).. controls +(108:0.2)  and +(-68:0.2)..
			(15.21,9.03).. controls +(112:0.2) and +(-90:0).. 
			(15.06,9.42).. controls +(90:0.1)  and +(145:0.1)..
			(15.46,9.37).. controls +(-30:0.1)  and +(111:0.7)..
			(16.21,7.9).. controls +(-69:0.3)  and +(124:0)..
			(16.44,7.33).. controls +(-63:0)  and +(-178:0.3)..
			 (16.78,7.29).. controls +(3:0.2) and +(-135:0)..
			(17.13,7.34).. controls +(56:0) and  +(-98:0.3)..
			 (17.25,7.88).. controls +(81:2.3) and +(-164:1.2)..
			(19.49,11.8).. controls +(16:0.1) and  +(180:0.2)..
			(20,11.85).. controls +(-4:0.3)  and +(153:0.3)..
			(20.61,11.68).. controls  +(-27:0.5) and +(124:0.4)..
			(21.5,10.93)..  controls +(-57:0.8) and +(92:1.2)..
			(22.22,8.22) -- (22.24,7.64) --  (22.42,7.67).. controls +(6:0.1) and +(180:0)..
			(22.65,7.69).. controls +(0:0) and  +(-141:0.1)..
			(22.9,7.86)..  controls +(45:0.1) and +(-166:0)..
			(23.16,8.06).. controls +(14:0)  and +(-177:0.7)..
			(24.52,8.15).. controls +(3:1.4) and  +(-176:1.4)..
			(27.81,8.34) -- (28.71,8.4) -- (28.94,8.28).. controls +(-31:0.2) and  +(180:0.1)..
			(29.35,8.15).. controls +(0:0.1) and  +(-135:0)..
			(29.55,8.18).. controls +(63:0) and  +(-98:0.3)..
			(29.65,8.76).. controls +(81:1.7)  and +(-155:0.8)..
			(31.08,11.74).. controls +(25:0.3) and  +(152:0.3)..
			(31.96,11.75)..  controls +(-28:0.7) and +(94:1.2)..
			(32.99,9.47).. controls  +(-86:0.3) and +(0:0)..
			(33.04,8.93).. controls +(-45:0.1) and +(-112:0.1).. 
			(33.4,9.24)  -- (33.44,9.36) -- (33.49,9.16).. controls +(-74:0.4) and +(71:0.2)..
			(33.53,8.56).. controls  +(-108:0.4) and +(41:0.3)..
			(33.04,7.77).. controls +(-140:0.2) and +(72:0).. 
			(32.79,7.51).. controls +(-108:0)  and +(37:0)..
			(32.71,7.39)..  controls +(-129:0.1) and +(83:0.3)..
			(32.6,6.96).. controls +(-97:0.8) and  +(48:0.6)..
			(31.85,5.06) --  (31.62,4.82) -- (31.15,4.79).. controls +(-178:0.3) and +(2:0.7)..
			(29.33,4.72).. controls +(-178:0.7) and  +(6:0.1)..
			(27.82,4.67).. controls  +(-174:0.1) and +(3:0.2)..
			(27.31,4.64).. controls +(-176:0.3)  and +(5:0.8)..
			 (24.5,4.41).. controls +(-175:0.4) and +(5:1)..
			(22.65,4.25)..  controls +(-175:0.3) and +(8:0.2)..
			(21.64,4.14).. controls +(-175:0.2) and +(6:0.1).. 
			(21.08,4.07).. controls +(-174:0.1) and +(23:0.1).. 
			(20.69,3.95).. controls  +(-152:0.1) and +(7:0.2)..
			(20.2,3.81).. controls  +(-172:0.7) and +(-1:3.2)..
			(12.03,3.71)  -- cycle
			(9.65,4.41).. controls +(2:0.5)  and +(-176:0.3)..
			 (11.15,4.48).. controls +(4:0.3) and +(-176:0.5)..
			(12.63,4.57).. controls +(3:1) and +(-176:0.7).. 
			(15.96,4.79).. controls +(4:0.3) and  +(-174:0.1)..
			(16.61,4.85)..  controls +(0:0.1) and +(-135:0)..
			(16.82,4.88).. controls +(45:0) and  +(-69:0.2)..
			(16.71,5.23) -- (16.6,5.57) -- (16.16,5.59).. controls +(178:1) and +(-79:0.1).. 
			(14.88,5.7).. controls  +(104:0) and +(-1:1)..
			 (13.72,5.76).. controls +(179:0.6) and +(-2:0.6)..
			(11.47,5.81).. controls +(178:0.6) and +(-2:0.2).. 
			(9.92,5.86)  -- (9.48,5.86) -- (9.36,5.58).. controls +(-114:0.4) and +(45:0.4)..
			 (8.7,4.6).. controls +(-135:0.2) and +(90:0)..
			(8.47,4.36).. controls +(0:0) and  +(180:0.1)..
			 (8.6,4.36).. controls +(8:0.1) and +(-178:0.5)..
			(9.65,4.41)  -- cycle
			(29.23,5.24).. controls +(124:0.3) and  +(-73:0.3)..
			(28.8,6.14) -- (28.73,6.4) -- (28.35,6.41) -- (27.98,6.42) -- (27.97,6.52).. controls  +(95:0.1) and +(-4:0.6)..
			(27.25,6.7)..  controls +(175:0.2) and +(-12:0.1)..
			(26.67,6.77)..  controls +(168:0.2) and +(54:0.1)..
			(26.18,6.75).. controls  +(-135:0) and +(0:0)..
			(26.07,6.7).. controls  +(180:0) and +(80:0.1)..
			(25.99,6.6).. controls +(-102:0.2)  and +(2:2.3)..
			(23.51,6.28).. controls  +(-178:0.7) and +(10:0.1)..
			(21.96,6.21)..  controls +(-170:0.1) and +(7:0.1)..
			(21.6,6.15).. controls +(-169:0.2) and  +(75:0.5)..
			(21.28,5.61).. controls  +(-105:0.2) and +(56:0)..
			 (21.18,5.28).. controls +(-108:0.1) and +(180:1.1)..
			(22.32,5.22).. controls  +(0:1.7) and +(177:0.9)..
			(28.05,5.11).. controls +(-4:0.4) and  +(180:0.1)..
			(28.97,5.06).. controls  +(0:0.1) and +(90:0)..
			(29.14,5.04).. controls +(-90:0) and +(180:0.1).. 
			(29.26,5.02)  -- (29.38,5.02) -- cycle
			(3.85,6.56).. controls +(173:0.2) and  +(-48:0.1)..
			(3.49,6.7).. controls +(130:0.1) and +(-76:0).. 
			(3.37,6.88)..  controls +(104:0) and +(0:0)..
			 (3.27,6.95) -- (3.19,6.95) -- (3.27,6.88).. controls +(-37:0.1) and +(126:0.1)..
			(3.43,6.69).. controls  +(-52:0.2) and +(180:0.3)..
			(3.92,6.53).. controls +(0:0.2)  and +(-7:0.2)..
			(3.85,6.56)  -- cycle
			(8.79,6.55).. controls +(173:0.1) and  +(6:0.1)..
			(8.49,6.55).. controls +(-173:0.1)  and +(180:0.2)..
			 (8.64,6.54).. controls +(0:0.2) and +(-6:0.1)..
			 (8.79,6.55) -- cycle
			(19.44,6.62).. controls  +(90:0) and +(27:0)..
			 (19.4,6.65).. controls +(-153:0) and +(90:0)..
			(19.37,6.57).. controls +(-90:0)  and +(-153:0)..
			(19.4,6.54)..  controls +(27:0) and +(-90:0)..
			(19.44,6.62) --  cycle
			(7.68,6.59).. controls +(180:0.2) and  +(0:0.2)..
			(7.07,6.59).. controls +(-177:0.2)  and +(180:0.3)..
			 (7.37,6.57).. controls +(0:0.3) and +(-3:0.2)..
			(7.68,6.59)  -- cycle
			(6.59,6.62)..  controls +(169:0.1) and +(10:0.1)..
			(6.38,6.62).. controls +(-170:0.1) and  +(180:0.1)..
			(6.47,6.61).. controls +(-5:0.1) and  +(-11:0.1)..
			(6.59,6.62) -- cycle 
			(19.79,6.65).. controls +(90:0) and +(0:0).. 
			(19.76,6.7).. controls +(180:0) and  +(108:0)..
			(19.74,6.65).. controls +(-56:0) and  +(180:0)..
			(19.78,6.59).. controls  +(0:0) and +(-90:0)..
			(19.79,6.65) --  cycle
			(19.85,6.91).. controls  +(63:0) and +(22:0.1)..
			(19.74,6.91).. controls +(-157:0.1)  and +(180:0.1)..
			(19.75,6.88).. controls +(0:0) and  +(-135:0)..
			(19.85,6.91) -- cycle 
			(19.83,7.14).. controls  +(82:0.1) and +(0:0)..
			(19.82,7.26).. controls  +(180:0) and +(90:0.1)..
			(19.79,7.16).. controls +(-90:0.1) and  +(27:0)..
			(19.75,7.04)..  controls +(-162:0) and +(180:0)..
			(19.76,7.02).. controls +(0:0) and  +(-96:0.1)..
			 (19.83,7.14) -- cycle
			(16.69,7.1).. controls +(0:0.1) and  +(-45:0)..
			(16.84,7.13).. controls +(117:0)  and +(180:0.1)..
			(16.91,7.16).. controls  +(10:0.1) and +(-12:0.1)..
			(16.88,7.19).. controls +(165:0.2) and +(53:0).. 
			(16.19,7.15).. controls +(-122:0.1) and +(-173:0.2).. 
			(16.37,7.07).. controls +(6:0.1) and  +(-174:0.1)..
			(16.69,7.1) --  cycle
			(19.72,7.18).. controls +(90:0.1) and  +(80:0.1)..
			(19.62,7.17).. controls  +(-101:0.1) and +(-121:0.1)..
			(19.65,7.16)..  controls +(53:0.1) and +(90:0)..
			 (19.69,7.15).. controls +(-90:0) and +(180:0)..
			(19.7,7.09).. controls +(0:0)  and +(-90:0)..
			(19.72,7.18)  -- cycle
			(19.32,7.25).. controls  +(90:0) and +(-18:0)..
			(19.27,7.33).. controls +(162:0) and  +(90:0.1)..
			(19.23,7.25)..  controls +(-90:0.1) and +(-162:0)..
			(19.27,7.17).. controls  +(18:0) and +(-90:0)..
			 (19.32,7.25) -- cycle
			(19.6,7.7).. controls +(117:0) and +(90:0).. 
			(19.59,7.65)..  controls +(-90:0.1) and +(-108:0)..
			(19.6,7.61).. controls +(63:0) and  +(-72:0)..
			(19.6,7.7)  -- cycle
			(26.07,7.62).. controls +(3:0.8)  and +(-90:0)..
			(27.02,7.73)..  controls +(90:0) and +(6:0.1)..
			(26.86,7.74).. controls  +(-174:0.1) and +(4:0.3)..
			(26.18,7.68)..  controls +(-176:0.7) and +(45:0)..
			 (25.3,7.6).. controls +(-124:0) and +(-176:0.7)..
			(26.07,7.62)  -- cycle
			 (31.43,7.77).. controls +(90:0) and +(0:0)..
			 (31.4,7.83).. controls +(180:0) and +(135:0)..
			(31.39,7.78).. controls +(-72:0) and +(45:0).. 
			(31.37,7.68).. controls +(-124:0)  and +(117:0)..
			(31.35,7.6).. controls +(-34:0)  and +(-90:0.1)..
			 (31.43,7.77) -- cycle
			(31.91,7.74)..  controls +(27:0) and +(7:0.1)..
			(31.79,7.74).. controls +(-174:0.1) and +(0:0).. 
			(31.6,7.73).. controls  +(-135:0) and +(90:0)..
			(31.57,7.69).. controls +(-90:0) and +(-158:0.1).. 
			(31.91,7.74) --  cycle
			 (28.56,7.8).. controls +(5:0.5) and +(180:0)..
			(29.5,7.87).. controls +(14:0) and +(-45:0).. 
			(29.54,7.9).. controls +(135:0) and  +(4:0.6)..
			(28.36,7.84).. controls +(-176:0.6)  and +(0:0)..
			(27.15,7.76).. controls +(180:0) and  +(117:0)..
			(27.11,7.72).. controls +(-56:0) and  +(-177:0.2)..
			(27.41,7.7)..  controls +(4:0.2) and +(-175:0.5)..
			(28.56,7.8) -- cycle 
			(31.31,7.82).. controls +(135:0.1) and +(113:0.1).. 
			(31.28,7.73)..  controls +(-72:0.1) and +(-108:0.1)..
			(31.32,7.73)..  controls +(76:0) and +(-45:0)..
			 (31.31,7.82) -- cycle
			(30.98,7.8).. controls +(90:0.1) and  +(56:0)..
			(30.94,7.84).. controls +(-146:0) and  +(117:0)..
			(30.92,7.76).. controls +(-56:0.1)  and +(-90:0.1)..
			(30.98,7.8)  -- cycle
			(31.26,7.9).. controls +(90:0) and +(0:0).. 
			 (31.22,7.93).. controls +(180:0) and +(117:0)..
			(31.2,7.9).. controls +(-45:0)  and +(180:0)..
			(31.24,7.86).. controls  +(0:0) and +(-90:0)..
			(31.26,7.9)  -- cycle
			(4.55,8.15).. controls +(150:0.1)  and +(90:0)..
			(4.45,8.15).. controls +(-90:0) and  +(180:0)..
			(4.53,8.11).. controls +(0:0.1)  and +(-31:0.1)..
			(4.55,8.15)  -- cycle
			 (11.81,8.14).. controls +(180:0) and +(0:0)..
			(11.67,8.14).. controls  +(-166:0) and +(180:0.1)..
			(11.73,8.12).. controls  +(0:0.1) and +(-18:0)..
			 (11.81,8.14) -- cycle
			(19.3,10.17).. controls +(90:0) and +(-27:0).. 
			(19.26,10.21)..  controls +(153:0) and +(90:0)..
			(19.23,10.19)..  controls +(-90:0) and +(180:0)..
			(19.26,10.16)..  controls +(0:0) and +(-90:0)..
			(19.3,10.17) -- cycle 
			(3.13,8.83).. controls +(97:0.5) and +(-121:0.1).. 
			(3.1,9.52).. controls  +(55:0.1) and +(89:0.5)..
			(3.17,9.06)..  controls +(-91:0.4) and +(-83:0.2)..
			(3.13,8.83) --  cycle
			 (6.25,9.21).. controls +(179:1.3) and +(-48:0.2)..
			(4.61,9.37).. controls +(130:0.2) and +(-113:0.1).. 
			(3.97,10.99).. controls  +(72:0) and +(-156:0.1)..
			(4.16,11.12).. controls +(22:0.1) and +(-108:0).. 
			 (4.36,11.24).. controls +(63:0) and +(-48:0.1)..
			(4.23,11.45).. controls +(131:0.3) and +(-92:0.2).. 
			(3.98,12.04)..  controls +(82:0.1) and +(128:0.1)..
			(4.07,12.01) -- (4.15,11.91)  -- (4.26,11.98).. controls +(40:0.1) and +(-124:0)..
			(4.41,12.12).. controls +(45:0.1)  and +(12:0.1)..
			 (4.37,12.16).. controls +(-165:0.1) and +(-70:0.2)..
			(4.06,12.45)..  controls +(108:0.1) and +(139:0.1)..
			(4.1,12.5).. controls +(-45:0.1)  and +(76:0)..
			 (4.15,12.37).. controls +(-104:0.1) and +(-135:0.1)..
			(4.37,12.21).. controls +(45:0.1) and  +(-38:0.1)..
			(4.33,12.38).. controls +(139:0.1)  and +(157:0.2)..
			(4.39,12.39)..  controls +(-23:0.2) and +(-180:3.3)..
			(7.99,12.35).. controls +(0:3.9)  and +(-173:2.8)..
			(14.08,12.68) -- (15.37,12.83) -- (15.61,12.98).. controls +(32:0.3) and  +(135:0)..
			(16.25,13.29).. controls +(-45:0)  and +(27:0.1)..
			 (16.08,13.19).. controls +(-157:0.2) and +(135:0)..
			(15.49,12.86).. controls +(-45:0)  and +(-171:1)..
			(16.54,13).. controls +(8:0.4)  and +(135:0)..
			 (17.2,13.08).. controls +(-45:0) and +(90:0.1)..
			(17.22,12.95).. controls +(-87:0.2)  and +(23:0.3)..
			(16.55,12.08).. controls +(-156:0.3) and +(-6:0.2).. 
			 (15.23,11.83).. controls +(173:0.2) and +(12:1.6)..
			(12.29,11.28)  -- (11.93,11.21) -- (11.79,10.88).. controls +(-115:0.2) and +(68:0.2)..
			(11.48,10.16).. controls +(-112:0.8) and  +(10:0.3)..
			(10.77,9.2).. controls +(-172:0.2)  and +(-1:1.8)..
			(6.25,9.21) -- cycle 
			(11.06,9.38).. controls +(45:0.1) and  +(0:0)..
			(11.11,9.45)..  controls +(180:0) and +(53:0.1)..
			(11.03,9.38)..  controls +(-135:0.1) and +(180:0)..
			 (10.98,9.31).. controls +(0:0) and +(-127:0.1)..
			 (11.06,9.38) -- cycle
			(4.76,9.36).. controls +(90:0) and +(56:0).. 
			(4.68,9.38).. controls +(-135:0)  and +(180:0)..
			(4.71,9.35).. controls +(0:0) and  +(-90:0)..
			 (4.76,9.36) -- cycle
			(9.63,9.4)..  controls +(31:0.1) and +(0:0)..
			(9.69,9.45).. controls +(180:0) and +(31:0.1).. 
			 (9.56,9.4).. controls +(-149:0.1) and +(180:0)..
			(9.51,9.35).. controls +(0:0)  and +(-149:0.1)..
			(9.63,9.4)  -- cycle
			(7.08,11.38)..  controls +(79:0.2) and +(98:0.2)..
			(7.18,11.25).. controls  +(-84:0.4) and +(-176:0.3)..
			(7.52,10.89)..  controls +(0:0.1) and +(-135:0)..
			(7.81,10.93).. controls +(63:0) and +(-98:0.1).. 
			(7.87,11.22)..  controls +(82:0.1) and +(-135:0)..
			 (7.95,11.52).. controls +(45:0) and +(56:0)..
			(7.96,11.51).. controls +(-117:0) and  +(153:0)..
			(7.98,11.45).. controls +(0:0) and  +(94:0.2)..
			(8.04,11.18) -- (8.07,10.92) -- (8.11,11.12).. controls +(77:0.5) and  +(0:0.5)..
			 (7.55,11.66).. controls +(180:0.5) and +(90:0.3)..
			(6.99,11.29)..  controls +(-90:0.1) and +(98:0.1)..
			(7.01,10.92).. controls +(-77:0.1) and +(-92:0.2).. 
			(7.05,11.04).. controls +(90:0.1)  and +(-104:0)..
			(7.08,11.38)  -- cycle
			(7.42,11.03).. controls  +(119:0.1) and +(-129:0.1)..
			(7.37,11.45).. controls  +(63:0) and +(180:0)..
			(7.46,11.48).. controls +(0:0)  and +(-146:0)..
			(7.58,11.51)..  controls +(45:0) and +(124:0)..
			(7.66,11.5).. controls +(-63:0)  and +(-101:0.1)..
			(7.7,11.52).. controls +(79:0.1) and +(90:0.1).. 
			(7.72,11.49).. controls +(-86:0.2) and +(31:0.1).. 
			 (7.61,11).. controls +(-144:0.1) and +(-63:0.1)..
			 (7.42,11.03) -- cycle
			(12.79,11.46).. controls +(15:0.2) and  +(180:0.1)..
			(13.3,11.57)..  controls +(0:0) and +(-112:0.1)..
			 (13.4,11.64).. controls +(76:0) and +(180:0)..
			(13.45,11.71).. controls +(0:0) and  +(-90:0)..
			(13.48,11.75).. controls +(90:0)  and +(45:0)..
			 (13.44,11.76).. controls +(-135:0) and +(-72:0)..
			(13.39,11.81).. controls +(108:0.1) and  +(90:0.1)..
			(13.27,11.83).. controls  +(-90:0) and +(-162:0)..
			 (13.31,11.8).. controls +(45:0) and +(45:0)..
			(13.32,11.79).. controls +(-108:0) and  +(0:0)..
			(13.24,11.74).. controls  +(180:0) and +(117:0)..
			(13.21,11.71).. controls +(-63:0)  and +(0:0)..
			 (13.19,11.67).. controls +(180:0) and +(45:0.1)..
			 (13.06,11.6).. controls +(-138:0.1) and +(-130:0.2)..
			(13.07,11.68).. controls  +(50:0.1) and +(-90:0)..
			(13.16,11.85).. controls +(90:0.1) and  +(0:0.1)..
			(13.07,11.92).. controls +(180:0.1)  and +(144:0.1)..
			 (13.03,11.81).. controls +(-45:0) and +(56:0)..
			(13.03,11.73).. controls +(-112:0.1)  and +(-56:0)..
			(12.98,11.71).. controls  +(117:0) and +(45:0.1)..
			(12.89,11.7)..  controls +(-139:0.1) and +(-120:0.1)..
			(12.62,11.68).. controls +(63:0) and +(-45:0).. 
			(12.58,11.77).. controls +(153:0) and +(180:0).. 
			(12.54,11.82).. controls +(0:0) and +(-34:0.1).. 
			(12.47,11.88).. controls  +(153:0.1) and +(0:0.1)..
			(12.24,11.94).. controls +(-172:0.1)  and +(0:0.1)..
			(11.97,11.93).. controls +(-173:0.1) and  +(90:0)..
			 (11.82,11.89).. controls +(-90:0.1) and +(180:0.1)..
			(12.31,11.36).. controls +(0:0)  and +(-165:0.2)..
			(12.79,11.46) --  cycle
			(13.72,11.69).. controls +(45:0)  and +(-56:0)..
			 (13.74,11.8).. controls +(108:0.1) and +(45:0.1)..
			(13.63,11.77).. controls  +(-129:0.1) and +(90:0)..
			 (13.55,11.66).. controls +(-90:0) and +(-129:0.1)..
			(13.72,11.69)  -- cycle
			(14.17,11.77).. controls +(135:0) and  +(27:0)..
			(14.11,11.77).. controls +(-162:0) and +(180:0).. 
			(14.13,11.75)..  controls +(0:0) and +(-45:0)..
			(14.17,11.77)  -- cycle
			(14.46,11.84).. controls +(56:0) and +(117:0).. 
			(14.51,11.86)..  controls +(-56:0.1) and +(-90:0.1)..
			 (14.68,11.93).. controls +(90:0) and +(0:0)..
			(14.64,11.99).. controls +(180:0)  and +(-104:0)..
			(14.63,12.06).. controls  +(76:0) and +(0:0)..
			(14.61,12.13)..  controls +(180:0) and +(80:0.1)..
			(14.56,12.02).. controls +(-90:0.1) and +(-8:0.1).. 
			(14.47,11.92).. controls  +(173:0.1) and +(82:0.1)..
			(14.32,11.81).. controls +(-90:0.1)  and +(-117:0.1)..
			 (14.46,11.84) -- cycle
			(14.85,12.12).. controls +(90:0.1)  and +(90:0)..
			 (14.78,12.14).. controls +(-90:0) and +(79:0.1)..
			(14.76,12.03).. controls +(-102:0.1)  and +(-131:0.1)..
			(14.8,12.01).. controls +(45:0) and  +(-90:0)..
			(14.85,12.12) --  cycle
			(15.68,12.15).. controls +(59:0.1) and +(0:0).. 
			(15.71,12.24).. controls  +(180:0) and +(90:0)..
			(15.49,12.05).. controls +(-90:0.1)  and +(-126:0.1)..
			(15.68,12.15) --  cycle
			 (12.17,12.28).. controls +(23:0.1) and +(24:0.1)..
			(11.73,12.28).. controls +(-162:0.1) and  +(174:0.1)..
			 (11.78,12.24).. controls +(-12:0.1) and +(-148:0.1)..
			(12.17,12.28) -- cycle 
			(15.88,12.33)..  controls +(68:0.1) and +(-63:0)..
			 (15.9,12.44).. controls +(135:0) and +(63:0.1)..
			(15.83,12.35).. controls  +(-124:0.1) and +(180:0)..
			(15.8,12.24).. controls +(0:0) and +(-121:0.1).. 
			(15.88,12.33) --  cycle
			(33.12,13.04).. controls +(127:0)  and +(-61:0.1)..
			(32.98,13.26)..  controls +(119:0.1) and +(-55:0.1)..
			(32.76,13.61).. controls +(124:0.3)  and +(-168:0.1)..
			(32.7,13.87).. controls +(22:0.1) and  +(-16:0.1)..
			(32.68,13.91).. controls +(157:0.1) and +(-107:0.1).. 
			(32.64,14.05)..  controls +(75:0.2) and +(-54:0.3)..
			(32.5,14.66).. controls +(124:0.6) and +(-48:0.7).. 
			 (31.16,16.35).. controls +(131:0.5) and +(-87:0.2)..
			(30.83,16.89).. controls +(96:0.1) and  +(0:0)..
			(30.82,17.07)..  controls +(0:0) and +(153:0.1)..
			 (30.93,17.03).. controls +(-18:0.1) and +(112:0.4)..
			 (31.28,16.49).. controls +(-70:0.1) and +(131:0.1)..
			(31.5,16.1).. controls  +(-46:0.3) and +(128:0.7)..
			(32.62,14.73)..  controls +(-54:0.3) and +(132:0.1)..
			(33.12,14.08).. controls +(-51:0.3) and +(90:0.3).. 
			(33.37,13.32).. controls  +(-90:0.2) and +(39:0.1)..
			(33.31,13.03).. controls  +(-143:0.1) and +(-49:0.1)..
			 (33.12,13.04) -- cycle
			(22.53,13.41).. controls +(125:0.2)  and +(-134:0.3)..
			 (22.66,13.82).. controls +(48:0.1) and +(34:0)..
			(22.79,13.99).. controls +(-162:0) and +(22:0.1).. 
			(22.55,13.9).. controls +(-158:0.2)  and +(-68:0.1)..
			(22.34,13.88).. controls +(108:0) and  +(-146:0)..
			(22.37,13.97).. controls +(34:0) and +(-5:0.1).. 
			 (22.24,14.01) -- (22.07,14.03) -- (22.28,14.03).. controls +(2:0.3) and +(-68:0.1)..
			(22.64,14.12)..  controls +(108:0) and +(2:0.3)..
			(22.13,14.15)..  controls +(-178:0.3) and +(-90:0)..
			(21.63,14.15).. controls  +(90:0.1) and +(-119:0.5)..
			(22.24,15.41).. controls +(66:0.1)  and +(-118:0.1)..
			(22.47,15.84).. controls +(62:1.1) and +(-155:0.4).. 
			 (23.66,17.39).. controls +(25:0.4) and +(180:1.1)..
			(25.31,17.58).. controls +(0:0.9)  and +(173:0.5)..
			(27.02,17.5).. controls +(-7:1.2) and  +(153:0.3)..
			(29.72,16.98).. controls +(-27:0.3) and +(104:0.3).. 
			(31.28,15.12).. controls  +(-76:0.1) and +(28:0.4)..
			(30.81,14.68).. controls +(-152:0.5) and  +(2:2)..
			 (28.33,14.37).. controls +(-178:1) and +(0:0.2)..
			(26.22,14.31).. controls +(-177:0.2) and  +(-27:0)..
			(25.86,14.32).. controls +(153:0) and +(-78:0.3).. 
			(25.72,14.87).. controls  +(101:0.6) and +(63:0)..
			(25.42,16.18).. controls  +(180:0) and +(102:0.3)..
			(25.51,15.66)..  controls +(-78:0.6) and +(69:0.1)..
			(25.67,14.34).. controls +(-104:0)  and +(-3:0.2)..
			(25.38,14.29).. controls +(180:0.2) and  +(-146:0)..
			(25.15,14.32).. controls +(31:0.1)  and +(0:0)..
			(25.16,14.35).. controls  +(166:0) and +(-127:0)..
			(25.15,14.41).. controls +(45:0) and  +(0:0)..
			(25.17,14.46)..  controls +(180:0) and +(90:0)..
			(24.91,14.14)..  controls +(-90:0.1) and +(-6:0.3)..
			(23.68,13.99) -- (23.48,14.01) -- (23.27,13.69).. controls +(-124:0.2) and  +(14:0)..
			 (22.98,13.35).. controls +(-162:0) and +(0:0.1)..
			(22.75,13.33).. controls +(180:0.1) and +(-56:0.1).. 
			(22.53,13.41) --  cycle
			 (25.1,14.48).. controls +(22:0.1) and +(-29:0.1)..
			(25.06,14.56).. controls +(153:0.1) and  +(20:0.1)..
			(24.83,14.56).. controls  +(-166:0.1) and +(90:0)..
			(24.68,14.49).. controls  +(-124:0) and +(-169:0.1)..
			(25.1,14.48) -- cycle 
			(25.14,14.6).. controls  +(135:0) and +(0:0)..
			 (25.08,14.64).. controls +(180:0) and +(135:0)..
			(25.07,14.6).. controls +(-63:0) and  +(180:0)..
			 (25.12,14.57).. controls +(0:0) and +(-63:0)..
			(25.14,14.6) -- cycle 
			(25.08,14.77).. controls +(90:0.1)  and +(45:0.1)..
			 (24.94,14.87).. controls +(-127:0.1) and +(129:0.1)..
			(24.95,14.76).. controls +(-49:0.1) and +(-90:0.1).. 
			(25.08,14.77) -- cycle 
			(24.83,14.87).. controls +(135:0) and +(18:0).. 
			(24.76,14.87).. controls +(-162:0) and +(180:0).. 
			(24.78,14.85).. controls +(0:0) and  +(-45:0)..
			(24.83,14.87) --  cycle
			(24.97,15.33).. controls  +(80:0.1) and +(-63:0)..
			(24.94,15.52).. controls +(112:0.1) and +(-102:0.1).. 
			(24.94,15.72).. controls +(75:0.1) and  +(-75:0.2)..
			(24.9,16.12).. controls +(108:0.3) and  +(7:0.1)..
			(24.73,16.37).. controls +(-176:0.1)  and +(86:0.1)..
			(24.46,16.12).. controls +(-100:0.1) and  +(180:0)..
			 (24.48,16.01).. controls +(0:0) and +(110:0.2)..
			 (24.65,15.61).. controls +(-73:0.4) and +(-174:0.1)..
			(24.87,15.21).. controls +(8:0.1)  and +(-96:0.1)..
			(24.97,15.33)  -- cycle
			(18.99,13.82).. controls  +(0:0) and +(180:0)..
			(19.11,13.82).. controls +(-14:0) and  +(0:0.1)..
			 (19.05,13.8).. controls +(180:0.1) and +(-162:0)..
			(18.99,13.82)  -- cycle
			(19.29,13.85)..  controls +(27:0) and +(162:0)..
			 (19.38,13.85).. controls +(-27:0) and +(0:0.1)..
			(19.33,13.83).. controls  +(180:0) and +(-153:0)..
			(19.29,13.85) -- cycle 
			(19.68,13.88).. controls +(27:0)  and +(162:0)..
			(19.77,13.88)..  controls +(-27:0) and +(0:0.1)..
			 (19.72,13.87).. controls +(180:0) and +(-162:0)..
			(19.68,13.88)  -- cycle
			(19.9,13.93).. controls +(11:0.1)  and +(-112:0.1)..
			 (20.03,14.04).. controls +(59:0.1) and +(90:0)..
			 (20.07,14.07).. controls +(-90:0.1) and +(0:0.1)..
			 (19.88,13.9).. controls +(180:0.1) and +(-164:0.1)..
			(19.9,13.93) --  cycle
			(9.46,14.31).. controls  +(90:0.2) and +(111:0.1)..
			 (9.56,14.44).. controls +(-74:0.1) and +(-80:0.1)..
			(9.58,14.42).. controls +(100:0.1) and +(-145:0.1).. 
			(9.7,14.59).. controls +(36:0.3) and  +(135:0)..
			(10.32,14.95)..  controls +(-45:0) and +(39:0.2)..
			 (9.7,14.32).. controls +(-144:0.1) and +(-56:0)..
			(9.58,14.29).. controls +(108:0)  and +(90:0)..
			 (9.57,14.27).. controls +(-76:0) and +(0:0)..
			(9.52,14.19).. controls +(-169:0.1)  and +(-90:0.1)..
			(9.46,14.31) --  cycle
			(9.74,14.27).. controls  +(48:0.3) and +(-140:0.6)..
			(11.29,15.66).. controls  +(40:2.1) and +(-170:1.1)..
			(15.12,17.86).. controls +(10:0.5)  and +(13:0.5)..
			(15.2,17.83).. controls +(-169:0.2)  and +(20:0.1)..
			(14.57,17.67).. controls  +(-162:0.3) and +(90:0.1)..
			(13.81,17.29).. controls +(-90:0) and +(33:0.2).. 
			 (13.45,17.02).. controls +(-148:0.9) and +(44:1.6)..
			(10.4,14.5).. controls +(-136:0.4) and +(-8:0.1).. 
			 (9.95,14.22).. controls +(173:0.1) and +(63:0)..
			(9.79,14.21).. controls +(-117:0)  and +(0:0)..
			(9.71,14.18).. controls +(180:0)  and +(-135:0.1)..
			 (9.74,14.27) -- cycle
			(10.46,14.77).. controls +(42:0.3) and  +(-140:0.3)..
			(11.33,15.54).. controls  +(41:0.7) and +(37:0.4)..
			(11.68,15.88).. controls +(-142:0.5) and +(44:0.5).. 
			(10.23,14.61).. controls  +(-135:0.6) and +(-136:0.8)..
			(10.46,14.77) --  cycle
			(11.98,16.08).. controls +(53:0.1) and  +(0:0)..
			(12.02,16.15).. controls  +(180:0) and +(45:0.1)..
			(11.94,16.08)..  controls +(-135:0) and +(180:0)..
			(11.9,16.01).. controls  +(0:0) and +(-143:0)..
			(11.98,16.08)  -- cycle
			(12.14,16.22)..  controls +(45:0.1) and +(39:0.1)..
			(12.13,16.23).. controls +(-143:0.1) and +(90:0).. 
			(12.07,16.17)..  controls +(-90:0) and +(-124:0.1)..
			(12.14,16.22)  -- cycle
			(10.67,14.3)..  controls +(37:0.1) and +(180:0)..
			 (10.8,14.35).. controls +(7:0.1) and +(-122:0.1)..
			(11.23,14.6).. controls  +(56:0) and +(-169:0.1)..
			(11.37,14.68).. controls +(22:0.1)  and +(-104:0)..
			(11.48,14.78).. controls  +(72:0.1) and +(143:0.1)..
			(11.56,14.81)..  controls +(-45:0) and +(-149:0.1)..
			(11.67,14.81).. controls +(39:0.1) and +(11:0.1).. 
			(11.67,14.84).. controls +(-162:0) and +(-63:0).. 
			(11.59,14.85).. controls  +(117:0) and +(180:0)..
			(11.6,14.88).. controls +(-27:0)  and +(-34:0)..
			(11.86,15.14)..  controls +(135:0) and +(162:0)..
			(11.89,15.14).. controls +(0:0) and  +(-127:0)..
			(11.99,15.2).. controls  +(53:0.1) and +(162:0)..
			(12.09,15.26).. controls +(-27:0) and +(-153:0).. 
			(12.16,15.27).. controls +(63:0) and  +(0:0)..
			(12.14,15.31).. controls +(180:0)  and +(-90:0)..
			(12.1,15.34).. controls +(90:0)  and +(180:0)..
			(12.15,15.38)..  controls +(0:0) and +(-90:0)..
			(12.21,15.42).. controls +(90:0)  and +(162:0)..
			(12.26,15.44).. controls  +(-22:0.1) and +(-59:0.1)..
			(12.3,15.5).. controls +(117:0) and  +(180:0)..
			(12.46,15.66).. controls  +(0:0) and +(-131:0.1)..
			(12.63,15.8).. controls  +(41:0.1) and +(-166:0)..
			(12.84,15.96).. controls +(18:0) and +(-108:0).. 
			(12.93,16.03)..  controls +(72:0) and +(180:0)..
			(13.01,16.08).. controls +(0:0) and +(-108:0).. 
			(13.1,16.14).. controls  +(72:0) and +(153:0)..
			 (13.16,16.17).. controls +(-34:0) and +(-72:0.1)..
			(13.18,16.23).. controls +(101:0.1) and +(117:0).. 
			 (13.18,16.27).. controls +(-72:0) and +(180:0)..
			(13.25,16.23)..  controls +(0:0) and +(-45:0)..
			(13.28,16.26)..  controls +(117:0) and +(180:0)..
			(13.34,16.29).. controls +(-8:0.1)  and +(-90:0.1)..
			(13.42,16.37).. controls +(104:0) and +(117:0).. 
			(13.43,16.41)..  controls +(-72:0) and +(-135:0.1)..
			 (13.51,16.43).. controls +(49:0.1) and +(159:0.1)..
			(13.91,16.56).. controls +(-27:0) and  +(-117:0)..
			 (14,16.6).. controls +(69:0.1) and +(153:0)..
			(14.16,16.73)..  controls +(-27:0) and +(-121:0.1)..
			(14.24,16.79).. controls +(54:0.1)  and +(90:0.1)..
			(14.39,16.82).. controls +(-90:0) and  +(180:0)..
			 (14.43,16.79).. controls +(0:0.1) and +(-23:0.1)..
			(14.39,16.9).. controls +(153:0) and  +(180:0.1)..
			(14.52,16.92)..  controls +(0:0.2) and +(-162:0.1)..
			 (14.99,16.99).. controls +(24:0.1) and +(-14:0.1)..
			(14.96,17.06).. controls +(165:0.1) and +(-171:0.4).. 
			(15.25,17.16).. controls +(7:0.2) and +(-90:0).. 
			(15.66,17.25)..  controls +(90:0) and +(6:0.1)..
			(15.51,17.27).. controls  +(-174:0.3) and +(-165:0.3)..
			 (15.56,17.39).. controls +(13:0.2) and +(-117:0)..
			(15.88,17.5).. controls  +(59:0.1) and +(-177:1.1)..
			(17.37,17.64).. controls +(3:0.9)  and +(176:1.1)..
			(21.77,17.6)  -- (22.46,17.56) -- (22.63,17.72).. controls +(45:0.1) and +(-157:0.1)..
			(22.93,17.95)..  controls +(29:0.2) and +(146:0)..
			(23.49,18.12)..  controls +(-63:0) and +(0:0)..
			(23.46,18.09).. controls  +(180:0.1) and +(43:0.2)..
			(22.86,17.75).. controls +(-137:0.5) and  +(62:2.5)..
			(20.94,14.59) -- (20.77,14.28) --  (20.04,14.3).. controls +(178:2.6) and +(5:1)..
			(10.83,14.26).. controls +(-175:0.2) and +(-140:0.1).. 
			(10.67,14.3) -- cycle 
			(12.82,14.5).. controls +(3:0.4) and +(-178:0.3).. 
			(14.2,14.56) -- (14.8,14.58) -- (14.98,15.08)..  controls +(70:0.8) and +(153:0)..
			(15.64,16.59).. controls  +(-27:0) and +(-72:0)..
			(15.65,16.61).. controls +(117:0) and  +(-122:0.1)..
			(15.71,16.8).. controls +(66:0.1) and  +(-45:0)..
			(15.77,16.98)..  controls +(180:0) and +(7:0.2)..
			(15.44,16.96).. controls +(-177:0.2) and +(0:0.1).. 
			(14.95,16.93).. controls +(-174:0.1)  and +(90:0)..
			(14.78,16.89).. controls  +(-90:0) and +(11:0.1)..
			(14.7,16.84).. controls +(-163:0.1)  and +(34:0.1)..
			 (14.11,16.58).. controls +(-141:0.1) and +(-8:0.1)..
			(13.91,16.53).. controls +(170:0.1)  and +(45:0)..
			 (13.77,16.5).. controls +(-127:0) and +(180:0)..
			(13.79,16.45).. controls +(0:0) and  +(-63:0)..
			(13.84,16.48)..  controls +(135:0) and +(180:0)..
			(13.86,16.51).. controls  +(0:0.1) and +(24:0.1)..
			(13.84,16.4).. controls +(-158:0.1) and +(0:0).. 
			(13.74,16.37).. controls +(135:0)  and +(45:0.1)..
			(13.64,16.3)..  controls +(-135:0.1) and +(-18:0)..
			(13.51,16.24).. controls +(146:0) and  +(108:0)..
			(13.49,16.23).. controls +(-63:0) and  +(22:0.1)..
			(13.43,16.15).. controls  +(-153:0) and +(36:0.1)..
			(13.15,15.96).. controls  +(-144:0.1) and +(37:0.1)..
			(12.88,15.77).. controls +(-149:0.1) and +(34:0).. 
			(12.75,15.67).. controls  +(-135:0) and +(-90:0)..
			(12.7,15.65).. controls +(90:0) and +(39:0.1).. 
			(12.62,15.59)..  controls +(-133:0.2) and +(-14:0)..
			(12.06,15.16)..  controls +(180:0) and +(135:0)..
			(12.04,15.14).. controls +(-90:0) and +(42:0.2).. 
			(11.83,14.91).. controls +(-136:0.3) and +(-45:0.1).. 
			(11.55,14.75).. controls  +(146:0.1) and +(130:0.1)..
			(11.54,14.73).. controls  +(-56:0.1) and +(40:0.2)..
			(11.45,14.54) -- (11.31,14.42) -- (11.68,14.44).. controls +(3:0.2) and +(-177:0.4).. 
			(12.82,14.5) -- cycle 
			(15.12,14.84).. controls +(77:0.1) and  +(-121:0.1)..
			(15.23,15.17).. controls +(63:0) and +(-112:0.3).. 
			(15.45,15.71)..  controls +(67:0.3) and +(-104:0)..
			(15.66,16.22).. controls  +(58:0.1) and +(-90:0)..
			(15.88,16.7)..  controls +(90:0) and +(0:0)..
			(15.85,16.72).. controls +(180:0.1) and  +(78:0.3)..
			(15.08,14.91).. controls +(-101:0.1) and +(59:0.1).. 
			(14.97,14.63).. controls +(-135:0) and +(180:0).. 
			(15,14.6).. controls  +(0:0) and +(-106:0.1)..
			 (15.12,14.84) -- cycle
			(17.12,15.04)  -- (17.12,15.18) -- (17.15,15.03).. controls +(-81:0.1) and +(-97:0.1)..
			 (17.18,14.97).. controls +(90:0.1) and +(124:0)..
			(17.21,15.03).. controls  +(-53:0) and +(-104:0)..
			(17.26,15.04).. controls +(56:0) and  +(153:0)..
			(17.32,15.08).. controls +(-27:0) and +(-102:0.1).. 
			(17.39,15.23).. controls +(78:0.6) and  +(-157:0.2)..
			(17.74,16.05).. controls +(20:0.1) and  +(160:0.1)..
			(18.97,16.07).. controls +(-11:0.1) and  +(-129:0.1)..
			(19.1,16.12).. controls +(45:0.1) and  +(166:0)..
			(19.23,16.17).. controls +(-14:0) and +(-121:0.1).. 
			(19.34,16.23) -- (19.38,16.31) -- (19.42,16.21) -- (19.47,16.11) -- (19.48,16.22).. controls +(80:0.1) and  +(27:0)..
			(19.46,16.31)..  controls +(-135:0) and +(-104:0)..
			(19.45,16.35).. controls +(72:0) and +(-90:0).. 
			 (19.47,16.44).. controls +(90:0) and +(153:0)..
			 (19.51,16.46).. controls +(-45:0) and +(-56:0)..
			(19.57,16.78)..  controls +(135:0) and +(-7:0.3)..
			(18.99,16.87).. controls +(173:0.7) and +(38:0.1).. 
			(17.59,16.9).. controls +(-140:0.1) and +(176:0.3).. 
			(17.89,16.82).. controls +(-2:0.4)  and +(96:0.1)..
			(18.27,16.71)..  controls +(-86:0.1) and +(23:0.1)..
			(18.14,16.31).. controls +(-162:0) and  +(2:0.3)..
			(17.46,16.27).. controls +(-179:0.4)  and +(90:0)..
			(16.83,16.21)..  controls +(-90:0) and +(59:0.1)..
			(16.78,16.08).. controls +(-119:0.2) and  +(94:0.1)..
			(16.59,15.4)..  controls +(-90:0.1) and +(56:0)..
			(16.56,15.23).. controls  +(-135:0) and +(84:0.1)..
			(16.5,15.03) --  (16.48,14.88) -- (16.8,14.89) -- (17.11,14.89) -- cycle
			(17.32,14.94).. controls  +(90:0) and +(27:0)..
			(17.29,14.97).. controls +(-153:0) and +(90:0).. 
			(17.25,14.92).. controls  +(-90:0) and +(180:0)..
			(17.29,14.88).. controls  +(0:0) and +(-90:0)..
			(17.32,14.94) -- cycle 
			 (11.87,15.03).. controls +(63:0) and +(18:0)..
			(11.84,15.04).. controls +(-162:0.1)  and +(180:0.1)..
			(11.82,14.99).. controls  +(0:0) and +(-117:0)..
			 (11.87,15.03) -- cycle
			(19.17,15.96).. controls +(108:0)  and +(45:0)..
			(19.13,16).. controls  +(-135:0) and +(117:0)..
			 (19.14,15.95).. controls +(-63:0.1) and +(-72:0.1)..
			(19.17,15.96) --  cycle
			(13.67,16.44).. controls +(63:0) and +(0:0).. 
			(13.64,16.47).. controls  +(180:0) and +(90:0)..
			(13.58,16.44).. controls +(-90:0) and  +(180:0)..
			(13.62,16.4).. controls  +(0:0) and +(-117:0)..
			 (13.67,16.44) -- cycle
			 (15.14,17.07).. controls +(90:0) and +(0:0)..
			(15.1,17.11).. controls +(180:0) and  +(117:0)..
			(15.08,17.07).. controls  +(-63:0) and +(180:0)..
			(15.12,17.04).. controls +(0:0) and  +(-90:0)..
			(15.14,17.07)  -- cycle
			(16.08,17.52)..  controls +(135:0) and +(18:0)..
			(16.01,17.52).. controls +(-162:0) and  +(180:0)..
			(16.03,17.5)..  controls +(0:0) and +(-63:0)..
			(16.08,17.52) --  cycle
			 (31.47,14.83).. controls +(0:0) and +(-135:0)..
			(31.53,14.89).. controls +(34:0.1) and +(45:0.1).. 
			(31.54,14.88).. controls +(-124:0.1) and +(-90:0).. 
			 (31.47,14.83) -- cycle
			(31.31,15.42) -- (31.24,15.5) -- (31.32,15.43)..  controls +(-37:0) and +(0:0)..
			(31.4,15.36).. controls +(-90:0) and  +(-49:0.1)..
			 (31.31,15.42) -- cycle
			(22.38,16.07).. controls +(62:1) and +(-153:0.4).. 
			(23.54,17.55)..  controls +(26:0.3) and +(180:0.3)..
			(24.5,17.74).. controls +(-3:0.2)  and +(6:0.3)..
			(24.4,17.7).. controls  +(-175:0.9) and +(60:1.3)..
			(22.53,16.25).. controls +(-120:0.3) and  +(27:0)..
			(22.2,15.7).. controls +(-135:0) and  +(-117:0.2)..
			(22.38,16.07)  -- cycle
			(31.04,15.76).. controls +(129:0.1) and  +(141:0.1)..
			(31.05,15.77)..  controls +(-40:0.1) and +(0:0)..
			(31.1,15.7).. controls +(0:0)  and +(-53:0.1)..
			(31.04,15.76) -- cycle 
			(30.72,16.15).. controls  +(129:0.1) and +(146:0.1)..
			(30.74,16.16).. controls  +(-40:0.1) and +(0:0)..
			(30.79,16.08).. controls  +(180:0) and +(-53:0.1)..
			 (30.72,16.15) -- cycle
			(30.46,16.45)  -- (30.39,16.52) -- (30.47,16.46).. controls +(-39:0.1) and +(0:0)..
			(30.55,16.38).. controls +(-90:0) and +(-53:0.1).. 
			(30.46,16.45) -- cycle 
			 (30.16,16.75).. controls +(129:0.1) and +(141:0.1)..
			(30.17,16.76).. controls  +(-40:0.1) and +(0:0)..
			(30.22,16.68).. controls +(180:0) and +(-53:0.1).. 
			(30.16,16.75) --  cycle
			(29.88,16.99).. controls +(130:0.1) and  +(146:0.1)..
			(29.89,17.01).. controls +(-40:0.1)  and +(0:0)..
			 (29.94,16.93).. controls +(180:0) and +(-56:0)..
			 (29.88,16.99) -- cycle
			(14.88,17.03).. controls +(18:0) and  +(135:0)..
			(14.95,17.02).. controls +(-45:0)  and +(0:0)..
			(14.9,17).. controls +(166:0) and  +(-146:0)..
			(14.88,17.03) -- cycle 
			(29.49,17.15).. controls +(163:0.1)  and +(-12:0.2)..
			(28.95,17.29).. controls +(169:0.2)  and +(-45:0)..
			(28.58,17.38).. controls +(135:0.1) and +(152:0.2).. 
			(29.66,17.13).. controls  +(-28:0.1) and +(-18:0.3)..
			(29.49,17.15)  -- cycle
			 (28.43,17.41).. controls +(27:0) and +(153:0)..
			(28.51,17.41)..  controls +(-18:0) and +(0:0.1)..
			(28.47,17.39).. controls +(180:0) and  +(-162:0)..
			(28.43,17.41) --  cycle
			(28.21,17.45).. controls  +(18:0) and +(153:0)..
			 (28.3,17.45).. controls +(-18:0) and +(0:0)..
			(28.26,17.43).. controls +(180:0.1) and +(-153:0).. 
			 (28.21,17.45) -- cycle
			 (28,17.48).. controls +(18:0) and +(153:0)..
			(28.09,17.48).. controls +(-18:0) and  +(0:0.1)..
			(28.05,17.47).. controls +(180:0.1) and  +(-153:0)..
			(28,17.48) --  cycle
			 (27.77,17.52).. controls +(14:0) and +(153:0)..
			(27.88,17.52).. controls +(-27:0)  and +(0:0.1)..
			(27.82,17.5).. controls +(180:0.1)  and +(-162:0)..
			(27.77,17.52) --  cycle
			(27.53,17.55).. controls +(18:0)  and +(162:0)..
			(27.65,17.55).. controls +(0:0) and  +(0:0.1)..
			 (27.59,17.54).. controls +(180:0.1) and +(180:0)..
			 (27.53,17.55) -- cycle
			(27.26,17.59).. controls +(11:0.1)  and +(162:0)..
			(27.4,17.59).. controls  +(-14:0) and +(0:0.1)..
			 (27.32,17.57).. controls +(180:0.1) and +(-166:0)..
			 (27.26,17.59) -- cycle
			(26.96,17.63).. controls +(0:0)  and +(180:0)..
			(27.09,17.63)..  controls +(-18:0) and +(0:0.1)..
			(27.02,17.61).. controls  +(180:0.1) and +(-162:0)..
			(26.96,17.63) -- cycle 
			(26.63,17.66).. controls +(14:0) and  +(162:0)..
			(26.77,17.66).. controls +(-18:0) and +(-7:0.1).. 
			(26.69,17.65)..  controls +(180:0.1) and +(-169:0.1)..
			 (26.63,17.66) -- cycle
			(25.42,17.71).. controls +(135:0) and +(180:0.2).. 
			(25.7,17.74).. controls +(-2:0.3) and +(3:0.4).. 
			(25.65,17.7)..  controls +(-174:0.1) and +(-18:0)..
			(25.42,17.71)  -- cycle
			(26.24,17.7)..  controls +(0:0) and +(180:0)..
			(26.38,17.7).. controls +(-18:0) and  +(0:0.1)..
			(26.3,17.68).. controls +(180:0.1)  and +(-166:0)..
			 (26.24,17.7) -- cycle
			 (25,17.77).. controls +(8:0.1) and +(170:0.1)..
			(25.23,17.77).. controls +(-8:0.1) and  +(0:0.1)..
			(25.12,17.75).. controls +(180:0.1) and  +(-170:0.1)..
			(25,17.77) -- cycle 
			(16.3,17.95).. controls +(37:0.1) and +(180:0.1).. 
			(16.94,18.06).. controls +(0:0.1)  and +(8:0.2)..
			(16.7,17.99).. controls +(-171:0.1) and  +(10:0.1)..
			(16.35,17.93).. controls +(-172:0.1)  and +(-143:0)..
			 (16.3,17.95) -- cycle
			(17.17,18.05)..  controls +(11:0.1) and +(166:0)..
			(17.33,18.05).. controls +(-11:0.1)  and +(0:0.1)..
			(17.25,18.03).. controls +(180:0.1) and +(-166:0).. 
			(17.17,18.05) --  cycle
			(17.53,18.09).. controls  +(0:0.1) and +(170:0.1)..
			 (17.74,18.08).. controls +(0:0.1) and +(0:0.1)..
			(17.62,18.07).. controls +(180:0.1)  and +(-172:0.1)..
			(17.53,18.09) -- cycle 
			 (18.26,18.13).. controls +(4:1.3) and +(-179:1.3)..
			(21.26,18.26)..  controls +(1:0.8) and +(-177:0.2)..
			(22.97,18.29).. controls +(2:0.3) and  +(9:0.5)..
			 (22.74,18.22).. controls +(-171:0.6) and +(-161:0.2)..
			(22.16,18.17).. controls +(14:0.1)  and +(4:0.6)..
			(19.55,18.17).. controls  +(-176:0.3) and +(2:0.3)..
			(18.43,18.12).. controls +(-178:0.6)  and +(-176:0.4)..
			 (18.26,18.13) -- cycle
			(21.41,18.12).. controls +(10:0.1) and  +(172:0.1)..
			(21.65,18.12).. controls  +(-7:0.1) and +(0:0.1)..
			 (21.54,18.11).. controls +(180:0.1) and +(-172:0.1)..
			(21.41,18.12)  -- cycle
			(20.77,18.15)..  controls +(18:0) and +(153:0)..
			(20.86,18.15).. controls  +(-27:0) and +(0:0)..
			 (20.82,18.14).. controls +(180:0.1) and +(-153:0)..
			(20.77,18.15) -- cycle; 
	\end{tikzpicture}
	 }
	\noindent 
	đồng.  Đến hết đợt thứ nhất anh Bình có tất cả $624$ triệu đồng.\\
	\textbf{Đợt thứ hai kế tiếp}:  do muốn rút ngắn thời gian mua xe thì số tiền còn lại anh tiếp tục tích góp với tháng đầu là $5$ triệu đồng và mỗi tháng tiếp theo số tiền gấp đôi tháng kề trước nó. Xét tính đúng sai của các khẳng định sau
	\choiceTF 
		{\True Đợt thứ nhất  anh Bình tích lũy tiền theo dãy số với cấp số cộng có công sai là $d=2$ triệu và $u_1=3$ triệu}
		{Anh Bình tích lũy tiền hết  đợt thứ nhất trong $25$ tháng}
		{\True Đợt thứ hai anh Bình tích lũy tiền theo dãy số với  cấp số nhân có công bội là $q=2$ triệu và $u_1=5$ triệu}
		{\True Để đủ tiền mua ô tô thì anh Bình thì anh Bình tích góp ít  nhất $31$ tháng}
	 \loigiai{
		 \textbf{Đợt thứ nhất:}\\
		Dãy số tiền tích lũy theo tháng là  một cấp số cộng với công sai $d=2$ triệu.
		Theo đề bài, tháng thứ $10$ góp $21$ triệu, ta có $u_{10} = u_1 + 9d \Rightarrow 21  = u_1+9\cdot 2 \Rightarrow u_1 = 3$ triệu.\\
		Tổng số tiền đợt một  là $624$ triệu. Số tháng tích lũy $n$ được tính như sau\\
		$S_n=624  \Leftrightarrow 624 = \dfrac{n[2u_1+(n-1)\cdot 2]}{2} \Rightarrow n=24$ tháng.\\
		\textbf{Đợt  thứ hai:}\\
		Số  tiền cần tích lũy thêm là $1259 - 624 = 635$ triệu đồng.\\
		Dãy số tiền tích lũy theo tháng là một cấp số nhân có số hạng đầu $v_1=5$ triệu và  công bội $q=2$.\\
		Số tháng $m$ cần tích lũy thêm là $635  = 5\cdot \dfrac{1-2^m}{1-2} \Rightarrow m=7$ tháng.\\
		\textbf{Tổng thời  gian:}\\
		Tổng cộng hai đợt cần có ít nhất  $24+7=31$ tháng.
		\begin{itemchoice}
			\itemch  Đợt thứ nhất là cấp số cộng với công sai $d=2$ triệu và số hạng đầu $u_1=3$ triệu.
			 \itemch Anh Bình tích lũy hết đợt thứ nhất trong $24$ tháng.
			\itemch Đợt thứ hai là cấp số nhân  có công bội $q=2$ và số hạng đầu là $5$ triệu.
			\itemch Để đủ tiền mua ô tô, anh Bình cần tích góp ít nhất $31$ tháng. 
		\end{itemchoice}
	 }
\end{ex}

% nguồn: Ok_dungsai_lop11_Lan2
\begin{ex}%[1D2V2-7]%[ma:bd6a67a4]
	Một sinh viên sau khi ra trường và xin vào làm cho một trung tâm với mức lương khởi điểm là $100$ triệu đồng một năm. Cứ sau mỗi năm, trung tâm trả thêm cho sinh viên $20$ triệu đồng. Gọi $u_n$ (triệu đồng) là số tiền lương mà sinh viên đó nhận được  ở năm thứ $n$. Xét tính đúng sai của các khẳng định sau
	 \choiceTF
		{\True Số tiền lương sinh viên nhận được ở năm thứ hai là  $120$ triệu đồng}
		{Số tiền  lương sinh viên nhận được ở năm thứ $10$ là $300$ triệu đồng}
		{Dãy số $(u_n)$ là cấp số cộng có  $u_1=120$ và công sai $d=20$}
		{Giả sử, mỗi năm bạn sinh viên chi tiêu tiết kiệm hết $70$ triệu đồng. Vậy sau ít nhất $12$ năm  thì sinh viên đó mua được căn chung cư $2$ tỉ đồng}
	\loigiai{ 
		Số tiền lương mỗi năm lập thành một cấp số  cộng $(u_n)$ có số hạng đầu $u_1=100$ và công sai $d=20$.\\
		Công thức số hạng tổng quát là $u_n  = u_1+(n-1)d=100+(n-1)\cdot 20=20n+80$.\\
		Số tiền lương năm thứ hai là  $u_2 = 20\cdot 2+80=120$ (triệu đồng).\\
		 Số tiền lương năm thứ $10$ là $u_{10} = 20\cdot 10+80=280$ (triệu đồng).\\
		Tổng  số tiền sinh viên tiết kiệm được sau $n$ năm (lấy tổng lương trừ đi tiền chi tiêu) là
		\[S = \dfrac{n}{2}[2u_1+(n-1)d] - 70n = \dfrac{n}{2}[2\cdot 100+(n-1)\cdot 20] - 70n =  10n^2+20n\, \text{(triệu đồng)}.\]
		 Để mua được căn chung cư giá $2$ tỉ đồng ($2\,000$ triệu), ta cần có
		 \[S \ge 2\,000 \Rightarrow 10n^2+20n - 2\,000 \ge 0 \Leftrightarrow \heva{&n \ge 13{,}1\\& n \le -15{,}1}.\]
		Vì $n$ là số năm nên $n>0$, do đó ta cần $n \ge 13{,}1$. Vậy sinh viên cần ít nhất $14$ năm để đủ tiền mua nhà. 
		\begin{itemchoice}
			\itemch Số tiền  lương năm thứ hai là $120$ triệu đồng.
			\itemch Số tiền lương năm thứ  $10$ là $280$ triệu đồng.
			\itemch  Cấp số cộng có $u_1=100$.
			\itemch  Cần ít nhất $14$ năm để mua được căn chung cư.
		\end{itemchoice}
	 }
\end{ex}

% nguồn: Ok_dungsai_lop11_Lan2
\begin{ex}%[1D2V2-7]%[ma:34c0f660]
	Khi kí kết hợp đồng lao động với người lao động, một doanh nghiệp để xuất hai phương  án trả lương như sau: Phương án $1$: Năm thứ nhất, tiền lương là $120$ triệu. Kể từ năm thứ hai trở đi, mỗi năm tiền lương được tăng $18$ triệu. Phương án $2$: Quý thứ nhất, tiền lương là $24$ triệu. Kể từ quý thứ hai trở đi, mỗi quý tiền lương được tăng $1{,}8$ triệu. Các mênh đề sau đúng hay sai?
	\choiceTF 
		{\True Trong phương  án $1$: dãy số tiền lương là cấp số cộng có số hạng đầu tiên là $u_1=120$, công sai $d_1=18$}
		{Trong phương  án $1$: tiền lương người lao động nhận được trong năm thứ ba là $174$ triệu}
		{\True Tổng tiền lương người lao động nhận được trong ba năm là $414$  triệu}
		{\True Nếu kí hợp đồng lao động trong ba năm, với mong muốn nhận được tổng số tiền lương cao  nhất thì người lao động nên chọn phương án $1$}
	 \loigiai{
		\begin{itemchoice}
			\itemch  \textbf{Đúng} \\
			Dãy số tiền lương là cấp số cộng có số hạng đầu tiên  là $u_1=120$, công sai $d_1=18$.
			\itemch \textbf{Sai}  \\
			Ta có $u_3=120+2\cdot 18=156$ (triệu). 
			 \itemch \textbf{Đúng} \\
			Trong phương án $1$: Tổng tiền lương người lao động nhận được trong ba năm  là
			\[S_3=\dfrac{\left(u_1+u_3\right) \cdot  3}{2}=\dfrac{(120+156) \cdot 3}{2}=414\text{ (triệu)}.\]
			\itemch \textbf{Đúng}  \\
			Trong phương án $2$: Dãy số tiền lương là cấp số cộng có số hạng đầu tiên là $v_1=24$, công sai $d_2=1{,}8$, lương  tăng theo quý. Ba năm tương ứng với $12$ quý.\\
			 Ta có $v_{12}=24+11 \cdot 1{,}8=43{,}8$ triệu.\\
			Tổng số tiền lương nhận được sau $3$  năm tương ứng với $12$ quý là
			\[S_{12}=\dfrac{\left(v_1+v_{12}\right)  \cdot 12}{2}=(24+43{,}8) \cdot 6=406{,}8 \text { (triệu)}.\]
			Ta  nhận thấy $406{,}8<414$.\\
			Vậy nếu kí hợp đồng lao động trong ba năm, với mong muốn nhận được tổng số tiền lương cao nhất thì người lao động nên  chọn phương án $1$.
		\end{itemchoice}
	 }
\end{ex}

% nguồn: Ok_dungsai_lop11_Lan2
\begin{ex}%[1D2V2-7]%[ma:6fd6582a]
	Một nhà hát có $25$ hàng ghế với  $16$ ghế ở hàng thứ nhất, $18$ ghế ở hàng thứ hai, $20$ ghế ở hàng thứ ba và cứ tiếp tục theo quy luật đó, tức là hàng sau nhiều hơn hàng liền trước nó $2$ ghế. Gọi $u_n$ (ghế) là tổng số ghế ở hàng thứ $n$.
	\choiceTF 
		{\True Dãy số $\left(u_n\right)$ là cấp số cộng có công  sai $d=2$}
		{\True $u_2=18$} 
		{Số ghế ở hàng thứ  $20$ nhỏ hơn $54$}
		{Tổng số ghế trong nhà hát nhiều hơn $1\,000$} 
	\loigiai{ 
		\begin{itemchoice}
			 \itemch \textbf{Đúng} \\
			Số ghế trong các dãy ghế liên tiếp lập thành một cấp số  cộng với số hạng đầu $u_1=16$ và công sai $d=2$.
			\itemch \textbf{Đúng} \\ 
			Ta  có $u_2=u_1+d=16+2=18$.
			\itemch \textbf{Sai}  \\
			Ta có $u_{20}=u_1+19 d=16+19\cdot 2=54$. 
			\itemch \textbf{Sai} \\ 
			Số ghế ở mỗi hàng của nhà hát lập thành một cấp số cộng, gồm $25$ số hạng, với số hạng đầu $u_1=16$ và công sai $d=2$.\\ Tổng các số hạng này là 
			\[S_{25}=u_1+u_2+\ldots+u_{25}=\dfrac{25}{2}\left[2 u_1+(25-1)  d\right]=\dfrac{25}{2}(2\cdot 16+24\cdot 2)=1\,000.\]
		\end{itemchoice}
	} 
\end{ex}


%%%%%%%%%%%%%%%% TRẢ LỜI NGẮN %%%%%%%%%%%%%%%%

% nguồn: Ok_TraLoiNgan_Lop11_Lan1
\begin{ex}%[1D2H2-7]%[ma:854e3fa2]
	Trong một khán phòng có tất cả $30$  dãy ghế, dãy đầu tiên có $15$ ghế, các dãy liền sau nhiều hơn dãy trước đó $4$ ghế, hỏi khán phòng đó có tất cả bao nhiêu ghế?
	 \shortans{2190}
	 \loigiai{
		Gọi $u_1, u_2, \ldots, u_{30}$ lần lượt là số ghế của dãy ghế  thứ nhất, dãy ghế thứ hai,..., dãy ghế thứ ba mươi. \\
		Khi đó, $\left(u_n \right)$ là một cấp số cộng có số hạng đầu $u_1=15$, công sai $d=4$ (trong đó $1\le n\le  30$). \\
		Gọi $S_{30}$ là  tổng số ghế trong khán phòng.\\
		Ta có $S_{30}=u_1+u_2+\cdots+u_{30}=\dfrac{30}{2}  \left[2u_1+(30-1)d\right]=15(2\cdot15+29\cdot4)=2190$.
	} 
\end{ex}

% nguồn: Ok_TraLoiNgan_Lop11_Lan2
\begin{ex}%[1D2H2-7]%[ma:4f9178a7]
	Một cơ sở khoan giếng đưa ra định mức giá như sau Giá của mét khoan đầu tiên là $100$ nghìn đồng và kể từ mét khoan thứ hai, giá của mỗi mét sau tăng thêm $30$ nghìn đồng so với giá của mét  khoan ngay trước đó. Một người cần khoan một giếng sâu $20$ m để lấy nước dùng cho sinh hoạt của gia đình. Hỏi sau khi hoàn thành việc khoan giếng, gia đình đó phải thanh toán cho cơ sở khoan giếng số tiền bao nhiêu nghìn đồng?
	\shortans{7,70} 
	 \loigiai{
		Gọi  $u_n$ là giá của mét khoan thứ $n$, trong đó $1 \le n \le 20$.\\
		Khi đó, $(u_n)$ là cấp số cộng có số hạng đầu $u_1=100$ và  công sai $d=30$.\\
		Số tiền mà gia đình phải thanh toán cho cơ sở khoan  giếng là\\
		$S_{20}=u_1+u_2+\ldots+u_{20}=\dfrac{20(2u_1+19d)}{2}=\dfrac{20(2\cdot 100+19\cdot 30)}{2}=7\,700$  (nghìn đồng).
	} 
\end{ex}

% nguồn: Ok_TraLoiNgan_Lop11_Lan2
\begin{ex}%[1D2H2-7]%[ma:e2510d6f]
	 \immini[thm]
	{Một bức tường trang trí có dạng hình thang, rộng $2{,}4$ m ở đáy và rộng $1{,}2$ m ở đỉnh (hình vẽ bên). Các viên gạch hình vuông có kích thước $10$ cm $\times$ $10$ cm phải được đặt sao cho mỗi hàng ở phía trên chứa ít hơn một viên so với hàng ở ngay phía dưới nó. Hỏi sẽ cần bao nhiêu viên gạch hình vuông như vậy để ốp hết bức tường đó?} 
	{\begin{tikzpicture}[>=stealth,line join=round,line cap=round,font=\small,scale=0.25]
		\def\bot  tomwidth{240}
		\def\to pwidth{120} 
		\def\tilewidth{10} 
		 \def\tileheight{10}
		 \def\numrows{13}
		\def\sf{0.1} 
		\coordinate  (A) at (\sf*\bot tomwidth/2,0);
		\coordinate (B) at (-\sf*\bot tomwidth/2,0); 
		 \coordinate (C) at (-\sf*\to pwidth/2,\sf*\numrows*\tileheight);
		 \coordinate (D) at (\sf*\to pwidth/2,\sf*\numrows*\tileheight);
		\draw[line width=1mm,fill=gray!20,opacity=0.5] (A) -- (B) -- (C) -- (D)  -- cycle;
		\foreach \row  [count=\i] in {1,...,\numrows}{
			\pgfmathsetmacro{\numtiles}{\bot  tomwidth/\tilewidth - (\i-1)}
			\pgfmathsetmacro{\currentwidth}{\numtiles  * \tilewidth}
			 \pgfmathsetmacro{\offset}{(\bot tomwidth - \currentwidth) / 2}
			 \foreach \tile [count=\j] in {1,...,\numtiles}{
			\pgfmathsetmacro{\xpos}{\sf*(\offset  + (\j-1)*\tilewidth - \bot tomwidth/2)}
			 \pgfmathsetmacro{\ypos}{\sf*((\row-1)*\tileheight)}
			\draw[fill=white,draw=black!50] (\xpos,\ypos) rectangle (\xpos +  \sf*\tilewidth,\ypos + \sf*\tileheight);
		} 
		} 
		\draw[<->] (-\sf*\bot tomwidth/2,-1) -- (\sf*\bot tomwidth/2,-1)  node[midway,below] {2,4 m};
		\draw[<->] (-\sf*\to pwidth/2,\sf*\numrows*\tileheight + 1) --  (\sf*\to pwidth/2,\sf*\numrows*\tileheight + 1) node[midway,above] {1,2 m};
		\draw[<->] (\sf*\bot tomwidth/2 + 1,0) --  (\sf*\bot tomwidth/2 + 1,\sf*\numrows*\tileheight) node[midway,right] {1,3 m};
	\end{tikzpicture}
	 }
	\shortans{234} 
	\loigiai{ 
		Đổi $2{,}4$  m $= 240$ cm; $1{,}2$ m $= 120$ cm.\\
		Số viên gạch ở hàng  đầu tiên (ứng với đáy lớn) là $u_1=240:10=24$.\\
		Số viên gạch ở  hàng trên cùng (ứng với đáy nhỏ) là $u_n=120:10=12$.\\
		Vì mỗi hàng ở phía trên chứa ít hơn một viên so với hàng ở ngay phía dưới nó nên ta thu được cấp số cộng có công  sai $d=-1$.\\
		 Như vậy $u_n=12=u_1+(n-1)(-1) \Rightarrow n=13$.\\
		Vậy số viên gạch hình  vuông cần thiết để ốp hết bức tường đó là\\
		$S_{13}=\dfrac{(u_1+u_{13})13}{2}=234$ (viên gạch). 
	} 
\end{ex}

% nguồn: Ok_TraLoiNgan_Lop11_Lan2
\begin{ex}%[1D2H2-7]%[ma:fa07eea7]
	Một chiếc đồng hồ đánh chuông, số tiếng chuông được đánh bằng số giờ mà đồng hồ chỉ tại thời điểm đánh chuông. Hỏi một ngày  đồng hồ đó đánh bao nhiêu tiếng chuông báo giờ
	\shortans{156} 
	 \loigiai{
		Gọi $u_n$  là số tiếng chuông ở giờ thứ $n$, trong đó $1\le n\le 12$. Sau $12$ giờ đồng hồ đánh được
		$S_{12}=u_1+\cdots+u_{12}=1+2+\cdots+12=6\cdot (1+12)=78$  tiếng.\\
		Một ngày đồng hồ đánh được $78{,}2$ = $156$ tiếng. 
	 }
\end{ex}

% nguồn: Ok_TraLoiNgan_Lop11_Lan2
\begin{ex}%[1D2H2-7]%[ma:7520fa59]
	Từ $0$ giờ đến $12$ giờ trưa, chuông của  một chiếc đồng hồ quả lắc sẽ đánh bao nhiêu tiếng, biết rằng nó chỉ đánh chuông báo giờ và số tiếng chuông bằng số giờ?
	 \shortans{78}
	 \loigiai{
		Vì đồng hồ đánh chuông báo giờ với số tiếng chuông bằng số giờ nên ta có\\ 
		Lúc $1$ giờ đánh $1$ tiếng  chuông.\\
		Lúc $2$ giờ đánh $2$  tiếng chuông.\\
		 $\cdots$\\
		Lúc  $12$ giờ đánh $12$ tiếng chuông.\\
		Do đó, từ $0$ giờ đến $12$ giờ trưa số tiếng chuông đồng hồ đánh $1+2+\cdots+12$.\\ 
		Đây là tổng $12$ số hạng của cấp số cộng có số hạng đầu $u_1=1$,  công sai $d=1$, nên ta có
		\[{S_{12}}=\dfrac{12}{2}\cdot \left[2\cdot u_1+(12-1)\cdot d\right]=6\cdot (2\cdot 1+11\cdot  1)=78.\]
		Vậy từ $0$ giờ đến $12$ giờ trưa số tiếng chuông đồng hồ  quả lắc đó đánh là $78$ tiếng.
	 }
\end{ex}

% nguồn: Ok_TraLoiNgan_Lop11_Lan2
\begin{ex}%[1D2H2-7]%[ma:7ae02caf]
	Người ta trồng $465$ cây trong một khu vườn hình tam giác như sau: Hàng thứ  nhất có $1$ cây, hàng thứ hai có $2$ cây, hàng thứ ba có $3$ cây ... Số hàng cây trong khu vườn là bao nhiêu?
	\shortans{30} 
	 \loigiai{
		Cách trồng $465$  cây trong một khu vườn hình tam giác như trên lập thành một cấp số cộng $(u_n)$ với số $u_n$ là số cây ở hàng thứ $n$ và $u_1=1$ và công sai $d=1$.\\
		Tổng  số cây trồng được là $465$ suy ra
		\begin{eqnarray*}
			S_n=465&\Leftrightarrow&  \dfrac{n(n+1)}{2}=465\\
			&\Leftrightarrow&  n^2+n-930=0\\
			 &\Leftrightarrow& \hoac{&n=30\\&n=-31\quad \text{(loại)}.}
		\end{eqnarray*}
		Như vậy số hàng cây trong  khu vườn là $30$.
	 }
\end{ex}

% nguồn: Ok_TraLoiNgan_Lop11_Lan2
\begin{ex}%[1D2H2-7]%[ma:32269b51]
	Câu 10[Dự án 11  HVA 2024-2025]
	Một rạp hát  có $15$ hàng ghế xếp theo hình quạt. Hàng thứ nhất có $15$ ghế, hàng thứ hai có $18$ ghế, hàng thứ ba có $21$ ghế, \dots, cứ thế tiếp tục cho đến hàng cuối cùng. Rạp hát có tất cả bao nhiêu ghế?
	\shortans{540} 
	 \loigiai{
		Số ghế ở mỗi hàng lập thành cấp số cộng, có  số hạng đầu ${u_1}=15$, công sai $ d=3$.\\
		 Tổng số ghế ở rạp hát là tổng $15$ số hạng đầu của cấp số cộng là\\
	${S_{15}}=15{u_1}+\frac{15\cdot  14\cdot d}{2}=15\cdot 15+\frac{15\cdot 14\cdot 3}{2}=540$.}
\end{ex}

% nguồn: Ok_TraLoiNgan_Lop11_Lan2
\begin{ex}%[1D2H2-7]%[ma:99eb3624]
	 \immini[thm]
	{Một chiếc đồng hồ báo thức, kể từ thời điểm $0$ (giờ) thì sau mỗi giờ thì số tiếng chuông được đánh đúng bằng số giờ  mà đồng hồ chỉ tại thời điểm đánh chuông. Hỏi một ngày đồng hồ đó đánh bao nhiêu tiếng chuông?}
	{\begin{tikzpicture}[scale=1.5,font=\footnotesize,line join=round,line cap=round,>=stealth]
		\draw[thick]  (0,0) circle (1);
		\foreach \num  in {1,...,12}{
			\node at ({0.92*sin (30*\num)},{0.92*cos (30*\num)})  {\num};
		} 
		\foreach \i in  {1,...,60}{
			\draw[thin] ({0.92*sin (6*\i)},{0.92*cos (6*\i)}) --  ({1*sin (6*\i)},{1*cos (6*\i)});
		} 
		\draw[thick] (0,0) -- ({0.5*sin (30*10)},{0.5*cos (30*10)}); 
		\draw[thick,blue] (0,0)  -- ({0.8*sin (30*2)},{0.8*cos (30*2)});
		\draw[red] (0,0) -- ({0.9*sin (30*7)},{0.9*cos (30*7)}); 
		 \fill (0,0) circle (0.03);
	\end{tikzpicture}
	 }
	\shortans{300} 
	\loigiai{ 
		Kể từ lúc $1$ (giờ) đến $24$ (giờ) số tiếng chuông được đánh lập thành cấp số cộng có $24$ số  hạng với $u_1=1$, công sai $d=1$.\\
		Vậy số tiếng chuông được đánh trong $1$ ngày là  $S=S_{24}=\dfrac{24}{2}(u_1+u_{24})=12(1+24)=300$.
	 }
\end{ex}

% nguồn: Ok_TraLoiNgan_Lop11_Lan2
\begin{ex}%[1D2H2-7]%[ma:650e9b42]
	\immini[thm]{Mặt cắt của một tổ ong có hình lưới tạo bởi các ô hình lục giác đều như hình vẽ. Từ một ô đầu tiên, bước thứ nhất, các ong thợ tạo ra vòng $1$ gồm $6$ ô lục giác; bước thứ hai, các ong thợ sẽ tạo ra vòng $2$ có $12$ ô bao quanh vòng $1$; bước thứ ba, các ong thợ sẽ tạo ra $18$ ô bao quanh vòng $2$; cứ thế tiếp tục như hình vẽ bên dưới. Tính số ô trên vòng thứ $8$ mà ong thợ tạo ra? 
	\shortans{48} 
	}{ 
	\begin{tikzpicture}[scale=0.4,font=\footnotesize,line join=round,line cap=round,>=stealth]
			\def\k{0.7} 
			 \node[draw,regular polygon,regular polygon sides=6,minimum size=\k cm,double,fill=orange] (hexagon) {};
			\def\r{1.33} 
			\foreach \angle in  {-90,90,30,-30,150,-150}{
				\node[draw,regular polygon,regular polygon sides=6,minimum size=\k cm,fill=orange!20] at ({1.2*\r*cos  (\angle)},{1.2*\r*sin (\angle)}) {};
			} 
			\foreach \angle in  {-90,90,30,-30,150,-150}{
				\node[draw,regular polygon,regular polygon sides=6,minimum  size=\k cm,fill= orange!50] at ({2.4*\r*cos (\angle)},{2.4*\r*sin (\angle)}) {};
			 }
			\foreach \angle  in {0,180,60,-60,120,-120}{
				 \node[draw,regular polygon,regular polygon sides=6,minimum size=\k cm,fill= orange!50] at ({2.1*\r*cos (\angle)},{2.1*\r*sin (\angle)}) {};
			 }
		\end{tikzpicture}
	} 
	\loigiai{ 
		Số  ô lục giác ở từng vòng tạo thành cấp số cộng có công sai $d=6$.\\
		Với cấp số  cộng có $u_1=6$; $u_2=12$; $u_3=18$.\\
		Vậy số ô trên vòng thứ $8$ mà ong thợ tạo ra  là $u_8=u_1+(n-1)d=6+(8-1)\cdot 6=48$.
	} 
\end{ex}

% nguồn: Ok_TraLoiNgan_Lop11_Lan2
\begin{ex}%[1D2H2-7]%[ma:cef62378]
	Câu 15[Dự án 11 HVA 2024-2025] 
	Trong sân vận động,  dãy đầu tiên có $15$ ghế, các dãy liền sau nhiều hơn dãy trước $4$ ghế, hỏi dãy ghế thứ $10$ bao nhiêu ghế?
	\shortans{51} 
	\loigiai{ 
		Trả lời:  51\\
		Gọi ${u_1}, {u_2}, \dots, {u_n}$ lần  lượt là số ghế của dãy ghế thứ nhất, dãy ghế thứ hai,… và dãy ghế thứ $n$.\\
		Do dãy đầu tiên có $15$ ghế, các dãy liền sau nhiều  hơn dãy trước $4$ ghế nên ta được dãy số $( {u_n})$ là cấp số cộng với ${u_1}=15$ và công sai $ d=4 $.\\
		Khi đó ${u_n}=15+(n-1)\cdot  4$.\\
		Số ghế của  dãy ghế thứ $10$ là ${u_{10}}=15+(10-1)\cdot 4=51$.\\
	} 
\end{ex}

% nguồn: Ok_TraLoiNgan_Lop11_Lan2
\begin{ex}%[1D2H2-7]%[ma:e026dac6]
	An định xếp  một hình tháp bởi các mảnh ghép tam giác.
	\begin{center}
		\begin{tikzpicture}[yscale=sqrt (.75),xslant=.5]
				 \def\n{4}
				\clip[preaction=draw] (0,0)--(\n,0)--(0,\n)--cycle; 
				\draw 
					(0,0) grid  (\n,\n)
					[xslant=-1]  (0,0) grid (\n,\n);
			\end{tikzpicture}
	\end{center}
	Tầng dưới cùng An xếp $35$ hình và tầng tiếp theo ít hơn tầng dưới nó hai hình. An  xếp cho đến khi không xếp lên được nữa. Hỏi An cần bao nhiêu mảnh ghép hình tam giác để xếp xong tháp?
	 \shortans{324}
	 \loigiai{
		Theo giả thiết tòa tháp được xếp bằng các hình tam giác, số các  tam giác xếp theo quy luật là một cấp số cộng.\\
		Gọi  số tam giác tầng trên cùng là $u_1$ thì ta có $u_1=1$, $d=2$.\\
		Gọi tầng dưới cùng là $u_n\Rightarrow u_n=35$.\\ 
		Ta có 
		\begin{eqnarray*}
			&&u_n=u_1+(n-1)d=35\\ 
			&\Leftrightarrow& 1+(n-1)\cdot 2=35\\ 
			 &\Leftrightarrow& n=18\\
			&\Rightarrow&  S_{18}=\dfrac{18(1+35)}{2}=324.
		\end{eqnarray*}
		Vậy tổng số tam  giác trong $18$ tòa tháp trên là $S=324$.\\
		Chú ý giáo viên có thể dùng tổng của dãy số lẻ $1+3+5+\cdots+35$ trong tiết quy nạp lớp  10.
	 }
\end{ex}

% nguồn: Ok_TraLoiNgan_Lop11_Lan2
\begin{ex}%[1D2H2-7]%[ma:8c4c2a8b]
	Một rạp hát có $18$ hàng ghế xếp theo hình quạt. Hàng thứ nhất có $16$ ghế, hàng thứ hai có  $20$ ghế, hàng thứ ba có $24$ ghế, ... cứ thế cho đến hàng cuối cùng. Hỏi tổng số ghế có trong rạp là bao nhiêu?
	 \shortans{900}
	 \loigiai{
		Gọi số ghế ở hàng thứ $n$ là $u_n$.\\ 
		 Ta có $u_1=16$, $u_2=20$, $u_3=24$, $\cdots$\\
		Nhận xét dãy  số $(u_n)$ là cấp số cộng có $u_1=16$, công sai $d=u_2-u_1=20-16=4$.\\
		Khi đó $S_{18}=\dfrac{n\left[2u_1+(n-1)d\right]}{2}=\dfrac{18(2\cdot 16+17\cdot  4)}{2}=900$.\\
		 Vậy tổng số ghế trong rạp hát là $900$ ghế.
	} 
\end{ex}

% nguồn: Ok_TraLoiNgan_Lop11_Lan2
\begin{ex}%[1D2H2-7]%[ma:9a1a866d]
	Một rạp hát có $30$ dãy ghế, dãy đầu tiên có $25$ ghế. Mỗi dãy sau có hơn dãy trước $3$ ghế. Hỏi rạp hát có tất cả  bao nhiêu ghế?
	\shortans{2,06} 
	 \loigiai{
		Số ghế của mỗi  dãy theo thứ tự đó lập thành một cấp số cộng có $30$ số hạng có công sai $d=3$ và $u_1=25$.
		Tổng số ghế là ${S_{30}}=u_1+u_2+\cdots +{u_{30}}=30u_1+\dfrac{30\cdot  29}{2}d=2\,055$.
	 }
\end{ex}

% nguồn: Ok_TraLoiNgan_Lop11_Lan2
\begin{ex}%[1D2H2-7]%[ma:ae516cd5]
	Nam chơi trò chơi xếp các que diêm thành tháp theo qui tắc  thể hiện như hình vẽ.
	\begin{center}
		\begin{tikzpicture}[scale=0.6,font=\footnotesize,line join=round,line cap=round,>=stealth]
				\draw[fill=black]  (0,0)--(0,1) circle (2pt)--(1,1) circle (2pt)--(1,0);
				\draw[fill=black] (2,0)--(2,1)  circle (2pt)--(5,1) circle (2pt)--(5,0)
					(3,0)--(3,1)  circle (2pt)--(3,2) circle (2pt)--(4,2) circle (2pt)--(4,1) circle (2pt)--(4,0);
				 \draw[fill=black]
					(6,0)--(6,1) circle  (2pt)--(11,1) circle (2pt)--(11,0)
					(7,0)--(7,1) circle (2pt)--(7,2) circle (2pt)--(8,2) circle (2pt)--(8,3) circle (2pt)--(9,3) circle (2pt)--(9,2) circle (2pt)--(10,2)  circle (2pt)--(10,1) circle (2pt)--(10,0)
					(8,0)--(8,1) circle  (2pt)--(8,2)
					(8,2)--(9,2) 
					(9,2)--(9,1) circle (2pt)--(9,0); 
				 \draw
					(0.5,-0.5)  circle (1pt) node {1 tầng}
					(3.5,-0.5) circle (1pt) node {2 tầng} 
					(8.5,-0.5) circle (1pt) node {3  tầng};
			\end{tikzpicture}
	\end{center}
	Tính số que diêm để bạn Nam  xếp được tháp có $10$ tầng.
	\shortans{210} 
	\loigiai{ 
		Tổng số que diêm  ở $1$ tầng là $u_1=3$.\\
		Tổng số que  diêm ở $2$ tầng là $u_1+u_2=3+7$.\\
		Tổng  số que diêm ở $3$ tầng là $u_1+u_2+u_3=3+7+11$.\\
		$\cdots$\\ 
		Ta có cấp số cộng  $(u_n)$ với $u_1=3$, $d=4$\\
		Để cần có $10$ tầng cần tổng ${S_{10}}=\dfrac{n}{2}\left[2u_1+(n-1)d\right]=\dfrac{10}{2}(2\cdot 3+9\cdot  4)=210$ que.
	 }
\end{ex}

% nguồn: Ok_TraLoiNgan_Lop11_Lan2
\begin{ex}%[1D2H2-7]%[ma:5a87403e]
	Litva sẽ tham gia vào  cộng đồng chung châu Âu sử dụng đồng Euro là đồng tiền chung vào ngày $01$ tháng $01$ năm $2\,015$ Để kỷ niệm thời khắc lịch sử này, chính quyền đất nước này quyết định dùng $122\,550$ đồng tiền xu Litas Lithuania cũ của đất nước để xếp một mô hình kim tự tháp. Biết rằng tầng dưới cùng có $4\,901$ đồng và cứ lên thêm một tầng thì số đồng xu giảm đi $100$ đồng. Hỏi mô hình Kim tự tháp này có tất cả bao nhiêu tầng?
	\shortans{50} 
	 \loigiai{
		Gọi $u_n(n\in \mathbb{N^*})$ số đồng xu tầng thứ $n$. \\Theo đề ta có $(u_n)$ là cấp số cộng có $u_1=4\,901$, công sai $d=-100$ và $S_n=122\,550$.\\ 
		Ta có 
		\begin{eqnarray*}
			&&S_n=nu_1+\dfrac{n(n-1)}{2}d\\ 
			&\Leftrightarrow&  122\,550=4\,901n+\dfrac{n(n-1)}{2}\cdot (-100)\\
			 &\Leftrightarrow& 122\,550=4\,901n-50n^2+50n\\
			&\Leftrightarrow& -50n^2+4\,951n-122\,550=0\\ 
			&\Leftrightarrow&  \hoac{&n=50\\&n=\dfrac{2\,451}{50}\quad \text{(loại)}.}
		\end{eqnarray*}
		Vậy mô hình kim  tự tháp có tổng cộng $50$ tầng.
	 }
\end{ex}

% nguồn: Ok_TraLoiNgan_Lop11_Lan2
\begin{ex}%[1D2H2-7]%[ma:a5fc091b]
	Một công ty thực hiện việc trả lương cho các công nhân theo phương thức sau Mức lương của quý làm việc đầu tiên cho công ty là $13{,}5$ triệu đồng/ quý, và kể từ quý làm việc thứ hai, mức lương sẽ  được tăng thêm $500\, 000$ đồng mỗi quý. Tính tổng số tiền lương một công nhân nhận được sau ba năm làm việc cho công ty. (đơn vị triệu đồng)
	 \shortans{195}
	\loigiai{ 
		Ta thấy tổng số tiền lương trong ba năm là tổng $12$ số hạng đầu tiên của cấp số cộng $(u_n)$ với $u_1=13{,}5$; $d=0{,}5$ nên tổng số tiền lương một công  nhân nhận được sau ba năm làm việc cho công ty là ${S_{12}}=\dfrac{12\left(2\cdot 13{,}5+(12-1)0{,}5\right)}{2}=195$ triệu đồng.
	} 
\end{ex}

% nguồn: Ok_TraLoiNgan_Lop11_Lan2
\begin{ex}%[1D2H2-7]%[ma:c80d11d4]
	Giá của một chiếc xe ô tô lúc mới mua là $680$ triệu đồng. Cứ sau mối năm sử dụng, giá của chiếc xe ô tô giảm $55$ triệu đồng. Sau $5$ năm giá còn lại của chiếc xe là bao nhiêu triệu đồng? (Làm  tròn đến hàng đơn vị)
	 \shortans{460}
	\loigiai{ 
		 Giá của chiếc xe sau $n$ năm là $u_n=680-55\cdot(n-1)$.\\
		Vậy sau $5$ năm sử dụng giá của chiếc xe  là ${u_5}=680-55\cdot(5-1)=460$ (triệu đồng).
	} 
\end{ex}

% nguồn: Ok_TraLoiNgan_Lop11_Lan2
\begin{ex}%[1D2H2-7]%[ma:8ed94be4]
	 Một ngôi nhà hình kim tự tháp (có gạch nâu ốp bên ngoài) được bao quanh bởi rất nhiều cây cối và là nơi tuyệt vời để nghỉ mát mùa hè. Ngôi nhà có chiều dài và chiều rộng đều là $6{,}8$ m, chiều cao là $2{,}72$ m. Khi xây dựng ngôi nhà, người chủ đã tính toán số viên gạch nâu hình hộp chữ nhật cần ốp tường, biết hàng trên ít hơn hàng dưới $1$ viên, hàng trên cùng là $1$ viên, kích thước viên gạch nâu hình hộp chữ nhật là $0{,}2$m -- $0{,}08$ m -- $1$ m. Hãy dự tính số viên gạch nâu ốp tường cả bốn mặt của ngôi nhà.
	 \shortans{2,38}
	\loigiai{ 
		Một bức tường có $2{,}72\div 0{,}08 = 34$ hàng gạch.\\ 
		Số gạch ở  mỗi hàng tạo thành một cấp số cộng với số hạng đầu $u_1=34$ và công sai $d=1$.\\
		Số viên gạch trên một bức tường là $S_{34}=34\cdot 1+\dfrac{34\cdot 33}{2}\cdot 1=595$ viên gạch.\\ 
		 Vì $4$ mặt đều bằng nhau nên có $4\cdot 595=2\,380$ viên gạch người chủ dự tính đặt mua.
	 }
\end{ex}

% nguồn: Ok_TraLoiNgan_Lop11_Lan2
\begin{ex}%[1D2H2-7]%[ma:407c8321]
	\immini[thm] 
	{Chuông đồng hồ ở một toà tháp đánh số tiếng đúng bằng số giờ và cứ mỗi $30$ phút không phải là giờ đúng thì đánh $1$ tiếng chuông. Hỏi bắt đầu từ lúc $1$ giờ đêm đến $12$ giờ trưa, chuông đồng hồ đó đã đánh tất cả bao nhiêu  tiếng?}
	{\begin{tikzpicture}[line cap=round,line join=round,scale=0.3]
		 \path[fill=black] (8.72,1.8).. controls +(52:0.1) and +(-90:1.3)..
			(8.79,3.2) -- (8.79,4.52) -- (8.59,4.73) -- (8.38,4.94)  -- (9,4.94).. controls +(0:0.6) and +(-104:0.1)..
			(9.63,5.02).. controls +(76:0)  and +(-91:0.9)..
			(9.68,6.76).. controls +(88:0.9) and +(-93:0.4).. 
			(9.75,9.07) -- (9.77,9.72) -- (11.67,9.72) -- (13.58,9.72) -- (13.59,9.57).. controls  +(-83:0.1) and +(91:1)..
			(13.65,7.67).. controls +(-89:1) and +(92:0.3).. 
			(13.71,5.43) --  (13.73,4.94) -- (14.36,4.94) -- (14.99,4.94) -- (14.78,4.73) -- (14.58,4.52) -- (14.58,3.18) -- (14.58,1.83) -- (14.66,1.78).. controls +(-36:0.1) and +(0:1.2)..
			(13.47,1.72) -- (12.2,1.72) -- (12.18,2.55).. controls  +(92:0.5) and +(-74:0.1)..
			(12.11,3.52).. controls +(105:0.2)  and +(0:0.1)..
			 (11.68,4.01).. controls +(180:0.1) and +(75:0.2)..
			 (11.26,3.52).. controls +(-104:0.1) and +(90:0.4)..
			(11.15,1.85) -- (11.15,1.72) -- (9.89,1.72) -- (8.64,1.72) -- cycle 
			(11.97,5.7)..  controls +(90:0.8) and +(-51:0.1)..
			 (11.89,6.57).. controls +(133:0.2) and +(47:0.2)..
			 (11.48,6.57).. controls +(-129:0.1) and +(90:0.8)..
			(11.4,5.7) -- (11.4,4.94) -- (11.69,4.94)  -- (11.97,4.94) -- cycle
			(11.04,8.19).. controls +(90:0.8) and +(-49:0.1).. 
			(10.97,9.07).. controls +(135:0.1)  and +(27:0.1)..
			(10.61,9.12).. controls +(-149:0.1)  and +(91:0.8)..
			(10.5,8.12) -- (10.51,7.4) -- (10.78,7.4) -- (11.04,7.4) -- cycle 
			(12.87,8.12).. controls  +(89:0.8) and +(-29:0.1)..
			(12.76,9.12).. controls  +(156:0.1) and +(45:0.1)..
			(12.41,9.07).. controls  +(-139:0.1) and +(90:0.8)..
			(12.33,8.2) -- (12.33,7.4) --  (12.6,7.4) -- (12.86,7.4) -- cycle
			(9.42,10.05) -- (9.26,10.24) -- (9.51,10.25) -- (9.76,10.26) -- (9.76,11.71) -- (9.76,13.15) -- (11.67,13.15) -- (13.58,13.15) -- (13.59,11.71) -- (13.6,10.28) -- (13.84,10.27) -- (14.08,10.26) -- (13.91,10.06)  -- (13.75,9.87) -- (11.67,9.87) -- (9.58,9.87) -- cycle
			(12.27,10.22).. controls +(26:0.5)  and +(-92:0.5)..
			(13.07,11.42).. controls +(87:1.1) and +(25:1.1).. 
			(11.04,12.78).. controls +(-155:0.3) and +(64:0.3).. 
			 (10.39,12.14).. controls +(-115:0.9) and +(162:0.9)..
			(11.28,10.13).. controls  +(-18:0.1) and +(180:0.2)..
			 (11.73,10.09).. controls +(4:0.3) and +(-155:0.2)..
			(12.27,10.22)  -- cycle
			(11.33,10.21)..  controls +(165:0.5) and +(-71:0.4)..
			(10.38,11.08)..  controls +(108:0.3) and +(-117:0.3)..
			(10.45,12.09).. controls +(64:0.6)  and +(169:0.6)..
			(11.87,12.83).. controls +(-11:0.5)  and +(115:0.4)..
			(12.87,12.09).. controls +(-63:0.2)  and +(90:0.3)..
			(12.99,11.51).. controls +(-90:0.3) and +(65:0.2).. 
			(12.9,10.97)..  controls +(-115:0.4) and +(14:0.4)..
			 (11.99,10.21).. controls +(-165:0.3) and +(-15:0.3)..
			 (11.33,10.21) -- cycle
			(12.14,10.34).. controls +(19:0.3) and +(-116:0.3).. 
			 (12.8,10.97).. controls +(62:0.2) and +(-90:0.3)..
			(12.92,11.51)..  controls +(90:0.2) and +(-65:0.2)..
			(12.82,12.02)..  controls +(115:0.6) and +(10:0.6)..
			 (11.46,12.76).. controls +(-169:0.4) and +(60:0.4)..
			(10.52,12.08).. controls +(-117:0.2) and  +(90:0.3)..
			(10.42,11.53).. controls +(-90:0.3) and +(114:0.2).. 
			(10.51,10.97).. controls +(-64:0.7) and +(-159:0.7).. 
			(12.14,10.34) --  cycle
			 (11.62,10.49).. controls +(90:0.3) and +(94:0.3)..
			(11.72,10.48)..  controls +(-90:0.2) and +(0:0.1)..
			 (11.67,10.29).. controls +(180:0) and +(-90:0.2)..
			(11.62,10.49)  -- cycle
			(11.03,10.46).. controls  +(117:0) and +(180:0.1)..
			(11.23,10.79).. controls +(0:0) and  +(62:0.2)..
			(11.18,10.58)..  controls +(-120:0.2) and +(-59:0.1)..
			(11.03,10.46) -- cycle 
			(12.14,10.58).. controls  +(120:0.2) and +(170:0.1)..
			(12.11,10.79).. controls +(0:0) and  +(122:0.1)..
			 (12.22,10.65).. controls +(-60:0.2) and +(-59:0.2)..
			(12.14,10.58) -- cycle 
			(10.6,10.9)..  controls +(117:0) and +(-151:0.1)..
			(10.74,11.03).. controls  +(29:0.2) and +(58:0.1)..
			(10.92,11.05)..  controls +(-120:0.1) and +(-63:0)..
			(10.6,10.9)  -- cycle
			(12.52,10.96).. controls +(149:0.1)  and +(-146:0)..
			(12.4,11.09).. controls +(45:0) and +(90:0.1).. 
			(12.72,10.93).. controls  +(-90:0.1) and +(-30:0.2)..
			(12.52,10.96) -- cycle 
			 (11.57,11.37).. controls +(146:0) and +(-90:0.1)..
			(11.51,11.5).. controls +(90:0.1) and +(-16:0.2).. 
			(11.2,11.77).. controls +(157:0.1)  and +(-135:0)..
			(11.02,12).. controls +(45:0)  and +(135:0.2)..
			(11.36,11.8).. controls +(-44:0.3)  and +(-168:0.5)..
			(12.23,11.69).. controls +(12:0.2) and  +(180:0)..
			(12.57,11.75).. controls +(-45:0) and  +(17:0.2)..
			(12.21,11.62).. controls +(-162:0.2) and +(90:0).. 
			(11.83,11.47).. controls +(-90:0) and +(0:0).. 
			(11.68,11.33)..  controls +(180:0) and +(-27:0)..
			(11.57,11.37) --  cycle
			(10.44,11.51)..  controls +(90:0) and +(180:0.1)..
			(10.63,11.55).. controls  +(0:0.1) and +(90:0)..
			(10.83,11.51).. controls +(-90:0)  and +(0:0.1)..
			(10.63,11.47).. controls  +(180:0.1) and +(-90:0)..
			 (10.44,11.51) -- cycle
			(12.51,11.51).. controls  +(90:0) and +(180:0.1)..
			(12.69,11.55).. controls +(0:0.1)  and +(90:0)..
			(12.87,11.51).. controls +(-90:0) and  +(0:0.1)..
			(12.69,11.47).. controls +(180:0.1) and +(-90:0).. 
			(12.51,11.51) -- cycle 
			(10.75,11.97).. controls +(152:0.2)  and +(-149:0.1)..
			(10.61,12.13).. controls +(37:0.1) and  +(90:0.1)..
			(10.94,11.94).. controls +(-90:0) and +(-25:0.2).. 
			(10.75,11.97) -- cycle 
			(12.39,11.93).. controls +(135:0) and +(-148:0.1).. 
			(12.52,12.06).. controls  +(31:0.2) and +(90:0.1)..
			(12.72,12.11).. controls +(-90:0) and  +(32:0.1)..
			(12.58,11.98).. controls  +(-151:0.2) and +(-63:0)..
			(12.39,11.93)  -- cycle
			(11.1,12.38).. controls  +(120:0.2) and +(180:0.1)..
			(11.07,12.58).. controls +(0:0) and +(90:0.1).. 
			(11.26,12.26).. controls  +(-90:0.1) and +(-59:0.2)..
			(11.1,12.38) --  cycle
			 (12.04,12.24).. controls +(90:0.1) and +(124:0)..
			(12.26,12.55).. controls +(-45:0.1) and +(0:0.1).. 
			 (12.09,12.22).. controls +(180:0) and +(-90:0)..
			 (12.04,12.24) -- cycle
			(11.62,12.52).. controls +(90:0.2) and +(180:0).. 
			(11.67,12.72)..  controls +(0:0.1) and +(90:0.2)..
			(11.72,12.54).. controls +(-94:0.3) and +(-90:0.3).. 
			(11.62,12.52) -- cycle 
			(9.42,13.42) -- (9.29,13.58) -- (11.69,13.58) -- (14.1,13.58) --  (13.94,13.42) -- (13.78,13.26) -- (11.67,13.26) -- (9.56,13.26) -- cycle
			 (9.75,13.95).. controls +(94:0.1) and +(-76:0)..
			(9.71,14.27).. controls +(117:0.1)  and +(180:0.3)..
			(9.99,14.33) -- (10.29,14.33) --  (10.29,14.93) -- (10.29,15.53) -- (10.08,15.62).. controls +(156:0.1) and +(-90:0)..
			(9.88,15.72)..  controls +(45:0) and +(-160:0.2)..
			(10.24,15.87) -- (10.58,16.01) --  (11.67,16.01) -- (12.75,16.01) -- (13.1,15.86) -- (13.45,15.71) -- (13.25,15.62) -- (13.04,15.54) -- (13.04,14.93) -- (13.04,14.33) -- (13.34,14.33).. controls +(0:0.2) and +(135:0)..
			 (13.66,14.31).. controls +(-45:0) and +(53:0.1)..
			(13.63,14.2).. controls +(-117:0.1) and +(90:0.1).. 
			(13.58,13.94).. controls +(-84:0.1) and  +(72:0)..
			(13.58,13.73).. controls  +(-90:0) and +(0:1.5)..
			(11.67,13.69) -- (9.76,13.69) -- cycle 
			 (11.04,14.81).. controls +(90:0.4) and +(-45:0.1)..
			(10.97,15.37).. controls +(135:0.1) and +(0:0).. 
			(10.85,15.44)..  controls +(180:0) and +(45:0.1)..
			(10.72,15.37).. controls +(-131:0.1)  and +(90:0.4)..
			(10.65,14.81) -- (10.65,14.33) -- (10.85,14.33)  -- (11.04,14.33) -- cycle
			(11.97,14.78).. controls +(90:0.6)  and +(0:0.3)..
			(11.66,15.47).. controls +(180:0)  and +(45:0.1)..
			 (11.49,15.39).. controls +(-135:0.1) and +(90:0.5)..
			(11.4,14.82) --  (11.4,14.33) -- (11.69,14.33) -- (11.97,14.33) -- cycle
			(12.68,14.86) -- (12.67,15.38) -- (12.57,15.42).. controls +(164:0.1)  and +(40:0.1)..
			(12.4,15.39).. controls +(-135:0.1) and +(90:0.4).. 
			 (12.33,14.83) -- (12.33,14.33) -- (12.51,14.33) -- (12.69,14.33) -- cycle
			(9.69,15.86).. controls +(141:0.1) and +(-158:0.5).. 
			(10.21,16.15).. controls +(24:0.1) and  +(-111:0.2)..
			(10.4,16.4).. controls  +(71:0.4) and +(-160:0.4)..
			 (11.26,17.26).. controls +(22:0.2) and +(-90:0.1)..
			(11.47,17.47).. controls  +(90:0.1) and +(-51:0.1)..
			(11.4,17.68) -- (11.33,17.77) -- (11.43,17.88).. controls  +(45:0.1) and +(-94:0.3)..
			(11.56,18.31).. controls  +(85:0.3) and +(-37:0.1)..
			(11.54,18.66).. controls +(153:0) and  +(180:0.1)..
			(11.55,18.69)..  controls +(0:0) and +(-90:0)..
			(11.62,18.76).. controls +(90:0) and  +(-153:0)..
			 (11.65,18.85).. controls +(27:0) and +(90:0)..
			(11.69,18.78)..  controls +(-90:0.1) and +(180:0.1)..
			(11.77,18.69)..  controls +(-10:0.1) and +(18:0.1)..
			(11.78,18.66).. controls +(-170:0.1)  and +(96:0.2)..
			 (11.74,18.43).. controls +(-85:0.1) and +(90:0.1)..
			(11.76,18.11).. controls +(-90:0.1) and +(135:0.1).. 
			 (11.87,17.89) -- (11.98,17.76) -- (11.9,17.67).. controls +(-129:0.1) and +(90:0.1)..
			(11.83,17.46).. controls +(-90:0.1)  and +(164:0.1)..
			(11.98,17.3).. controls +(-16:0.4) and  +(116:0.5)..
			(12.92,16.43) -- (13.04,16.17)  -- (13.33,16.05).. controls +(-23:0.3) and +(15:0.1)..
			(13.58,15.85)..  controls +(-166:0) and +(-23:0.3)..
			(13.04,16.03) -- (12.56,16.24) -- (11.7,16.25) -- (10.84,16.27) -- (10.36,16.06).. controls +(-157:0.6) and  +(-30:0.1)..
			(9.69,15.86) -- cycle; 
	\end{tikzpicture}
	} 
	 \shortans{89}
	\loigiai{ 
		Lúc $1$ giờ đêm, toà tháp đánh $1$ tiếng chuông; lúc $2$ giờ đêm, toà tháp đánh $2$ tiếng chuông...; lúc $12$ giờ trưa, toà tháp đánh $12$ tiếng chuông. Ngoài ra, mỗi $30$ phút không phải là giờ đúng thì đánh $1$ tiếng chuông (có $11$ lần như thế  từ $1$ giờ đến $12$ giờ).\\
		Vậy tổng số tiếng  chuông là
		\[S=(1+2+3+\ldots+12)+1\cdot  11=\dfrac{(1+12)\cdot 12}{2}+11=89~\text{(tiếng chuông).}\]
	 }
\end{ex}

% nguồn: Ok_TraLoiNgan_Lop11_Lan2
\begin{ex}%[1D2H2-7]%[ma:65d75d77]
	\immini[thm] 
	{Một người muốn  mua một thanh gỗ để cắt ra làm các thanh ngang của một cái thang. Biết rằng chiều dài các thanh ngang của cái thang đó (từ bậc dưới cùng) lần lượt là $45$ cm, $43$ cm, $41$ cm, \ldots, $31$ cm. Tính chiều dài thanh gỗ mà người đó cần mua, giả sử chiều dài các mối nối (phần gỗ bị cắt thành mùn cưa) là không đáng kể.}
	{\begin{tikzpicture}[>=stealth,line join=round,line cap=round,font=\footnotesize,scale=1]
		\def\bot  tomwidth{2.5}
		\def\to  pwidth{1.5}
		 \def\height{4}
		 \def\rungs{5}
		\pgfmathsetmacro{\dxleft}{(\bot  tomwidth - \to pwidth)/2}
		\pgfmathsetmacro{\dxright}{\bot tomwidth  - \dxleft}
		\draw[line width=3pt,brown!70!black,line  cap=round]
			 (0,0) -- (\dxleft,\height);
		\draw[line width=3pt,brown!70!black,line cap=round] 
			(\bot tomwidth,0)  -- (\dxright,\height);
		\foreach \i in  {1,...,\rungs}{
			 \pgfmathsetmacro{\y}{\i*\height/(\rungs+1)}
			 \pgfmathsetmacro{\xleft}{\dxleft*\y/\height}
			\pgfmathsetmacro{\xright}{\bot tomwidth  - \xleft}
			\draw[line width=2pt,gray!30!black,line  cap=round]
			(\xleft,\y) --  (\xright,\y);
		} 
		\draw[gray!30]  (-0.5,0) -- (\bot tomwidth+0.5,0);
	\end{tikzpicture}
	 }
	\shortans{304} 
	 \loigiai{
		Chiều dài các thanh ngang là dãy cấp số cộng có số hạng  đầu là $45$, công sai là $d=45-43=2$.\\
		 Vậy
		 $u_n=45-2(n-1)=47-2n$.\\
		Khi $u_n=31 \Leftrightarrow n=8$ nên cái  thang có $8$ bậc.\\
		Chiều dài thanh gỗ  mà người đó cần mua là $S_8=\dfrac{8\cdot(45+31)}{2}=304$.\\
		Vậy  chiều dài thanh gỗ là $304$ cm.
	 }
\end{ex}

% nguồn: Ok_TraLoiNgan_Lop11_Lan2
\begin{ex}%[1D2H2-7]%[ma:c80ddccc]
	Trong một khán phòng có tất cả $30$ dãy ghế, dãy đầu tiên có $15$ ghế, các dãy liền sau nhiều hơn dãy trước đó $4$ ghế,  hỏi khán phòng đó có tất cả bao nhiêu ghế?
	\shortans{2,19} 
	 \loigiai{
		Gọi $u_1; u_2; \ldots; u_{30}$ lần lượt là số ghế của dãy ghế thứ nhất, dãy ghế thứ hai, ..., dãy ghế  thứ ba mươi. Khi đó $(u_n)$ là một cấp số cộng có số hạng đầu $u_1=15$, công sai $d=4$ (trong đó $1 \le n \le 30$). Gọi $S_{30}$ là tổng số ghế trong khán phòng.\\
		Ta có 
		\begin{eqnarray*}
			S_{30}&=&u_1+u_2+\ldots+u_{30}\\ 
			 &=&\dfrac{30}{2}[2u_1+(30-1)d]\\
			&=&15(2\cdot 15+29\cdot  4)\\
			&=&2\,190. 
		\end{eqnarray*}
	 }
\end{ex}

% nguồn: Ok_TraLoiNgan_Lop11_Lan2
\begin{ex}%[1D2H2-7]%[ma:48dc2996]
	Câu 12[Dự án 11  HVA 2024-2025]
	Cho cấp số cộng $(u_n)$ có công sai $d=2$ và tổng của $2$5  số hạng đầu là 625. Tính ${u_5}$.
	\shortans{9} 
	 \loigiai{
		 Tổng của $25$ số hạng đầu là:\\
		${S_{25}}=25\cdot {u_1}+\frac{25\cdot 24\cdot  d}{2}=25\cdot {u_1}+600=625\Rightarrow {u_1}=1$.\\
	Vậy ${u_5}={u_1}+4d=1+4\cdot 2=9$.} 
\end{ex}

% nguồn: Ok_TraLoiNgan_Lop11_Lan2
\begin{ex}%[1D2H2-7]%[ma:0f8bc216]
	Giá của một chiếc xe ô tô lúc mới mua là $680$  triệu đồng. Cứ sau mỗi năm sử dụng, giá của chiếc xe ô tô giảm $55$ triệu đồng. Tính giá còn lại của chiếc xe sau $5$ năm sử dụng.
	 \shortans{405}
	 \loigiai{
		Giá của chiếc xe ô tô sau một năm sử dụng  là $680-55=625$.\\
		Giá của chiếc  xe ô tô sau mỗi năm sử dụng lập thành một cấp số cộng với số hạng đầu là $u_1=625$ và công sai $d=-55$.\\
		Do đó,  giá của chiếc ô tô sau $5$ năm sử dụng là $u_5=u_1+4d=405$.
	 }
\end{ex}

% nguồn: Ok_TraLoiNgan_Lop11_Lan2
\begin{ex}%[1D2H2-7]%[ma:8a007361]
	 Câu 1[Dự án 11 HVA 2024-2025]
	Người ta trồng $3003$ cây theo dạng một hình tam giác như sau: hàng thứ nhất trồng $1$ cây, hàng thứ hai trồng $2$ cây,  hàng thứ ba trồng $3$ cây\dots. Cứ tiếp tục trồng như thế cho đến khi hết số cây. Số hàng cây được trồng là
	 \shortans{77}
	\loigiai{ 
		Gọi số cây ở hàng thứ $ n$  là ${u_n}$.\\
		Ta có: ${u_1}=1$, ${u_2}=2$, ${u_3}=3$, \dots và $S={u_1}+{u_2}+{u_3}+ \dots  +{u_n}=3003$.\\
		Nhận xét dãy số $( {u_n})$ là cấp số cộng có ${u_1}=1$, công sai $  d=1$.\\
		Khi đó  $S=\frac{n[2{u_1}+( n-1)d]}{2}=3003$.\\
		Suy  ra $\frac{n[2\cdot 1+( n-1)1]}{2}=3003$\\
		$\Leftrightarrow$ $ n( n+1)=6006$ $\Leftrightarrow$ ${n^2}+n-6006=0$ $\Leftrightarrow$ $\hoac{& n=77 \\ & n=-78  }$ $\Rightarrow$ $ n=77$.\\
		 Vậy số hàng cây được trồng là $77$.
	 }
\end{ex}

% nguồn: Ok_TraLoiNgan_Lop11_Lan2
\begin{ex}%[1D2H2-7]%[ma:bc64a8a1]
	 \immini[thm]
	{Một  ruộng bậc thang có thửa thấp nhất (bậc thứ nhất) nằm ở độ cao $950$ m so với mực nước biển, độ chênh lệch giữa thửa trên và thửa dưới trung bình là $1{,}5$ m. Hỏi thửa ruộng ở bậc thứ $12$ có độ cao là bao nhiêu mét so với mực nước biển? (\textit{Kết quả làm tròn đến hàng đơn vị.})}
	{\begin{tikzpicture}[line cap=round,line join=round,scale=0.4]
		\definecolor{cd0d0d0}{RGB}{208,208,208} 
		 \definecolor{cb9b9b9}{RGB}{185,185,185}
		\definecolor{ca3a3a3}{RGB}{163,163,163} 
		 \definecolor{c8c8c8c}{RGB}{140,140,140}
		 \definecolor{c767676}{RGB}{118,118,118}
		 \definecolor{c606060}{RGB}{96,96,96}
		\definecolor{c494949}{RGB}{73,73,73} 
		\path[fill=cd0d0d0] (5.02,1.26).. controls  +(27:0) and +(135:0)..
			(5.07,1.26).. controls +(-45:0)  and +(-18:0)..
			(5.04,1.25).. controls  +(180:0) and +(-153:0)..
			(5.02,1.26)  -- cycle
			(6.08,1.29).. controls +(0:0) and  +(180:0)..
			(6.14,1.29).. controls +(-27:0) and +(0:0).. 
			(6.11,1.27).. controls +(180:0)  and +(-153:0)..
			(6.08,1.29) -- cycle 
			(0.93,2.24).. controls +(90:0) and  +(90:0)..
			(0.94,2.25).. controls +(-63:0)  and +(45:0)..
			(0.94,2.2).. controls +(-135:0) and +(-90:0).. 
			(0.93,2.24) --  cycle
			(12.7,3.08).. controls +(90:0) and  +(117:0)..
			(12.72,3.1).. controls  +(-63:0) and +(90:0)..
			(12.72,3.05).. controls +(-135:0)  and +(-90:0)..
			(12.7,3.08) --  cycle
			(2.37,5.05).. controls +(56:0) and  +(180:0)..
			(2.43,5.11).. controls +(0:0)  and +(45:0)..
			(2.39,5.05).. controls  +(-124:0) and +(0:0)..
			(2.33,5)..  controls +(0:0) and +(-135:0)..
			(2.37,5.05) -- cycle 
			(12.46,5.73).. controls +(121:0.1) and +(135:0.1).. 
			(12.47,5.74)..  controls +(-37:0) and +(0:0)..
			(12.5,5.69).. controls +(0:0)  and +(-45:0)..
			(12.46,5.73) --  cycle
			 (9.78,5.92).. controls +(27:0) and +(153:0)..
			(9.85,5.92).. controls +(-27:0)  and +(0:0)..
			(9.82,5.9).. controls +(180:0)  and +(-153:0)..
			(9.78,5.92) --  cycle
			(8.24,7.29) -- (8.24,7.77) --  (8.16,7.91).. controls +(120:0.3) and +(-4:0.3)..
			(7.55,8.3)..  controls +(176:0.1) and +(19:0.4)..
			(7.01,8.18).. controls +(-162:0.2) and  +(-90:0)..
			(6.62,8.06).. controls +(180:0)  and +(-131:0.2)..
			(6.88,8.37).. controls +(48:0.5) and +(-85:0.3).. 
			(7.24,9.11)..  controls +(95:0.2) and +(-36:0.2)..
			(6.97,9.64).. controls +(144:0.1)  and +(-20:0.2)..
			 (6.48,9.87).. controls +(159:0.2) and +(-153:0)..
			(6.13,10.02).. controls +(0:0) and +(-172:0.2).. 
			(6.52,10.08).. controls +(8:0.4) and  +(-131:0.2)..
			(7.24,10.35)..  controls +(49:0.2) and +(-90:0.2)..
			(7.43,10.86).. controls  +(90:0.2) and +(-60:0.5)..
			(7.12,11.63).. controls +(121:0.1) and +(0:0).. 
			(7.02,11.82).. controls  +(45:0) and +(148:0.2)..
			(7.32,11.65)..  controls +(-32:0.5) and +(180:0.2)..
			(8,11.37).. controls +(0:0.3) and  +(-115:0.3)..
			(8.68,11.8).. controls  +(74:0.1) and +(-101:0.2)..
			(8.8,12.3) -- (8.87,12.67) -- (9.01,12.27)..  controls +(-72:0.5) and +(151:0.3)..
			 (9.45,11.57).. controls +(-27:0.1) and +(180:0.2)..
			 (9.8,11.5).. controls +(0:0.3) and +(-141:0.5)..
			 (10.51,11.84).. controls +(40:0.2) and +(90:0)..
			(10.82,12.08).. controls +(-90:0) and +(69:0.2).. 
			(10.67,11.69).. controls  +(-111:0.3) and +(90:0.1)..
			(10.52,11.15).. controls +(-90:0.3) and +(153:0.3).. 
			(10.94,10.49).. controls  +(-27:0.2) and +(177:0.4)..
			(11.52,10.4) -- (11.95,10.38) -- (11.58,10.19).. controls +(-152:0.6) and  +(90:0.4)..
			(10.95,9.38).. controls +(-90:0.3) and +(140:0.9).. 
			 (11.73,8.39).. controls +(-63:0) and +(-12:0.2)..
			(11.4,8.44).. controls  +(168:0.5) and +(25:0.2)..
			(10.56,8.45).. controls +(-153:0.2) and +(68:0.2).. 
			(10.21,8.07).. controls +(-115:0.2)  and +(100:0.4)..
			 (10.21,7.3).. controls +(-79:0.2) and +(90:0)..
			(10.26,6.99).. controls +(180:0)  and +(-50:0.2)..
			(10,7.29).. controls +(127:0.2)  and +(-34:0.1)..
			(9.66,7.67).. controls +(144:0.3) and  +(29:0.3)..
			(8.87,7.7).. controls +(-151:0.2) and  +(60:0.5)..
			(8.42,7.13)  -- (8.23,6.81) -- cycle
			(9.54,8.67).. controls +(26:0.3) and  +(-128:0.2)..
			(10.05,9.08).. controls +(55:0.3)  and +(-86:0.3)..
			 (10.26,9.86).. controls +(93:0.2) and +(-65:0.2)..
			(10.15,10.3)..  controls +(117:0.5) and +(0:0.5)..
			(9.05,10.99).. controls +(-179:0.5) and +(63:0.5).. 
			(7.97,10.3).. controls +(-115:0.2) and +(87:0.2).. 
			(7.86,9.86)..  controls +(-93:0.4) and +(139:0.4)..
			(8.26,8.88).. controls +(-41:0.3) and  +(-178:0.4)..
			 (9.12,8.58).. controls +(3:0.2) and +(-155:0.2)..
			(9.54,8.67) --  cycle
			(1.51,9.53).. controls +(90:0)  and +(-135:0)..
			(1.55,9.58)..  controls +(31:0.1) and +(45:0.1)..
			(1.56,9.57).. controls  +(-127:0.1) and +(-90:0)..
			 (1.51,9.53) -- cycle
			(5.06,10.16)..  controls +(90:0.1) and +(90:0)..
			(5.07,10.21).. controls +(-72:0) and  +(72:0)..
			(5.07,10.11).. controls +(-90:0)  and +(-90:0.1)..
			(5.06,10.16) --  cycle
			(2.98,10.89).. controls  +(27:0) and +(135:0)..
			(3.04,10.89).. controls  +(-90:0) and +(0:0)..
			(3,10.88).. controls +(180:0) and  +(-153:0)..
			(2.98,10.89)  -- cycle
			(4.1,11.13).. controls  +(27:0) and +(135:0)..
			(4.15,11.13).. controls +(-45:0) and +(0:0).. 
			(4.11,11.11).. controls  +(162:0) and +(-162:0)..
			 (4.1,11.13) -- cycle;
		\path[fill=cb9b9b9] (2.77,1.26).. controls +(27:0)  and +(90:0)..
			(2.82,1.26).. controls  +(-45:0) and +(-18:0)..
			(2.79,1.25).. controls +(180:0)  and +(-153:0)..
			(2.77,1.26) --  cycle
			(4.94,1.26).. controls +(27:0) and +(135:0).. 
			(4.99,1.26)..  controls +(-45:0) and +(-18:0)..
			(4.96,1.25)..  controls +(180:0) and +(-153:0)..
			(4.94,1.26)  -- cycle
			(6,1.29).. controls +(27:0) and +(180:0).. 
			(6.05,1.29).. controls +(-45:0) and +(0:0).. 
			(6.02,1.27).. controls  +(180:0) and +(-153:0)..
			(6,1.29)  -- cycle
			(12.78,1.52).. controls +(90:0) and  +(117:0)..
			(12.8,1.54).. controls +(-63:0)  and +(90:0)..
			 (12.8,1.49).. controls +(-135:0) and +(-90:0)..
			(12.78,1.52) -- cycle 
			(12.76,2.03).. controls +(90:0)  and +(117:0)..
			(12.77,2.04)..  controls +(-63:0) and +(45:0)..
			(12.77,1.99).. controls +(-135:0)  and +(-108:0)..
			(12.76,2.03)  -- cycle
			(12.73,2.53).. controls  +(90:0) and +(90:0)..
			 (12.74,2.56).. controls +(-72:0) and +(63:0)..
			(12.74,2.48).. controls +(-90:0) and  +(-90:0)..
			(12.73,2.53) --  cycle
			(12.7,3.01).. controls +(90:0) and +(117:0).. 
			(12.72,3.02)..  controls +(-63:0) and +(90:0)..
			 (12.72,2.97).. controls +(-135:0) and +(-90:0)..
			 (12.7,3.01) -- cycle
			 (12.68,3.53).. controls +(90:0) and +(117:0)..
			(12.69,3.55).. controls  +(-63:0) and +(45:0)..
			(12.69,3.5)..  controls +(-135:0) and +(-117:0)..
			(12.68,3.53) -- cycle 
			(12.62,4.54)..  controls +(90:0) and +(117:0)..
			(12.64,4.56).. controls +(-63:0) and  +(0:0)..
			(12.64,4.51).. controls  +(-135:0) and +(-90:0)..
			(12.62,4.54)  -- cycle
			(12.6,5.04).. controls  +(90:0) and +(117:0)..
			 (12.61,5.06).. controls +(-63:0) and +(45:0)..
			(12.61,5.01)..  controls +(-135:0) and +(-90:0)..
			 (12.6,5.04) -- cycle
			 (3.85,5.93).. controls +(45:0) and +(180:0)..
			(3.91,5.98).. controls +(0:0) and  +(56:0)..
			(3.88,5.93).. controls +(-135:0) and +(0:0).. 
			(3.82,5.87)..  controls +(180:0) and +(-135:0)..
			 (3.85,5.93) -- cycle
			(9.44,5.89)..  controls +(27:0) and +(135:0)..
			(9.49,5.89).. controls +(-45:0) and +(-27:0).. 
			(9.46,5.88).. controls +(180:0) and +(-153:0).. 
			(9.44,5.89) --  cycle
			(9.89,5.92).. controls +(27:0) and  +(135:0)..
			(9.95,5.92).. controls +(-27:0) and  +(0:0)..
			(9.92,5.9)..  controls +(180:0) and +(-153:0)..
			 (9.89,5.92) -- cycle
			(11.93,5.92).. controls +(18:0) and  +(153:0)..
			(12.03,5.92).. controls  +(-18:0) and +(0:0.1)..
			(11.97,5.91)..  controls +(180:0.1) and +(-166:0)..
			(11.93,5.92) --  cycle
			(6.35,7.24)..  controls +(135:0.1) and +(141:0.1)..
			(6.36,7.25)..  controls +(-45:0) and +(0:0)..
			(6.4,7.21).. controls +(-90:0) and  +(-53:0.1)..
			 (6.35,7.24) -- cycle
			 (8.9,8.59).. controls +(27:0) and +(153:0)..
			(8.98,8.59).. controls +(-45:0) and +(0:0).. 
			(8.93,8.58).. controls  +(180:0) and +(-162:0)..
			(8.9,8.59) --  cycle
			(9.16,8.59).. controls +(18:0) and +(153:0).. 
			(9.24,8.59).. controls  +(-27:0) and +(0:0.1)..
			(9.19,8.58)..  controls +(180:0) and +(-153:0)..
			(9.16,8.59)  -- cycle
			 (8.75,8.62).. controls +(0:0) and +(135:0)..
			 (8.8,8.61).. controls +(0:0) and +(0:0)..
			(8.77,8.6)..  controls +(180:0) and +(-153:0)..
			(8.75,8.62)  -- cycle
			(9.33,8.62).. controls +(0:0)  and +(135:0)..
			(9.39,8.61).. controls +(0:0) and  +(0:0)..
			(9.35,8.6).. controls +(180:0)  and +(-153:0)..
			(9.33,8.62) --  cycle
			(1.39,8.76).. controls +(118:0.1)  and +(-106:0.2)..
			(1.36,9.24).. controls +(72:0.1) and +(-132:0.2).. 
			(1.55,9.55).. controls  +(43:0.2) and +(-94:0.2)..
			(1.71,9.86).. controls +(88:0.3) and +(180:0.3).. 
			(2.3,10.45)..  controls +(0:0.1) and +(-113:0.1)..
			 (2.46,10.53).. controls +(70:0.1) and +(-157:0.1)..
			(2.76,10.82).. controls  +(25:0.2) and +(153:0.1)..
			(3.26,10.83) --  (3.36,10.78) -- (3.43,10.85).. controls +(45:0.1) and +(-174:0.1)..
			(3.87,11.09).. controls +(7:0.1)  and +(117:0.2)..
			 (4.03,10.95).. controls +(-63:0.2) and +(72:0.2)..
			(4.09,10.24)..  controls +(-104:0.1) and +(0:0)..
			(4.05,10.1).. controls +(-45:0) and  +(150:0.1)..
			 (4.19,10.02).. controls +(-30:0.6) and +(17:0.6)..
			 (4.04,8.72).. controls +(-165:0.1) and +(0:1)..
			(2.67,8.68) --  (1.43,8.68) -- cycle
			(8.51,8.7).. controls  +(117:0) and +(153:0)..
			(8.53,8.71).. controls  +(-11:0.1) and +(0:0)..
			(8.55,8.68)..  controls +(180:0) and +(-45:0)..
			(8.51,8.7) --  cycle
			(8.2,8.93).. controls +(131:0.1) and +(180:0).. 
			(8.1,9.05).. controls  +(0:0) and +(135:0.1)..
			 (8.23,8.93).. controls +(-49:0.1) and +(0:0)..
			(8.33,8.81)..  controls +(180:0) and +(-45:0.1)..
			(8.2,8.93) -- cycle 
			(9.9,8.93).. controls +(51:0.1) and  +(180:0)..
			(10,9.02).. controls +(0:0) and +(45:0.1).. 
			(9.92,8.93).. controls +(-135:0.1)  and +(0:0)..
			(9.82,8.84).. controls +(180:0) and  +(-135:0.1)..
			 (9.9,8.93) -- cycle
			(5.45,8.98).. controls +(90:0) and  +(117:0)..
			(5.47,9).. controls +(-63:0) and  +(45:0)..
			(5.47,8.95)..  controls +(-135:0) and +(-90:0)..
			(5.45,8.98) -- cycle 
			 (5.45,9.09).. controls +(90:0) and +(117:0)..
			(5.47,9.11)..  controls +(-63:0) and +(0:0)..
			(5.47,9.06)..  controls +(-135:0) and +(-90:0)..
			(5.45,9.09) -- cycle 
			(7.86,9.6)..  controls +(90:0) and +(117:0)..
			(7.88,9.64).. controls  +(-90:0) and +(90:0)..
			(7.88,9.57)..  controls +(-117:0) and +(-90:0)..
			(7.86,9.6) -- cycle 
			 (7.86,9.97).. controls +(90:0) and +(117:0)..
			(7.88,10.01).. controls +(-90:0) and  +(90:0)..
			(7.88,9.94).. controls +(-117:0) and  +(-90:0)..
			(7.86,9.97) --  cycle
			(9.89,10.66).. controls  +(127:0.1) and +(143:0.1)..
			 (9.9,10.67).. controls +(-39:0.1) and +(0:0)..
			(9.94,10.61).. controls +(180:0) and  +(-56:0)..
			 (9.89,10.66) -- cycle
			(9.28,10.97).. controls +(27:0) and +(135:0).. 
			 (9.33,10.97).. controls +(-45:0) and +(-18:0)..
			(9.3,10.96).. controls +(180:0) and  +(-153:0)..
			 (9.28,10.97) -- cycle
			(9.02,11).. controls +(18:0) and  +(153:0)..
			 (9.12,11).. controls +(-18:0) and +(0:0.1)..
			(9.06,10.99)..  controls +(180:0.1) and +(-166:0)..
			(9.02,11)  -- cycle;
		\path[fill=ca3a3a3] (4.82,1.26).. controls +(18:0)  and +(153:0)..
			(4.9,1.26).. controls +(-27:0) and  +(0:0)..
			(4.86,1.25).. controls  +(180:0.1) and +(180:0)..
			(4.82,1.26) -- cycle 
			(5.85,1.29).. controls +(0:0) and +(180:0).. 
			(5.95,1.29).. controls +(-27:0) and +(0:0).. 
			(5.9,1.28).. controls +(180:0.1) and +(-153:0).. 
			 (5.85,1.29) -- cycle
			(8.01,1.29).. controls +(27:0)  and +(180:0)..
			(8.06,1.29).. controls +(-45:0) and +(0:0).. 
			(8.03,1.27).. controls +(180:0) and  +(-153:0)..
			(8.01,1.29)  -- cycle
			(12.76,1.91).. controls +(90:0.1) and  +(108:0)..
			 (12.77,1.95).. controls +(-63:0) and +(63:0)..
			(12.77,1.86).. controls  +(-108:0) and +(-90:0.1)..
			 (12.76,1.91) -- cycle
			(12.73,2.39).. controls +(90:0) and +(90:0).. 
			(12.74,2.43)..  controls +(-72:0) and +(63:0)..
			(12.74,2.35)..  controls +(-90:0) and +(-90:0)..
			(12.73,2.39) --  cycle
			(12.71,2.88).. controls  +(90:0.1) and +(108:0)..
			(12.72,2.93).. controls +(-90:0) and  +(90:0)..
			(12.72,2.84)..  controls +(-108:0) and +(-90:0)..
			 (12.71,2.88) -- cycle
			(12.68,3.41).. controls +(90:0)  and +(90:0)..
			(12.69,3.46).. controls  +(-72:0) and +(63:0)..
			(12.69,3.37).. controls  +(-90:0) and +(-90:0.1)..
			(12.68,3.41) --  cycle
			(12.65,3.9).. controls +(90:0.1)  and +(117:0)..
			 (12.67,3.94).. controls +(-90:0) and +(45:0)..
			(12.66,3.86).. controls +(-90:0) and +(-90:0).. 
			(12.65,3.9) --  cycle
			(12.63,4.42)..  controls +(90:0) and +(108:0)..
			(12.64,4.46).. controls +(-63:0) and +(72:0).. 
			(12.64,4.37).. controls +(-117:0) and  +(-90:0)..
			(12.63,4.42)  -- cycle
			(12.6,4.92).. controls  +(90:0) and +(90:0)..
			(12.61,4.97).. controls +(-72:0) and  +(72:0)..
			 (12.61,4.87).. controls +(-90:0) and +(-90:0)..
			(12.6,4.92) --  cycle
			(12.57,5.42).. controls +(90:0)  and +(108:0)..
			(12.59,5.47).. controls +(-90:0) and +(90:0).. 
			(12.59,5.38).. controls +(-108:0) and  +(-90:0)..
			(12.57,5.42)  -- cycle
			(7.67,5.63).. controls +(127:0.1)  and +(143:0.1)..
			(7.68,5.64).. controls +(-45:0.1) and  +(0:0)..
			(7.72,5.58).. controls +(180:0) and +(-56:0).. 
			 (7.67,5.63) -- cycle
			(7.34,5.99) -- (7.26,6.07) -- (7.35,5.99).. controls +(-41:0.1)  and +(0:0)..
			(7.43,5.9).. controls +(180:0)  and +(-51:0.1)..
			(7.34,5.99) -- cycle 
			(10,5.92).. controls +(27:0) and  +(180:0)..
			(10.05,5.92).. controls +(-45:0)  and +(0:0)..
			(10.01,5.9).. controls +(180:0)  and +(-153:0)..
			(10,5.92) -- cycle 
			(11.7,5.92).. controls +(14:0) and +(162:0).. 
			(11.83,5.92).. controls +(-14:0) and  +(0:0.1)..
			(11.76,5.91).. controls +(180:0.1)  and +(-166:0)..
			(11.7,5.92) --  cycle
			(2.68,8.33)..  controls +(0:0.4) and +(180:0.4)..
			(4.11,8.33).. controls  +(-2:0.4) and +(0:0.8)..
			(3.4,8.32).. controls +(180:0.8) and +(-178:0.4).. 
			(2.68,8.33) --  cycle
			(1.53,8.54) -- (1.47,8.61) -- (1.54,8.55).. controls +(-41:0.1) and  +(0:0)..
			 (1.61,8.47).. controls +(180:0) and +(-45:0.1)..
			 (1.53,8.54) -- cycle
			(5.21,8.5).. controls +(90:0) and  +(-146:0)..
			 (5.26,8.55).. controls +(37:0) and +(53:0.1)..
			(5.27,8.54).. controls +(-135:0.1)  and +(-90:0)..
			 (5.21,8.5) -- cycle
			(8.86,8.6).. controls +(170:0.4) and +(-66:0.4).. 
			(7.98,9.27).. controls +(116:0.6)  and +(-143:0.6)..
			(8.36,10.76).. controls +(35:0.3)  and +(166:0.3)..
			(9.34,10.95).. controls  +(-13:0.7) and +(76:0.6)..
			(10.23,9.49)..  controls +(-105:0.2) and +(48:0.2)..
			(9.94,8.98).. controls  +(-133:0.4) and +(-8:0.4)..
			(8.86,8.6)  -- cycle
			(2.03,8.67).. controls  +(2:0.3) and +(178:0.3)..
			(3.26,8.67)..  controls +(0:0.3) and +(0:0.7)..
			(2.65,8.66).. controls  +(180:0.7) and +(180:0.3)..
			(2.03,8.67)  -- cycle
			(3.96,8.7).. controls +(0:0)  and +(90:0)..
			 (4.01,8.69).. controls +(0:0) and +(0:0)..
			(3.98,8.68).. controls +(180:0)  and +(-153:0)..
			(3.96,8.7) -- cycle 
			 (4.31,8.88).. controls +(37:0) and +(180:0)..
			 (4.39,8.94).. controls +(0:0) and +(45:0)..
			(4.34,8.88).. controls +(-135:0.1) and +(0:0).. 
			(4.26,8.81).. controls +(0:0) and  +(-127:0.1)..
			(4.31,8.88) --  cycle
			(4.53,9.3).. controls +(90:0)  and +(90:0)..
			(4.54,9.33).. controls  +(-63:0) and +(63:0)..
			(4.54,9.25)..  controls +(-135:0) and +(-90:0)..
			 (4.53,9.3) -- cycle
			(4.53,9.5).. controls +(90:0)  and +(117:0)..
			(4.54,9.53).. controls  +(-63:0) and +(45:0)..
			(4.54,9.47)..  controls +(-117:0) and +(-90:0)..
			(4.53,9.5) --  cycle
			(5.18,9.6).. controls  +(127:0.1) and +(143:0.1)..
			(5.19,9.61).. controls +(-34:0)  and +(90:0)..
			(5.24,9.56).. controls +(-90:0)  and +(-51:0.1)..
			(5.18,9.6) -- cycle 
			(4.32,9.93) -- (4.27,9.99) --  (4.33,9.94).. controls +(-37:0) and +(0:0)..
			 (4.39,9.88).. controls +(-90:0) and +(-50:0.1)..
			(4.32,9.93) -- cycle 
			 (4.17,10.03).. controls +(149:0.1) and +(180:0)..
			(4.12,10.08).. controls +(0:0)  and +(135:0)..
			(4.21,10.03).. controls +(-41:0.1) and +(-35:0.1).. 
			(4.17,10.03)  -- cycle
			(4.05,10.12)..  controls +(90:0) and +(117:0)..
			(4.07,10.14).. controls +(-90:0) and  +(45:0)..
			(4.06,10.09).. controls  +(-90:0) and +(-90:0)..
			 (4.05,10.12) -- cycle
			(4.14,10.52).. controls  +(90:0.1) and +(101:0.1)..
			 (4.15,10.59).. controls +(-90:0.1) and +(90:0)..
			(4.15,10.43).. controls +(-104:0) and  +(-90:0.1)..
			(4.14,10.52) -- cycle 
			 (4.1,10.76).. controls +(90:0) and +(117:0)..
			(4.12,10.77)..  controls +(-63:0) and +(45:0)..
			(4.12,10.72).. controls +(-135:0)  and +(-90:0)..
			(4.1,10.76) -- cycle 
			(4.76,10.84) -- (4.7,10.91) -- (4.77,10.85).. controls +(-41:0.1) and  +(0:0)..
			(4.83,10.77).. controls +(0:0) and  +(-45:0.1)..
			 (4.76,10.84) -- cycle
			(3.97,11.06).. controls +(124:0)  and +(0:0)..
			 (3.95,11.11).. controls +(0:0) and +(124:0)..
			(4,11.06).. controls  +(-72:0) and +(0:0)..
			 (4.01,11.01).. controls +(180:0) and +(-72:0)..
			(3.97,11.06) --  cycle;
		\path[fill=c8c8c8c] (11.35,0.81).. controls  +(45:0) and +(135:0)..
			(11.4,0.81).. controls +(-45:0)  and +(-18:0)..
			(11.36,0.8)..  controls +(180:0) and +(-162:0)..
			(11.35,0.81) --  cycle
			(2.91,1.26)..  controls +(45:0) and +(135:0)..
			(2.96,1.26).. controls +(-90:0)  and +(-18:0)..
			(2.92,1.25).. controls +(180:0)  and +(-162:0)..
			(2.91,1.26) -- cycle 
			(4.66,1.26).. controls +(18:0) and +(162:0).. 
			(4.76,1.26).. controls +(0:0) and +(0:0).. 
			(4.71,1.25)..  controls +(180:0) and +(180:0)..
			 (4.66,1.26) -- cycle
			(5.72,1.29)..  controls +(18:0) and +(180:0)..
			(5.8,1.29).. controls +(-27:0)  and +(-11:0.1)..
			(5.75,1.28).. controls +(180:0) and +(-153:0).. 
			(5.72,1.29)  -- cycle
			(11.71,1.37).. controls +(180:0.1) and +(-11:0.4).. 
			 (10.95,1.5).. controls +(170:2.2) and +(0:2)..
			(7.14,1.83).. controls +(-180:3.9) and  +(-16:0.4)..
			(2.5,1.93).. controls +(162:0.2) and +(-90:0.1).. 
			(2.24,2.21)  -- (2.25,2.33) -- (2.17,2.28).. controls +(-148:0.4) and +(90:0.2)..
			(1.69,1.66) -- (1.69,1.54) --  (1.61,1.55).. controls +(169:0.4) and +(-90:0.3)..
			 (0.95,2.27).. controls +(90:0.1) and +(-119:0.4)..
			 (1.19,2.84).. controls +(58:0.1) and +(-114:0.3)..
			(1.51,3.53)..  controls +(67:0.9) and +(-139:0.4)..
			(2.47,5.12).. controls +(40:0.1)  and +(-153:0.3)..
			(3.11,5.48).. controls  +(26:0.3) and +(-144:0.1)..
			 (3.76,5.83).. controls +(36:0.1) and +(-128:0.4)..
			(4.3,6.45).. controls +(51:0.8) and  +(-154:0.3)..
			(5.22,7.38).. controls +(26:0.4) and +(137:0.4).. 
			 (6.35,7.24).. controls +(-45:0.1) and +(124:0.2)..
			(6.66,6.82).. controls +(-56:0.3)  and +(133:0.5)..
			(7.32,6.01)..  controls +(-47:0.4) and +(135:0)..
			(7.78,5.53).. controls +(0:0)  and +(-162:0.1)..
			(7.97,5.59)..  controls +(22:0.7) and +(179:1.6)..
			(11.26,5.92).. controls +(-1:0.8)  and +(157:0.1)..
			(12.27,5.85).. controls +(-24:0.3)  and +(93:0.6)..
			 (12.59,5.06).. controls +(-86:0.3) and +(90:0.1)..
			(12.63,4.39) --  (12.65,4.21) -- (12.51,4.18).. controls +(-167:0.2) and +(-20:0.7)..
			(11.5,4.39).. controls  +(161:0.4) and +(-10:0.1)..
			(10.69,4.65).. controls +(169:0.2)  and +(0:1.1)..
			(9.17,4.69) -- (7.84,4.7) -- (7.66,4.76).. controls +(160:0.2) and  +(-90:0.2)..
			(7.37,5.1).. controls  +(90:0.1) and +(0:0)..
			(7.35,5.21).. controls +(180:0)  and +(69:0.1)..
			(7.16,4.89).. controls +(-113:0.1) and +(90:0.1).. 
			(7.11,4.6).. controls +(-92:0.3) and +(170:0.4).. 
			(7.59,4.15).. controls +(-5:0.1)  and +(180:0.8)..
			 (8.96,4.13).. controls +(0:0.8) and +(170:0.1)..
			(10.36,4.1).. controls +(-8:0.2) and +(160:0.7).. 
			(11.39,3.78).. controls  +(-19:0.5) and +(-177:0.2)..
			(12.1,3.6).. controls +(5:0.2) and +(-128:0.2).. 
			(12.58,3.82).. controls +(45:0.1) and +(90:0).. 
			(12.65,3.89)..  controls +(-90:0) and +(92:0.3)..
			(12.68,3.33) -- (12.71,2.78) -- (12.61,2.74).. controls  +(-170:0.1) and +(0:0.2)..
			(12.21,2.71).. controls +(178:0.3) and +(-14:0.8).. 
			(11.1,2.92).. controls +(166:0.7) and  +(-5:0.3)..
			 (9.96,3.15).. controls +(177:0.2) and +(-1:1)..
			(7.83,3.19).. controls +(179:1)  and +(-4:0.1)..
			(5.78,3.23).. controls  +(174:0.5) and +(-90:0.2)..
			(5.08,3.6).. controls +(90:0.1)  and +(0:0)..
			(5.06,3.7).. controls +(180:0)  and +(61:0.1)..
			(4.78,3.39).. controls +(-121:0.1) and +(83:0.1).. 
			(4.71,3.16)..  controls +(-102:0.3) and +(169:0.5)..
			(5.24,2.68).. controls +(-10:0.2) and  +(-179:1)..
			(6.6,2.65)..  controls +(1:1.8) and +(172:0.3)..
			(9.55,2.62).. controls +(-4:0.1) and  +(166:0.5)..
			(10.68,2.38).. controls +(-14:1)  and +(-170:0.3)..
			(12.11,2.17)..  controls +(8:0.2) and +(-128:0.1)..
			(12.67,2.42)  -- (12.72,2.49) -- (12.74,2.35).. controls +(-90:0.1) and +(93:0.2)..
			(12.77,1.82) -- (12.79,1.43) --  (12.72,1.42).. controls +(-167:0.3) and +(-9:0.3)..
			(11.71,1.37)  -- cycle
			 (12.65,4.02).. controls +(90:0) and +(90:0)..
			(12.66,4.05).. controls +(-45:0)  and +(63:0)..
			 (12.66,3.99).. controls +(-90:0) and +(-90:0)..
			(12.65,4.02) --  cycle
			(1.76,8.4).. controls +(157:0.1)  and +(-45:0.1)..
			(1.54,8.55) --  (1.45,8.65) -- (2.71,8.66) -- (3.98,8.67) -- (4.13,8.74).. controls +(24:0.2) and +(-114:0.2)..
			(4.49,9.12).. controls  +(66:0.3) and +(-38:0.3)..
			(4.26,9.99).. controls +(140:0.1) and  +(0:0)..
			(4.11,10.08)..  controls +(180:0.1) and +(-106:0.2)..
			(4.12,10.29).. controls  +(73:0.2) and +(-66:0.2)..
			(4.05,10.95) -- (3.97,11.12) -- (4.15,11.11).. controls  +(-4:0.7) and +(67:0.7)..
			(4.96,9.75).. controls  +(-117:0) and +(153:0)..
			 (5.03,9.69).. controls +(-20:0.1) and +(115:0.1)..
			 (5.39,9.34).. controls +(-63:0.2) and +(68:0.2)..
			(5.41,8.79).. controls +(-111:0.2) and  +(24:0.2)..
			(5.04,8.4) -- (4.91,8.33) -- (3.4,8.33) -- (1.89,8.33) -- cycle; 
		\path[fill=c767676]  (11.48,0.81).. controls +(14:0) and +(166:0)..
			 (11.64,0.81).. controls +(0:0) and +(0:0.1)..
			(11.55,0.8).. controls  +(180:0.1) and +(180:0.1)..
			 (11.48,0.81) -- cycle
			(3,1.26).. controls  +(27:0) and +(153:0)..
			 (3.08,1.26).. controls +(-45:0) and +(0:0)..
			 (3.03,1.25).. controls +(180:0) and +(180:0)..
			(3,1.26) --  cycle
			(4.43,1.26).. controls  +(14:0) and +(166:0)..
			 (4.57,1.26).. controls +(0:0) and +(0:0.1)..
			(4.5,1.25).. controls +(180:0.1)  and +(180:0)..
			(4.43,1.26) --  cycle
			(8.51,1.26).. controls +(27:0) and +(135:0).. 
			(8.57,1.26).. controls  +(-90:0) and +(-18:0)..
			(8.53,1.25).. controls +(180:0) and +(-153:0).. 
			 (8.51,1.26) -- cycle
			 (5.51,1.29).. controls +(14:0) and +(180:0)..
			(5.64,1.29).. controls +(-14:0) and +(0:0.1).. 
			(5.57,1.28).. controls +(180:0.1) and +(-166:0).. 
			(5.51,1.29)  -- cycle
			 (8.14,1.29).. controls +(0:0) and +(180:0)..
			 (8.21,1.29).. controls +(-27:0) and +(0:0)..
			(8.18,1.27).. controls +(180:0) and +(-153:0).. 
			(8.14,1.29)  -- cycle
			(6.97,1.32).. controls +(0:0.1)  and +(180:0.1)..
			(7.31,1.32).. controls  +(-6:0.1) and +(0:0.2)..
			(7.13,1.31).. controls +(180:0.2)  and +(-174:0.1)..
			(6.97,1.32)  -- cycle
			(12.78,1.38).. controls +(90:0) and  +(117:0)..
			(12.8,1.41).. controls +(-90:0)  and +(90:0)..
			(12.8,1.34).. controls +(-117:0)  and +(-90:0)..
			(12.78,1.38)  -- cycle
			(11.72,1.37).. controls +(27:0) and  +(180:0)..
			(11.77,1.37).. controls +(-45:0) and  +(0:0)..
			(11.73,1.35).. controls +(180:0)  and +(-162:0)..
			(11.72,1.37)  -- cycle
			(10.6,1.55).. controls +(27:0)  and +(135:0)..
			(10.66,1.55).. controls +(-90:0) and +(0:0).. 
			(10.62,1.54).. controls +(180:0) and  +(-153:0)..
			(10.6,1.55) -- cycle 
			 (9.68,1.71).. controls +(27:0) and +(135:0)..
			(9.73,1.71).. controls  +(-45:0) and +(-18:0)..
			(9.69,1.7).. controls +(180:0)  and +(-153:0)..
			(9.68,1.71) --  cycle
			 (9.52,1.74).. controls +(27:0) and +(180:0)..
			(9.57,1.74)..  controls +(-45:0) and +(0:0)..
			(9.54,1.72).. controls  +(180:0) and +(-153:0)..
			(9.52,1.74) --  cycle
			(4.09,1.79)..  controls +(5:0.1) and +(175:0.1)..
			(4.53,1.79)..  controls +(0:0.1) and +(0:0.2)..
			(4.31,1.78).. controls +(180:0.2) and +(180:0.1).. 
			(4.09,1.79) --  cycle
			(9.15,1.79)..  controls +(27:0) and +(135:0)..
			(9.2,1.79)..  controls +(-45:0) and +(0:0)..
			(9.17,1.77)..  controls +(166:0) and +(-153:0)..
			(9.15,1.79) --  cycle
			(3.25,1.82)..  controls +(0:0) and +(180:0)..
			(3.35,1.82).. controls +(-18:0)  and +(0:0.1)..
			 (3.29,1.81).. controls +(180:0) and +(-162:0)..
			(3.25,1.82)  -- cycle
			(5.59,1.82).. controls  +(0:0) and +(180:0.1)..
			(5.79,1.82).. controls +(-11:0.1) and +(0:0.1).. 
			(5.69,1.81).. controls +(180:0.1) and  +(-170:0.1)..
			 (5.59,1.82) -- cycle
			(8.84,1.82).. controls +(0:0)  and +(180:0)..
			(8.92,1.82).. controls +(-27:0) and  +(0:0)..
			(8.88,1.8).. controls +(180:0.1) and +(-153:0).. 
			(8.84,1.82) --  cycle
			(2.93,1.84)..  controls +(27:0) and +(153:0)..
			(3,1.84).. controls +(0:0) and  +(0:0)..
			 (2.96,1.83).. controls +(180:0) and +(180:0)..
			(2.93,1.84) -- cycle 
			(6.69,1.85).. controls +(0:0.1) and  +(180:0.1)..
			(6.97,1.85).. controls +(-8:0.1)  and +(0:0.2)..
			(6.83,1.84).. controls +(180:0.2)  and +(-173:0.1)..
			(6.69,1.85) -- cycle 
			(8.25,1.84).. controls  +(11:0.1) and +(169:0.1)..
			(8.42,1.84).. controls  +(0:0.1) and +(0:0.1)..
			 (8.33,1.83).. controls +(180:0.1) and +(180:0.1)..
			(8.25,1.84) --  cycle
			(2.75,1.87).. controls +(27:0) and  +(135:0)..
			(2.8,1.87)..  controls +(-45:0) and +(0:0)..
			(2.76,1.85).. controls  +(162:0) and +(-153:0)..
			(2.75,1.87) -- cycle 
			(1.88,2.08).. controls  +(45:0) and +(0:0)..
			(1.95,2.14)..  controls +(0:0) and +(45:0)..
			 (1.91,2.08).. controls +(-135:0.1) and +(0:0)..
			 (1.83,2.01).. controls +(180:0) and +(-135:0.1)..
			 (1.88,2.08) -- cycle
			(11.66,2.16).. controls +(27:0) and +(153:0).. 
			(11.73,2.16).. controls +(-27:0)  and +(0:0)..
			 (11.69,2.15).. controls +(180:0) and +(-153:0)..
			 (11.66,2.16) -- cycle
			 (11.95,2.16).. controls +(27:0) and +(153:0)..
			(12.02,2.16).. controls  +(-27:0) and +(0:0)..
			(11.99,2.15)..  controls +(180:0) and +(-153:0)..
			(11.95,2.16) -- cycle 
			(2.22,2.24)..  controls +(72:0) and +(108:0)..
			(2.24,2.25).. controls  +(-63:0) and +(45:0)..
			 (2.24,2.2).. controls +(-135:0) and +(-90:0)..
			(2.22,2.24) -- cycle 
			(12.63,2.41).. controls +(45:0)  and +(0:0)..
			(12.69,2.46).. controls  +(0:0) and +(45:0)..
			(12.66,2.41).. controls +(-135:0)  and +(0:0)..
			 (12.6,2.35).. controls +(180:0) and +(-124:0)..
			 (12.63,2.41) -- cycle
			(9.71,2.61).. controls +(45:0) and +(135:0).. 
			(9.76,2.61).. controls +(-90:0) and  +(-18:0)..
			(9.72,2.6).. controls +(180:0) and  +(-162:0)..
			(9.71,2.61)  -- cycle
			(6.07,2.64).. controls +(0:0.1) and  +(180:0.1)..
			(6.48,2.64).. controls +(-5:0.1) and  +(0:0.2)..
			(6.27,2.63).. controls  +(180:0.2) and +(-175:0.1)..
			(6.07,2.64) --  cycle
			(9.47,2.64).. controls  +(0:0) and +(180:0)..
			(9.53,2.64).. controls +(-27:0)  and +(0:0)..
			(9.5,2.62).. controls  +(180:0) and +(-153:0)..
			(9.47,2.64)  -- cycle
			(5.43,2.66).. controls +(18:0) and  +(153:0)..
			(5.51,2.66).. controls +(0:0) and  +(0:0)..
			 (5.46,2.65).. controls +(180:0) and +(180:0)..
			(5.43,2.66) -- cycle 
			(7.44,2.67).. controls +(0:0.1) and  +(180:0)..
			(7.64,2.67).. controls +(-11:0.1)  and +(0:0.1)..
			(7.54,2.65).. controls +(180:0.1) and  +(-169:0.1)..
			(7.44,2.67) --  cycle
			(9.12,2.66).. controls +(18:0) and  +(162:0)..
			(9.23,2.66).. controls  +(0:0) and +(0:0.1)..
			(9.17,2.65).. controls +(180:0.1) and  +(180:0)..
			(9.12,2.66) -- cycle 
			(11.95,2.72).. controls +(0:0)  and +(90:0)..
			 (12,2.71).. controls +(0:0) and +(0:0)..
			(11.97,2.7).. controls +(180:0) and +(-153:0).. 
			(11.95,2.72) -- cycle 
			(4.69,3.04).. controls +(90:0) and +(135:0).. 
			(4.7,3.08).. controls +(-63:0) and  +(63:0)..
			 (4.7,3.01).. controls +(-117:0) and +(-90:0)..
			(4.69,3.04) -- cycle 
			(10.06,3.14).. controls  +(18:0) and +(153:0)..
			(10.15,3.14).. controls +(-18:0)  and +(0:0.1)..
			 (10.11,3.13).. controls +(180:0) and +(-153:0)..
			 (10.06,3.14) -- cycle
			(9.39,3.17)..  controls +(0:0.1) and +(180:0.1)..
			(9.65,3.17).. controls +(-8:0.1) and +(0:0.2).. 
			 (9.51,3.16).. controls +(180:0.1) and +(-173:0.1)..
			(9.39,3.17) -- cycle 
			 (6.74,3.19).. controls +(2:0.3) and +(178:0.3)..
			 (7.92,3.19).. controls +(0:0.3) and +(0:0.7)..
			 (7.33,3.19).. controls +(180:0.7) and +(180:0.3)..
			(6.74,3.19) -- cycle 
			 (5.81,3.22).. controls +(27:0) and +(153:0)..
			(5.88,3.22).. controls  +(-27:0) and +(0:0)..
			(5.85,3.21).. controls +(180:0) and +(-135:0).. 
			(5.81,3.22) --  cycle
			 (5.6,3.25).. controls +(0:0) and +(90:0)..
			(5.65,3.24).. controls +(0:0) and +(0:0).. 
			(5.62,3.23).. controls  +(180:0) and +(-153:0)..
			(5.6,3.25) --  cycle
			(10.29,4.12)..  controls +(27:0) and +(135:0)..
			(10.34,4.12).. controls  +(-45:0) and +(0:0)..
			(10.3,4.1)..  controls +(180:0) and +(-153:0)..
			(10.29,4.12)  -- cycle
			(7.71,4.15).. controls  +(0:0) and +(153:0)..
			(7.79,4.14).. controls  +(0:0) and +(0:0)..
			(7.74,4.13).. controls +(180:0)  and +(-162:0)..
			(7.71,4.15)  -- cycle
			(9.1,4.15)..  controls +(0:0.1) and +(180:0.1)..
			(9.31,4.15).. controls  +(-11:0.1) and +(0:0.1)..
			(9.19,4.14).. controls +(180:0.1)  and +(-172:0.1)..
			(9.1,4.15)  -- cycle
			(9.86,4.15).. controls +(0:0.1) and +(180:0).. 
			(10.02,4.15).. controls +(-14:0) and  +(0:0.1)..
			(9.94,4.14).. controls  +(180:0.1) and +(-166:0)..
			(9.86,4.15) --  cycle
			(12.3,4.15)..  controls +(0:0) and +(135:0)..
			(12.35,4.14).. controls +(-45:0)  and +(0:0)..
			(12.31,4.13).. controls +(180:0) and +(-153:0).. 
			 (12.3,4.15) -- cycle
			(7.19,4.31)..  controls +(127:0.1) and +(143:0)..
			(7.2,4.32)..  controls +(-39:0.1) and +(0:0)..
			(7.24,4.26)..  controls +(180:0) and +(-56:0)..
			 (7.19,4.31) -- cycle
			(7.1,4.56).. controls  +(90:0.1) and +(117:0)..
			 (7.11,4.6).. controls +(-72:0) and +(63:0)..
			(7.11,4.52).. controls +(-117:0) and +(-90:0).. 
			(7.1,4.56) --  cycle
			(8.41,4.68).. controls +(0:0.1) and  +(180:0.1)..
			(8.74,4.68)..  controls +(-6:0.1) and +(0:0.2)..
			 (8.57,4.67).. controls +(180:0.2) and +(-174:0.1)..
			(8.41,4.68) --  cycle
			(10.47,4.67).. controls +(27:0) and  +(153:0)..
			(10.54,4.67).. controls  +(0:0) and +(0:0)..
			 (10.5,4.66).. controls +(180:0) and +(180:0)..
			(10.47,4.67) -- cycle 
			(7.93,4.7)..  controls +(27:0) and +(135:0)..
			 (7.98,4.7).. controls +(-45:0) and +(-18:0)..
			 (7.95,4.69).. controls +(180:0) and +(-153:0)..
			(7.93,4.7) --  cycle
			(9.64,4.7).. controls +(6:0.1)  and +(175:0.1)..
			(10.01,4.7).. controls +(0:0.1) and +(0:0.2).. 
			(9.83,4.69).. controls +(180:0.2) and +(180:0.1).. 
			(9.64,4.7) -- cycle 
			(7.32,5.19).. controls  +(45:0) and +(180:0)..
			(7.36,5.21).. controls +(0:0) and  +(90:0)..
			 (7.38,5.19).. controls +(-90:0) and +(0:0)..
			(7.34,5.17)..  controls +(180:0) and +(-135:0)..
			 (7.32,5.19) -- cycle;
		 \path[fill=c606060] (9.95,1.03).. controls +(170:0.7) and +(-6:0.2)..
			 (8.39,1.27).. controls +(174:0.4) and +(2:0.7)..
			(5.42,1.29).. controls +(-177:0.2) and  +(2:0.4)..
			(4.38,1.25) -- (3.71,1.24) -- (3.67,1.3).. controls +(127:0.1)  and +(-115:0.2)..
			(3.68,1.66) -- (3.73,1.79) -- (4.45,1.79)..  controls +(0:0.4) and +(-179:0.7)..
			 (6.46,1.82).. controls +(1:2) and +(170:1.7)..
			(10.45,1.57)..  controls +(-10:1.3) and +(-170:0.4)..
			(12.57,1.38) -- (12.77,1.42) -- (12.79,1.34).. controls +(-79:0.3) and  +(0:0.8)..
			(11.51,0.82).. controls  +(178:0.3) and +(-10:1.1)..
			 (9.95,1.03) -- cycle
			(3.13,1.26).. controls +(27:0) and  +(153:0)..
			(3.22,1.26).. controls +(0:0)  and +(0:0)..
			(3.18,1.25)..  controls +(180:0.1) and +(180:0)..
			(3.13,1.26)  -- cycle
			(3.14,1.82).. controls +(0:0)  and +(0:0)..
			(3.19,1.82).. controls +(-45:0) and +(0:0).. 
			(3.16,1.8)..  controls +(180:0) and +(-153:0)..
			(3.14,1.82)  -- cycle
			(11.62,2.17).. controls +(173:0.1) and +(-14:0.5).. 
			(10.65,2.39).. controls +(166:0.5)  and +(-12:0.1)..
			(9.66,2.63)..  controls +(170:0.1) and +(0:1)..
			(7.99,2.66).. controls  +(180:1.7) and +(-44:0.3)..
			(6.23,2.84).. controls  +(131:0.1) and +(-84:0.1)..
			(6.13,3.06) -- (6.13,3.19) -- (7.5,3.18).. controls  +(0:0.8) and +(178:0.4)..
			(9.55,3.15) -- (10.23,3.12)  -- (11.1,2.91).. controls +(-13:0.8) and +(180:0.2)..
			 (12.21,2.7).. controls +(0:0.1) and +(-166:0.1)..
			(12.59,2.73) -- (12.73,2.77) -- (12.73,2.64).. controls +(-90:0.3) and +(4:0.5).. 
			(11.93,2.16)..  controls +(-174:0.1) and +(-7:0.1)..
			(11.62,2.17) --  cycle
			(5.58,2.66).. controls +(18:0) and  +(153:0)..
			(5.68,2.66).. controls +(0:0)  and +(0:0.1)..
			(5.62,2.65).. controls +(180:0.1) and  +(180:0)..
			(5.58,2.66)  -- cycle
			(5.29,2.69).. controls +(27:0)  and +(135:0)..
			(5.34,2.69).. controls  +(-45:0) and +(0:0)..
			(5.3,2.67)..  controls +(162:0) and +(-153:0)..
			(5.29,2.69)  -- cycle
			 (4.87,3.53).. controls +(0:0) and +(-135:0.1)..
			(4.95,3.61) -- (5.04,3.69) -- (4.96,3.6).. controls +(-131:0.1) and  +(-90:0)..
			(4.87,3.53)  -- cycle
			(11.58,3.73)..  controls +(161:1.1) and +(-5:0.3)..
			(10.2,4.13).. controls +(175:0.2) and  +(2:0.5)..
			 (9.19,4.14).. controls +(-179:0.7) and +(-32:0.1)..
			(8.41,4.18).. controls +(141:0.1) and  +(-82:0.2)..
			(8.21,4.54) -- (8.19,4.67)  -- (9.26,4.68).. controls +(1:0.7) and +(172:0.2)..
			(10.56,4.66).. controls +(-8:0.2) and  +(159:0.7)..
			(11.72,4.3).. controls +(-20:0.5) and +(-161:0.2).. 
			 (12.6,4.19).. controls +(27:0) and +(90:0.1)..
			(12.65,4.07).. controls +(-90:0.1) and +(56:0.1).. 
			(12.59,3.85)..  controls +(-120:0.1) and +(6:0.2)..
			(12.11,3.61).. controls  +(-176:0.2) and +(-20:0.4)..
			 (11.58,3.73) -- cycle
			(7.84,4.15).. controls +(0:0) and +(180:0).. 
			(7.93,4.15)..  controls +(-18:0) and +(0:0.1)..
			(7.88,4.13).. controls  +(180:0) and +(-162:0)..
			 (7.84,4.15) -- cycle
			(7.59,4.17).. controls +(27:0)  and +(135:0)..
			(7.64,4.17).. controls  +(-45:0) and +(-18:0)..
			(7.6,4.16).. controls +(180:0) and +(-153:0).. 
			(7.59,4.17)  -- cycle
			(7.42,4.87) -- (7.37,4.93)  -- (7.43,4.88).. controls +(-39:0.1) and +(0:0)..
			 (7.48,4.82).. controls +(180:0) and +(-45:0)..
			(7.42,4.87) --  cycle;
		\path[fill=c494949] (2.83,1.27)..  controls +(177:0.7) and +(-67:0.2)..
			(1.72,1.54)..  controls +(114:0.1) and +(-126:0.2)..
			 (1.82,1.97).. controls +(56:0.1) and +(180:0.1)..
			(2.2,2.28)..  controls +(0:0) and +(90:0.1)..
			(2.22,2.19).. controls  +(-90:0.1) and +(155:0.1)..
			(2.4,1.95).. controls +(-24:0.2)  and +(177:0.4)..
			(3.31,1.81).. controls +(-2:0.2) and  +(90:0)..
			(3.73,1.78)..  controls +(0:0) and +(67:0.1)..
			(3.68,1.66).. controls +(-113:0.2) and  +(127:0.1)..
			(3.67,1.3) --  (3.71,1.24) -- (3.49,1.25).. controls +(180:0.1) and +(-2:0.2)..
			(2.83,1.27) --  cycle
			(5.6,2.66).. controls +(176:0.4) and  +(-55:0.1)..
			(4.75,2.89).. controls +(126:0.1) and  +(-106:0.2)..
			(4.75,3.28).. controls +(72:0.1) and  +(-132:0.1)..
			(4.93,3.57) -- (5.08,3.72) -- (5.08,3.59).. controls +(-90:0.1)  and +(140:0.1)..
			 (5.15,3.4).. controls +(-45:0.2) and +(178:0.4)..
			 (5.98,3.2) -- (6.13,3.19) -- (6.13,3.06).. controls +(-84:0.1) and +(131:0.1)..
			(6.23,2.84).. controls +(-45:0.1) and +(149:0.1).. 
			(6.42,2.7)  -- (6.52,2.65) -- (6.13,2.65).. controls +(180:0.2) and +(0:0.1)..
			(5.6,2.66) --  cycle
			(7.75,4.16).. controls +(176:0.3)  and +(-55:0.1)..
			(7.17,4.36).. controls +(128:0.2)  and +(-123:0.3)..
			(7.27,5.07) -- (7.34,5.2) -- (7.36,5.06).. controls +(-90:0.1) and +(135:0).. 
			(7.42,4.87).. controls +(-49:0.1)  and +(175:0.2)..
			(7.97,4.69)  -- (8.19,4.67) -- (8.21,4.54).. controls +(-82:0.1) and +(141:0.1)..
			 (8.41,4.18) -- (8.48,4.13) -- (8.23,4.13).. controls +(180:0.1) and +(-10:0.1)..
			 (7.75,4.16) -- cycle;
	\end{tikzpicture}
	} 
	\shortans{967} 
	 \loigiai{
		Kí hiệu $u_n$ là chiều cao so  với mực nước biển của thửa ruộng ở bậc thứ $n$.\\
		Khi  đó, dãy số $(u_n)$ là một cấp số cộng với $u_1=950$ và $d=1{,}5$.\\
		Ta có  $u_{12}=u_1+11d=950+11\cdot 1{,}5=966{,}5$.\\
		Vậy thửa ruộng  ở bậc thứ $12$ có độ cao $966{,}5$ m so với mực nước biển.
	} 
\end{ex}

% nguồn: Ok_TraLoiNgan_Lop11_Lan1
\begin{ex}%[1D2V2-7]%[ma:120ba0b6]
	Trong một hội chợ đón xuân, một gian hàng sữa muốn xếp $900$ hộp sữa theo quy luật là hàng trên cùng có 1 hộp sữa, mỗi hàng ngay phía dưới  lần lượt được xếp nhiều hơn 2 hộp so với hàng trên nó (tham khảo hình vẽ dưới). Hỏi hàng dưới cùng có bao nhiêu hộp sữa?
	\shortans{59} 
	\loigiai{ 
		Xét  cấp số cộng với số hạng đầu $u_1=1$, công sai $d=2$.\\
		Khi đó, tổng của $n$ số hạng  đầu cấp số cộng là
		$$S_n=\dfrac{n}{2}  \left[2u_1+(n-1)d\right]\Leftrightarrow 900=\dfrac{n}{2} [2\cdot 1+(n-1)\cdot 2]\Leftrightarrow 1800=2n^2 \Leftrightarrow n^2=900.$$
		 Suy ra $n=30$.\\
		Vậy số hộp sữa của  dãy cuối cùng là $u_{30}=u_1+29d=1+29\cdot2=59$.
	 }
\end{ex}

% nguồn: Ok_TraLoiNgan_Lop11_Lan1
\begin{ex}%[1D2V2-7]%[ma:bd80a68c]
	Trên một bàn cờ có nhiều ô vuông, người ta đặt $7$ hạt dẻ vào ô đầu tiên, sau đó đặt tiếp vào ô thứ hai số hạt nhiều hơn ô  thứ nhất là $5$, tiếp tục đặt vào ô thứ ba số hạt nhiều hơn ô thứ hai là $5$,... và cứ thế tiếp tục đến ô thứ $n$. Biết rằng đặt hết số ô trên bàn cờ người ta phải sử dụng $25\,450$ hạt. Hỏi bàn cờ đó có bao nhiêu ô vuông?
	 \shortans{100}
	\loigiai{ 
		Số hạt dẻ trên mỗi ô (bắt đầu từ ô thứ nhất)  theo thứ tự đó lập thành cấp số cộng $(u_n)$, có $u_1=7,d=5$. \\
		Gọi $n$ là  số ô trên bàn cờ thì $u_1+u_2+...+u_n=25\,450=S_n$.\\
		Ta  có
		\begin{align*}
			&~ 25\,450=S_n=nu_1+\dfrac{n(n-1)}{2}d=7n+\dfrac{n^2-n}{2}\cdot  5 \\
			\Leftrightarrow &~  5n^2+9n-50900=0\\
			\Leftrightarrow &~  n=100.
		\end{align*}
	 }
\end{ex}

% nguồn: Ok_TraLoiNgan_Lop11_Lan1
\begin{ex}%[1D2V2-7]%[ma:8d3dde00]
	Khán đài D của một  sân vận động có $20$ hàng ghế xếp theo hình quạt. Hàng thứ nhất có $13$ ghế, hàng thứ hai có $16$ ghế, hàng thứ ba có $19$ ghế, $\ldots$, cứ thế tiếp tục cho đến hàng cuối cùng. Số ghế ở hàng cuối cùng là
	 \shortans{741}
	 \loigiai{
		Đây là cấp số cộng với $u_1 = 13$, $d = 3$.\\ Khi đó $u_{20} = u_1(20-1)d=13\cdot  19 \cdot 3 = 741$ ghế.
	} 
\end{ex}

% nguồn: Ok_TraLoiNgan_Lop11_Lan1
\begin{ex}%[1D2V2-7]%[ma:e0a65899]
	Một bức tường trang trí có dạng hình thang cân, rộng  $2{,}4$m ở đáy và rộng $1{,}2$m ở đỉnh. Các viên gạch hình vuông có kích thước $10\, \text{cm}\times 10\, \text{cm}$ phải được đặt sao cho mỗi hàng ở phía trên chứa ít hơn một viên so với hàng ở ngay phía dưới nó. Hỏi cần bao nhiêu viên gạch hình vuông như vậy để ốp hết bức tường đó?
	 \shortans{234}
	\loigiai{Đổi $2{,}4\,\text{m}=240\,\text{cm}$, $1{,}2\,\text{m}=120\,\text{cm}$.\\ 
		 Số viên gạch ở hàng đầu tiên (ứng với đáy lớn) là $u_1=\dfrac{240}{10}=24$.\\
		Số viên gạch ở  hàng trên cùng (ứng với đáy nhỏ) là $u_n=\dfrac{120}{10}=12$.\\
		Vì mỗi  hàng ở phía trên chứa ít hơn một viên so với hàng ở ngay phía dưới nó nên ta thu được một cấp số cộng có công sai $d=-1$.\\
		Như vậy $u_n=12=u_1+(n-1) d \Leftrightarrow n=13$.\\ 
		Vậy số viên gạch  hình vuông cần thiết để ốp hết bức tường đó là\\
		 $$S_{13}=\dfrac{13\left(u_1+u_{13}\right)}{2}=234\, \text{ viên gạch.}$$
	 }
\end{ex}

% nguồn: Ok_TraLoiNgan_Lop11_Lan1
\begin{ex}%[1D2V2-7]%[ma:48307e76]
	Một gia đình cần khoan một cái giếng để lấy nước. Họ thuê một đội khoan giếng nước đến để khoan giếng nước. Biết giá của một mét khoan đầu tiên là $75000$ đồng, kể từ mét khoan thứ hai giá của mỗi mét khoan tăng lên $6000$ đồng so với giá của mét khoan trước đó. Biết cần phải khoan sâu xuống $80$ m mới có nước. Số tiền phải trả (đơn  vị triệu đồng) là bao nhiêu?
	\shortans{25} 
	\loigiai{ 
		Số tiền phải trả để khoan giếng là một cấp số cộng có số hạng đầu  $u_1=75000$; công sai $ d=6000$.\\
		Số tiền phải trả để khoan $ 80m$ là $$S_{80}=\dfrac{80}{2}\left[2u_1+79d\right]=40\left[2.75000+79\cdot6000\right]=24\,960\,000 \, \text{đồng}\approx 25  \, \text{triệu đồng}.$$
	 }
\end{ex}

% nguồn: Ok_TraLoiNgan_Lop11_Lan1
\begin{ex}%[1D2V2-7]%[ma:0776a34e]
	Ông Sơn trồng cây trên một mảnh đất hình tam giác theo quy luật: ở hàng thứ nhất có $1$ cây, ở hàng thứ hai có $2$ cây, ở hàng thứ ba có $3$ cây, $\ldots$ , ở hàng thứ $n$ có  $n$ cây. Biết rằng ông đã trồng hết $11325$ cây. Hỏi số hàng cây được trồng theo cách trên là bao nhiêu?
	 \shortans{150}
	\loigiai{ 
		 Điều kiện: $n\in\mathbb{N}^*$.\\
		Số  cây ông Sơn trồng ở các hàng lần lượt tạo thành cấp số cộng có công sai $d = 1$.\\
		Ta  có
		\begin{eqnarray*}
			&&1+  2 + 3 + \cdots + n = 11325\\
			&\Leftrightarrow&  \dfrac{(1+n)n}{2} = 11325\\
			&\Leftrightarrow&n^2 +n  - 22650 = 0\Leftrightarrow \hoac{&n = 150 \text{ (thỏa mãn)}\\&n= -151 \text{ (loại)}.}
		\end{eqnarray*}
		Vậy  ông Sơn trồng $150$ hàng cây.
	} 
\end{ex}

% nguồn: Ok_TraLoiNgan_Lop11_Lan1
\begin{ex}%[1D2V2-7]%[ma:2dab16fe]
	Một rạp xiếc có $35$ dãy ghế, dãy đầu tiên có $18$ ghế. Mỗi dãy sau có hơn  dãy trước $4$ ghế. Hỏi rạp xiếc có tất cả bao nhiêu ghế?
	 \shortans{3010}
	\loigiai{ 
		Số ghế  của mỗi dãy (bắt đầu từ dãy đầu tiên) theo thứ tự đó lập thành một cấp số cộng có $35$ số hạng với công sai $d=4$ và $u_1=18$.\\
		Tổng số  ghế là $S_{35}=u_1+u_2+\cdots+u_{35}=\dfrac{35\cdot [2\cdot 18+(35-1)\cdot 4]}{2}=3010$.
	} 
\end{ex}

% nguồn: Ok_TraLoiNgan_Lop11_Lan1
\begin{ex}%[1D2V2-7]%[ma:9d6972e2]
	Một cơ sở khoan giếng đưa ra định mức giá như sau: Giá của mét khoan đầu tiên là $100$ nghìn đồng và kể từ mét khoan thứ hai, giá của mỗi mét sau tăng thêm $30$ nghìn đồng so với giá của mét khoan ngay trước đó. Một người cần khoan một giếng sâu $20$ m để lấy nước dùng cho sinh hoạt của gia đình. Hỏi sau khi hoàn thành việc khoan giếng, gia  đình đó phải thanh toán cho cơ sở khoan giếng số tiền bao nhiêu nghìn đồng?
	\shortans{7,70} 
	\loigiai{ 
		Gọi $u_n$ là giá của mét khoan thứ  $n$, trong đó $1\leq n\leq 20$.\\
		Khi đó, $\left(u_n\right)$ là  cấp số cộng có số hạng đầu $u_1=100$ và công sai $d=30$.\\
		Số tiền mà gia đình phải thanh toán cho cơ sở khoan  giếng là\\
	$$S_{20}=u_1+u_2+\ldots+u_{20}=\dfrac{20\left(2u_1+19d\right)}{2}=\dfrac{20(2\cdot 100+19\cdot 30)}{2}=7\,700\text{ (nghìn  đồng).}$$}
\end{ex}

% nguồn: Ok_TraLoiNgan_Lop11_Lan1
\begin{ex}%[1D2V2-7]%[ma:97f2384d]
	Người ta trồng cây  theo dạng một hình tam giác như sau: hàng thứ nhất trồng $1$ cây, hàng thứ hai trồng $3$ cây, hàng thứ ba trồng $5$ cây,... cứ tiếp tục trồng như thế cho đến khi hết số cây là $6561$ cây. Tính số hàng cây đã được trồng.
	 \shortans{81}
	\loigiai{ 
	Gọi $n$ là số hàng cây  được trồng $\left(n \in \mathbb{N}^*\right)$.\\
	Ta thấy,  số cây mỗi hàng tạo thành một cấp số cộng với $u_1=1, d=2$ và $S_n=6\,561$. \\
	Theo yêu  cầu bài toán, ta có:
	$$ \dfrac{n\left[2 u_1+(n-1) \cdot d\right]}{2}=\dfrac{n[2 \cdot  1+(n-1) \cdot 2]}{2}=6\,561.$$
	 Giải phương trình trên, ta được $n=81$(hàng).
	} 
\end{ex}

% nguồn: Ok_TraLoiNgan_Lop11_Lan1
\begin{ex}%[1D2V2-7]%[ma:d23afeed]
	Người ta trồng cây theo hình tam giác, với quy luật: ở  hàng thứ nhất có $1$ cây, ở hàng thứ hai có $2$ cây, ở hàng thứ ba có $3$ cây,…ở hàng thứ $n$ có $ n$ cây. Biết rằng người ta trồng hết $4950$ cây. Hỏi số hàng cây được trồng theo cách trên là bao nhiêu?
	\shortans{99} 
	\loigiai{ 
		Cách trồng cây theo quy luật trên lập thành một cấp số cộng với số hạng đầu là  số cây ở hàng 1, $u_1=1$ và công sai $ d=1$.\\
		 Ta có
		 \allowdisplaybreaks
		\begin{eqnarray*}
			&&S_n=4950  \Leftrightarrow 4950=\dfrac{n}{2}\left[2u_1+(n-1)d\right]\\
			&\Leftrightarrow& 4950=\dfrac{n}{2}\left[2\cdot1+(n-1)\cdot 1\right]\\ 
			&\Leftrightarrow&{n^2}+n-9900=0\Leftrightarrow \hoac{&n=99\\&n=-100\,(\text{loại}).} 
		\end{eqnarray*}
		Vậy để trồng hết $4950$ cây thì  cần $99$ hàng cây.
	 }
\end{ex}

% nguồn: Ok_TraLoiNgan_Lop11_Lan1
\begin{ex}%[1D2V2-7]%[ma:fc03dc59]
	Người ta thiết kế số ghế ngồi trên khán đài một sân vận động bóng đá như sau. Hàng ghế đầu tiên gần sân bóng đá nhất có $1600$ ghế. Kể từ hàng thứ hai trở đi, mỗi hàng liền sau hơn hàng liền  trước 400 ghế. Muốn sức chứa trên khán đài có ít nhất $222000$ ghế thì cần phải thiết kế ít nhất bao nhiêu hàng ghế?
	 \shortans{150}
	 \loigiai{
		Gọi $n$ là số hàng ghế cần thiết  kế và $u_1$ số ghế ở hàng đầu tiên\\
		Theo đề ta có $u_1=1600$, $u_2=u_1+400$ và $S_n\ge 222000\Leftrightarrow  n.u_1+\dfrac{n.\left(n-1\right)d}{2}\ge 222000$\\
		$\Leftrightarrow  1600n+\dfrac{n\left(n-1\right)400}{2}\ge 222000\Leftrightarrow n^2+2800-444000 \ge 0 \Leftrightarrow n\ge 150,24$.\\
	Vậy suy ra cần thiết kế ít nhất $150$ hàng  ghế.}
\end{ex}

% nguồn: Ok_TraLoiNgan_Lop11_Lan1
\begin{ex}%[1D2V2-7]%[ma:cb82e619]
	Tham gia tổ chức sự kiện âm nhạc BORNPINK tại Việt Nam, ban tổ chức ban đầu dự kiến sẽ sắp xếp vị trí các khu ghế ngồi, với mỗi khu được chia vị trí theo quy tắc số hàng ghế sau nhiều hơn hàng liền trước $3$ ghế. Giả sử khu A bắt buộc có $15$ hàng và khu A chỉ cho phép tối đa có  $1\,000$ ghế ngồi, vậy khu A sẽ có nhiều nhất bao nhiêu ghế ngồi cho khán giả theo các quy tắc trên?
	\shortans{990} 
	 \loigiai{
		Đây  là bài toán cấp số cộng.\\
		 Với $d=3$, n là số hàng, ${u_n}$ là số ghế của hàng thứ $n$.\\
		 ${u_n}={u_1}+\left(n-1 \right)d$\\
		${S_n}=\dfrac{n}{2}\left[2{u_1}+\left( n-1 \right)d \right]$\\ 
		Do khu A bắt buộc có $15$ hàng, nên $n=15$, và tổng số ghế tối  đa $1\,000$ ghế ngồi, nên ta có\\
		${S_{15}}=\dfrac{15}{2}\left[2{u_1}+\left(  15-1 \right)3 \right]\leq 1000 \\[1mm]
		\Rightarrow  u_1 \leq 45,67$\\
		Do số ghế phải là số nguyên lớn nhất $\Rightarrow u_1=45$.  \\
		Khi này, số ghế ngồi nhiều nhất mà khu A sẽ có là ${{S}_{15}}=\dfrac{15}{2}\left[2\cdot45+\left( 15-1 \right)3 \right]=990$  (ghế ngồi).
	} 
\end{ex}

% nguồn: Ok_TraLoiNgan_Lop11_Lan1
\begin{ex}%[1D2V2-7]%[ma:6630b92b]
	Người ta trồng  $3\,003$ cây theo một hình tam giác như sau: hàng thứ nhất trồng $1$ cây, hàng thứ hai trồng $2$ cây, hàng thứ ba trồng $3$ cây, $\ldots$ Cứ như thế, số cây ở hàng sau nhiều hơn hàng kề trước $1$ cây. Hỏi có tất cả bao nhiêu hàng cây được trồng?
	\shortans{77} 
	 \loigiai{
		Ta thấy số hàng cây trên lập thành cấp số cộng $(u_n)$, với $u_1=1$ và $d=1$.\\ 
		Gọi $u_n$ với $n\in \mathbb{N}^*$ là số  hàng cây do đó ta có
		 \allowdisplaybreaks
		\begin{eqnarray*}
			&&3\,003=\dfrac{n}{2}\cdot [2\cdot  1+(n-1)\cdot 1]\\
			 &\Leftrightarrow &n(n+1)=6\,006\\
			&\Leftrightarrow  &n^2+n-6\,006=0\\
			&\Leftrightarrow &\hoac{&n=77 \text{ (thỏa  mãn)}\\&n=-78\text{ (loại).}}
		\end{eqnarray*}
		Vậy có  $77$ hàng cây.
	} 
\end{ex}

% nguồn: Ok_TraLoiNgan_Lop11_Lan1
\begin{ex}%[1D2V2-7]%[ma:9e8dd3cd]
	Tỉ lệ tăng dân số mỗi năm của  một tỉnh A từ năm $2010$ đến năm $2019$ là $0{,}4\%$. Vì thực hiện chính sách về dân số nên tỉnh A dự kiến từ năm $2020$ đến $2035$ tỉ lệ tăng dân số mỗi năm chỉ còn lại $0{,}35\%$. Theo thống kê số dân tỉnh A năm $2021$ nhiều hơn năm $2017$ là $40\,500$ người. Ước tính số dân tỉnh A vào năm $2035$ (làm tròn đến hàng nghìn) là $\overline{abcd000}$. Hỏi $\overline{abcd}$ bằng bao nhiêu?
	 \shortans{2,86}
	 \loigiai{
		 Gọi dân số tỉnh A năm $2010$ là $a$ (người, $a \in \mathbb{N}^*$).\\
		Số dân mỗi năm của tỉnh A từ năm $2010$ đến $2019$ là một cấp số nhân $\left(u_n\right)$ với số hạng đầu $u_1=a$ và  công bội $q=1+0{,}4\%=1{,}004$.\\
		 Số dân tỉnh A năm $2017$ là $u_8=u_1 \cdot q^7=a \cdot 1{,}004^7$.\\
		Số dân tỉnh A năm $2019$ là $u_{10}=u_1 \cdot q^9=a \cdot  1{,}004^9$.\\
		Số dân tỉnh A năm $2020$ là $a \cdot 1{,}004^9 \cdot 1{,}0035$.\\ 
		Số dân mỗi năm của tỉnh A từ năm $2020$ đến $2035$ là một cấp số nhân $\left(v_n\right)$ với  số hạng đầu $v_1=a \cdot 1{,}004^9 \cdot 1{,}0035$ và công bội $q'=1+0{,}35 \%=1{,}0035$.\\
		Số dân của  tỉnh A năm $2021$ là $a \cdot 1{,}004^9 \cdot 1{,}0035^2$.\\
		Vì số dân tỉnh  A năm $2021$ nhiều hơn năm $2017$ là $40500$ người nên ta có
		\[ 
		a \cdot 1{,}004^9 \cdot 1{,}0035^2-a \cdot 1{,}004^7=40\;500 \Rightarrow a  \approx 2\,610\,894.
		 \]
		Số dân tỉnh A năm $2035$ là $2\,610\,894\cdot  1{,}004^9 \cdot 1{,}0035^{16} \approx 2\,862\,000$ người.\\
		 Vậy $\overline{abcd}=2\,862$.
	 }
\end{ex}

% nguồn: Ok_TraLoiNgan_Lop11_Lan1
\begin{ex}%[1D2V2-7]%[ma:b28f03e7]
	Công ty cây xanh X trồng $ 496 $ cây hoa trong một khu vườn hình tam giác như sau: hàng thứ nhất trồng một cây hoa, kể từ hàng thứ hai trở đi số cây hoa trồng mỗi hàng nhiều hơn một cây so với hàng liền trước nó. Hỏi công ty cây xanh X trồng được bao nhiêu hàng cây trong  khu vườn hình tam giác đó.
	\shortans{31} 
	\loigiai{Gọi $ u_n $ là số cây hoa được trồng ở hàng  thứ $ n $, với $ n \in \mathbb{N}^{*} $.\\
		Vì hàng thứ nhất trồng một cây  hoa, kể từ hàng thứ hai trở đi số cây hoa trồng mỗi hàng nhiều hơn một cây so với hàng liền trước nó nên $ (u_n)$ là cấp số cộng với số hạng đầu $ u_1=1 $, công sai $ d=1 $.\\
		Tổng số cây hoa trồng được ở $ n $  hàng là
		$$ S_n=\dfrac{\left(2u_1+(n-1)d\right)n}{2}=\dfrac{\left(2+n-1\right)n}{2}=\dfrac{\left(n+1\right)n}2.$$ 
		Mà công ty cần trồng $ 496  $ cây hoa nên
	$ \dfrac{\left(n+1\right)n}2=496  \Leftrightarrow n^2+n-992=0 \Leftrightarrow n=31$ do $ (n\in \mathbb{N}^{*}).$}
\end{ex}

% nguồn: Ok_TraLoiNgan_Lop11_Lan1
\begin{ex}%[1D2V2-7]%[ma:1a4b87a7]
	Một kiến trúc sư thiết kế một hội trường với $15$ ghế ngồi ở hàng thứ nhất, $18$ ghế ngồi ở hàng thứ hai, $21$ ghế ngồi ở hàng thứ ba, và cứ như vậy (số ghế ở hàng sau nhiều hơn $3$ ghế so với số ghế ở hàng liền trước nó). Nếu muốn hội trường  đó có sức chứa ít nhất $870$ ghế ngồi thì kiến trúc sư đó phải thiết kế tối thiểu bao nhiêu hàng ghế?\\
	 \shortans{20}
	\loigiai{ 
		Số ghế ở mỗi hàng của hội trường lập thành một cấp  số cộng với số hạng đầu $u_1 = 15$ và công sai $d = 3$.\\
		Giả sử cần thiết kế tối thiếu $n$ hàng ghế để hội trường có sức chứa ít nhất $870$ ghế  ngồi.\\
		Ta có $S_n=\dfrac{n}{2}\left[2  u_1+(n-1) d\right]=\dfrac{n}{2}\left[2\cdot15+(n-1)\cdot3\right]\geq 870$.\\
		Biến đổi bất phương trình này, ta được $$n(30+3 n-3) \geq 1740  \Leftrightarrow 3n^2+27n-1740 \geq 0\Rightarrow n \ge 20.$$
		Vậy cần thiết  kế tối thiểu $20$ hàng ghế để thỏa mãn yêu cầu bài toán.
	 }
\end{ex}

% nguồn: Ok_TraLoiNgan_Lop11_Lan1
\begin{ex}%[1D2V2-7]%[ma:da1be441]
	Chu vi của một đa  giác bằng $158$ cm; các cạnh của đa giác lập thành một cấp số cộng với công sai $3$ cm. Biết cạnh dài nhất là $44$ cm. Số cạnh của đa giác đó bằng
	\shortans{4} 
	 \loigiai{
		 Gọi số cạnh đa giác là $n$ và các cạnh lần lượt là $u_1; u_2;\ldots; u_n$. \\
		Theo đề ra, ta có  $u_n=44 \Leftrightarrow u_1+(n-1)d=44 \Leftrightarrow u_1=47-3n$.
		\begin{eqnarray*}
			&&S_n=\dfrac{u_1+u_n}{2}\cdot n  \Leftrightarrow 158=\dfrac{91-3n}{2}\cdot n
			\Leftrightarrow  3n^2-91n+316=0 \Leftrightarrow \hoac{&n=4 \text{ (nhận)} \\ &n=\dfrac{79}{3} \text{ (loại).}}
		\end{eqnarray*}
		Vậy  đa giác có $4$ cạnh.
	} 
\end{ex}

% nguồn: Ok_TraLoiNgan_Lop11_Lan1
\begin{ex}%[1D2V2-7]%[ma:ed070674]
	Một rạp hát có $18$ hàng ghế xếp theo hình quạt. Hàng thứ  nhất có $16$ ghế, hàng thứ hai có $20$ ghế, hàng thứ ba có $24$ ghế, $\ldots$ cứ thế cho đến hàng cuối cùng. Hỏi tổng số ghế có trong rạp là bao nhiêu?
	\shortans{900} 
	\loigiai{ 
		Gọi số ghế ở hàng  thứ $n$ là $u_n$.\\
		Ta có $u_1=16$;  $u_2=20$, $u_3=24$.\\
		Nhận  xét dãy số $(u_n)$ là cấp số cộng có $u_1=16$, công sai $d=u_2-u_1=20-16=4$.\\
		 Khi đó
		\begin{equation*}
			 S_{18}=\dfrac{n\left[2u_1+(n-1)d\right]}{2}=\dfrac{18\cdot(2\cdot16+17\cdot4)}{2}=900.
		\end{equation*}
		Vậy  tổng số ghế trong rạp là $900$ ghế.
	} 
\end{ex}

% nguồn: Ok_TraLoiNgan_Lop11_Lan1
\begin{ex}%[1D2V2-7]%[ma:c76db75d]
	Một kiến trúc sư thiết kế một hội trường với $15$ ghế ngồi ở hàng thứ nhất, $18$ ghế ngồi ở hàng thứ hai, $21$ ghế ngồi ở hàng thứ ba, và cứ như vậy (số ghế ở hàng sau nhiều hơn $3$ ghế so với  số ghế ở hàng liền trước nó). Nếu muốn hội trường đó có sức chứa ít nhất $870$ ghế ngồi thì kiến trúc sư đó phải thiết kế tối thiểu bao nhiêu hàng ghế?
	\shortans{20} 
	\loigiai{ 
		Số ghế ở mỗi hàng của nhà hát lập thành một cấp số cộng với số hạng đầu  $u_1=15$ và công sai $d=3$. Gọi $n$ là số hàng ghế trong hội trường. \\
		Tổng các số hạng  này là
		$S_{n}=\dfrac{n}{2}\left[2u_1+(n-1)d\right]=\dfrac{n}{2}\left[2\cdot  15+(n-1)\cdot 3\right]=\dfrac{n}{2}(3n+27)\ge 870$.\\
		Do đó $3n^2+27n-1740\ge0$.\\ 
		 Giải bất phương trình bậc hai này ta được $n\le-29 \text{\ (loại)}$ và $n\ge 20$.\\
		Vậy phải thiết kế tối thiểu $20$ hàng ghế. 
	} 
\end{ex}

% nguồn: Ok_TraLoiNgan_Lop11_Lan1
\begin{ex}%[1D2V2-7]%[ma:bd127be8]
	Bác Bình muốn trồng cây cà phê trên một ngọn đồi như sau: Từ trên đỉnh đồi  trồng hàng thứ nhất $2$ cây; đi xuống hàng thứ hai $5$ cây; đi xuống hàng thứ ba $8$ cây; $\ldots$; đi xuống hàng cuối cùng dưới chân đồi trồng $26$ cây. Theo cách trồng như trên bác bình trồng được bao nhiêu cây cà phê?
	\shortans{126} 
	 \loigiai{
		Gọi $u_1;u_2;u_3; \ldots;  u_n$ lần lượt là số cây hàng thứ nhất; hàng thứ hai; hàng thứ ba; \ldots; hàng thứ $n$.
		 Ta có $u_1=2; u_2=5 \Rightarrow d=3$.\\
		Theo cách trồng của bác Bình số cây mỗi hàng lập  thành một cấp số cộng với số hạng đầu $u_1=2$ và công sai $d=3$. Thật vậy, ta có
		\begin{itemize}
			\item $u_1=  2$ cây.
			\item $u_2=u_1+d = 2+3=5$ cây. 
			\item  $u_3=u_1+2d=5+3=8$ cây.
			\item  \ldots
			\item $u_n  = u_1+(n-1)d = 26$ cây.
		\end{itemize}
		Suy ra $2+(n-1)\cdot 3 =26 \Leftrightarrow n =  9$. Vậy có $9$ hàng cây.\\
		Do đó, số cây cà phê bác Bình trồng được là 
		\[S_9  =\dfrac{9(2+26)}{2}=126 \text{ cây.} \]
	 }
\end{ex}

% nguồn: Ok_TraLoiNgan_Lop11_Lan1
\begin{ex}%[1D2V2-7]%[ma:653906c9]
	Vào đầu mỗi tháng, ông An đều gửi vào ngân hàng số tiền cố định $30$ triệu đồng theo hình thức lãi kép với lãi suất $0{,}6\%$ mỗi tháng. Tính số  tiền ông An có được sau tháng sau tháng thứ hai
	 \shortans{60,5}
	 \loigiai{
		Số tiền  ông An có được:\\
		$\bullet$ Sau tháng thứ nhất là $T_1=30+30 \cdot  \dfrac{0{,}6}{100}=30{,}18$ (triệu đồng).\\
		$\bullet$ Sau tháng thứ hai  là $T_2=30+30 \left(1+\dfrac{0{,}6}{100}\right)+\left[30+30\left(1+\dfrac{0{,}6}{100}\right)\right]\dfrac{0{,}6}{100} \approx 60{,}54$ (triệu đồng).
	} 
\end{ex}

% nguồn: Ok_TraLoiNgan_Lop11_Lan1
\begin{ex}%[1D2V2-7]%[ma:c2db0b8e]
	Một gia đình cần khoan một cái giếng để lấy nước. Họ thuê một đội khoan giếng đến để khoan giếng nước. Biết giá của một mét khoan đầu tiên là $75\,000$ nghìn đồng, kể từ mét khoan thứ hai giá mỗi mét khoan tăng lên $6\,000$ nghìn đồng so với giá của mét  khoan trước đó. Biết cần phải khoang sâu $80\mathrm{\,m}$ thì mới có nước. Vậy gia đình cần phải trả bao nhiêu tiền để khoan cái giếng đó? (kết quả làm tròn đến hàng đơn vị của triệu đồng).
	\shortans{25} 
	\loigiai{ 
		Giá tiền khoan mỗi mét giếng (bắt đầu từ mét đầu tiên) lập thành cấp số  cộng $(u_n)$ có số hạng đầu $u_1=75\,000$ và công sai $d=6\,000$.\\
		Do cần khoan $80\mathrm{\,m}$ nên tổng số tiền cần trả là\\ $$S_{80}=u_1+u_2+\ldots+u_{80}=\dfrac{80}{2}\left(2\cdot 75\,000+79\cdot 6\,000 \right)=24\,960\,000\approx  25\,\text{(triệu đồng)}.$$
	 }
\end{ex}

% nguồn: Ok_TraLoiNgan_Lop11_Lan2
\begin{ex}%[1D2V2-7]%[ma:cb156ca2]
	Chiều cao (đơn vị: centimét) của một đứa  trẻ $n$ tuổi phát triển bình thường được cho bởi công thức $x_n=75+5(n-1)$.
	Một đứa trẻ phát triển bình thường có chiều cao  năm $3$ tuổi là bao nhiêu centimét? \\
	\begin{center}
		\begin{tikzpicture}[rotate=90,x=1cm,y=0.4cm,scale=0.8]
				\draw[thick]  (0,0) -- (10,0);
				\draw[thick,fill=pink!30,draw=pink!60]  (-0.2,-0.3) rectangle (11,4);
				 \foreach \x in {0,0.5,...,10}{
					 \draw[thick,red] (\x,0) -- (\x,1.5);
				 }
				\foreach \x in {0,0.1,...,10}{ 
					\draw[thin,red] (\x,0) -- (\x,0.55); 
				 }
				\node[left,above]  at (10.2,0.8) {cm};
			\end{tikzpicture}
	\end{center}
	\shortans{85} 
	 \loigiai{
		Chiều cao $3$ năm tuổi của một đứa bé phát triển  bình thường là ${x_3}=75+5\cdot\left(3-1\right)=85~\text{cm}$.
	} 
\end{ex}

% nguồn: Ok_TraLoiNgan_Lop11_Lan2
\begin{ex}%[1D2V2-7]%[ma:e39fb49c]
	Một người muốn mua một thanh gỗ đủ để cắt ra làm các thanh ngang của một cái thang. Biết rằng chiều dài các thanh ngang của cái thang  đó (từ bậc dưới cùng) lần lượt là $45; \, 43; \, 41; \, \ldots \,; 31$ cm.
	\begin{center}
		\begin{tikzpicture}[line width=6pt,line cap=round]
				\draw (0,0) -- (0,5); 
				\draw  (2,0) -- (2,5);
				\foreach \y  in {0.5,1.5,2.5,3.5,4.5}{
					\draw[line width=5pt]  (0,\y) -- (2,\y);
				 }
			\end{tikzpicture}
	\end{center}
	Chiếc thang đó  có bao nhiêu bậc?
	 \shortans{8}
	\loigiai{ 
		Ta  có
		\[ 
		u_1=45;\, d=-2;\,  u_n=31.
		\] 
		Dùng công  thức
		\[ 
		u_n=u_1+(n-1)d \Rightarrow  31=45+(n-1)(-2) \Rightarrow (n-1)=7\Rightarrow n=8.
		 \]
		Vậy thang có $8$ bậc. 
	} 
\end{ex}

% nguồn: Ok_TraLoiNgan_Lop11_Lan2
\begin{ex}%[1D2V2-7]%[ma:ddf28e0d]
	Anh Nam được nhận vào làm việc ở một công ty về công nghệ với mức lương khởi điểm là $100$ triệu đồng một năm. Công ty sẽ tăng thêm lương cho  Anh Nam mỗi năm là $20$ triệu đồng. Tính tổng số tiền lương mà anh Nam nhận được sau $10$ năm làm việc cho công ty đó.
	 \shortans{1,90}
	 \loigiai{
		Số tiền lương anh  Nam nhận được mỗi năm lập thành một cấp số cộng gồm $10$ số hạng, với số hạng đầu $u_1=100$, công sai $d=20$.\\
		 Tổng $10$ số hạng đầu của cấp số cộng này là ${S_{20}}=\dfrac{10}{2}\left(2\cdot 100+(10-1)\cdot 20\right)=1\,900$.\\
		Vậy số tiền lương anh Nam nhận  được trong $10$ năm là $1\,900$ triệu đồng
	} 
\end{ex}

% nguồn: Ok_TraLoiNgan_Lop11_Lan2
\begin{ex}%[1D2V2-7]%[ma:8457f55f]
	Vào năm $2\,023$, nhiệt độ trung bình của  thành phố A là khoảng $29{,}5^\circ \mathrm{C}$. Giả sử do biến đổi khí hậu nên mỗi năm nhiệt độ trung bình của thành phố A đều tăng thêm khoảng $0{,}1^\circ \mathrm{C}$. Hãy ước tính kể từ năm nào thì nhiệt độ trung bình của thành phố A đạt từ $35^\circ \mathrm{C}$ trở lên.
	 \shortans{2,08}
	 \loigiai{
		Theo bài toán, nhiệt độ trung bình ở mỗi năm của  thành phố A lập thành cấp số cộng với công
		sai là $d=0{,}1\left(^\circ C\right)$ và $u_1=29{,}5^\circ \mathrm{C}$ là nhiệt độ trung bình  của thành phố A vào năm $2\,023$.\\
		Giả sử số hạng thứ $n$ của cấp số  cộng có giá trị lớn hơn hoặc bằng $35$.\\
		Tức là $u_n\ge 35$ hay  $u_1+(n-1)\cdot d\ge 35\Leftrightarrow 29{,}5+(n-1)\cdot 0{,}1\ge 35\Leftrightarrow n\ge 56$.\\
		Do đó, kể từ số hạng thứ $56$ trở đi thì chúng đều có giá trị  lớn hơn hoặc bằng $35$.\\
		Ta có $u_1$ là nhiệt độ  trung bình của thành phố A vào năm $2\,023$.\\
		Nên ${u_{55}}$ là nhiệt độ trung bình của  thành phố A vào năm $(2\,023+56-1)=2\,078$.\\
		Vậy kể từ năm $2\,078$ thì nhiệt độ trung bình của thành phố A đạt từ $35^\circ \mathrm{C}$ trở  lên.
	 }
\end{ex}

% nguồn: Ok_TraLoiNgan_Lop11_Lan2
\begin{ex}%[1D2V2-7]%[ma:673d7859]
	Hùng đang tiết kiệm để mua một cây guitar. Trong tuần đầu tiên, anh ta để dành $42$ đô la, và trong mỗi tuần tiếp theo, anh ta đã thêm $8$ đô la  vào tài khoản tiết kiệm của mình. Cây guitar Hùng cần mua có giá $400$ đô la. Hỏi vào tuần thứ bao nhiêu thì anh ấy có đủ tiền để mua cây guitar đó?
	 \shortans{46}
	\loigiai{ 
		Xét $1$ cấp số cộng  có $u_1=42$, $d=8$.\\
		 Sau $n$ tuần anh ta phải có số tiền là $u_n=42+(n-1)\cdot 8\ge 400$.
		\[n-1\ge\dfrac{400-42}{8}\Leftrightarrow  n\ge\dfrac{400-42}{8}+1=45\cdot 75.\]
		Vậy kể  cả tuần đầu thì tuần thứ $46$ anh ta có đủ tiền để mua cây guitar đó.
	} 
\end{ex}

% nguồn: Ok_TraLoiNgan_Lop11_Lan2
\begin{ex}%[1D2V2-7]%[ma:e4ff66d0]
	Vào năm $2\,020$, dân số của  một thành phố là khoảng $1{,}2$ triệu người. Giả sử mỗi năm, dân số của thành phố này tăng thêm khoảng $30$ nghìn người. Hãy ước tính dân số của thành phố này vào năm $2\,030$. (đơn vị triệu người)
	\shortans{1,5} 
	\loigiai{ 
		Ta có $1{,}2$ triệu người $=1\,200$ nghìn  người.\\
		Dân số mỗi năm của thành phố từ  năm $2\,020$ đến năm $2\,030$ lập thành một cấp số cộng, gồm $11$ số hạng $(2\,030-2\,020+1=11)$, với số hạng đầu $u_1=1\,200$ và công sai $d=30$.\\
		Ta có ${u_{11}}=u_1+10d=1\,500$.\\ 
		Vậy dân số của thành phố này vào năm $2\,030$ khoảng $1\,500$ nghìn người hay $1{,}5$ triệu  người.
	 }
\end{ex}

% nguồn: Ok_TraLoiNgan_Lop11_Lan2
\begin{ex}%[1D2V2-7]%[ma:89f797bb]
	Khi một vận động viên nhảy dù nhảy ra khỏi máy bay, giả sử quãng đường người ấy rơi tự do (tính theo feet) trong mỗi giây liên tiếp theo thứ tự trước khi bung dù lần lượt là $16; \,48;\, 80;\, 112;\, 144; \ldots$ (các quãng đường này tạo thành cấp số  cộng).
	Công sai của cấp số  cộng trên là bao nhiêu?
	\shortans{32} 
	\loigiai{ 
		Công sai của cấp số cộng  trên là $d=32$.
	} 
\end{ex}

% nguồn: Ok_TraLoiNgan_Lop11_Lan2
\begin{ex}%[1D2V2-7]%[ma:55808ba2]
	Người ta trồng $3\,003$ cây theo một hình tam giác như sau hàng thứ nhất trồng $1$ cây, hàng thứ hai trồng $2$ cây, hàng thứ ba trồng $3$  cây. Hỏi có tất cả bao nhiêu hàng cây?
	\shortans{77} 
	 \loigiai{
		Số cây mỗi hàng (bắt đầu từ hàng thứ nhất) lập thành một cấp số cộng $(u_n)$ có $u_1=1$; $d=1$. Giả sử có $n$ hàng  cây thì $u_1+u_2+\ldots+u_n=3\,003=S_n$.\\
		Ta có $3\,003=S_n=nu_1+\dfrac{n(n-1)}{2}d \Leftrightarrow n^2+n-6\,006=0 \Leftrightarrow n=77$. 
	 }
\end{ex}

% nguồn: Ok_TraLoiNgan_Lop11_Lan2
\begin{ex}%[1D2V2-7]%[ma:bada578b]
	Trên một bàn cờ có nhiều ô vuông, người ta đặt $7$ hạt dẻ vào ô đầu tiên, sau đó đặt tiếp vào ô thứ hai số hạt nhiều hơn ô thứ nhất là $5$, tiếp tục đặt vào ô thứ ba số hạt nhiều hơn ô thứ hai là $5$, \ldots\, và cứ thế tiếp tục  đến ô thứ $n$. Biết rằng đặt hết số ô trên bàn cờ người ta phải sử dụng $25\,450$ hạt. Hỏi bàn cờ đó có bao nhiêu ô vuông?
	\shortans{100} 
	 \loigiai{
		Số hạt dẻ trên mỗi ô (bắt đầu từ ô thứ nhất) theo thứ tự đó lập thành cấp số cộng $(u_n)$ có $u_1=7, d=5$. Gọi $n$ là  số ô trên bàn cờ thì $u_1+u_2+\ldots+u_n=25\,450=S_n$.\\
		Ta  có
		\begin{eqnarray*}
			&&25\,450=S_n\\ 
			&\Leftrightarrow&nu_1+\dfrac{n(n-1)}{2}d=25\,450\\ 
			 &\Leftrightarrow&7n+\dfrac{n^2-n}{2}\cdot 5=25\,450\\
			 &\Leftrightarrow&5n^2+9n-50\,900=0\\
			 &\Rightarrow&n=100.
		\end{eqnarray*} $25\,450=S_n=nu_1+\dfrac{n(n-1)}{2}d=7n+\dfrac{n^2-n}{2}\cdot 5 \Leftrightarrow 5n^2+9n-50\,900=0 \Leftrightarrow n=100$.
	 }
\end{ex}

% nguồn: Ok_TraLoiNgan_Lop11_Lan2
\begin{ex}%[1D2V2-7]%[ma:1e79273f]
	Công ty A tuyển một kĩ sư xây dựng với mức lương năm đầu là $180$ triệu đồng/năm và cam  kết sau mỗi năm, tiền lương sẽ tăng thêm $8$ triệu đồng/năm so với năm liền trước đó. Hỏi sau bao nhiêu năm thì tổng tiền lương của người kĩ sư đó bằng $2\,160$ triệu đồng.
	\shortans{10} 
	 \loigiai{
		Số tiền lương của kĩ sư lập thành một cấp số cộng với $u_1=180$ và $d=8$.\\ 
		Theo đề  bài ta có $S_n=2\,160$, với $n\in \mathbb{N^*}$.\\
		Mặt khác, $S_n=nu_1+\dfrac{n(n-1)}{2}d=180n+4n(n-1)=4n^2+176n$.\\ 
		 Như vậy
		\begin{eqnarray*}
			 S_n=2\,160&\Leftrightarrow& 4n^2+176n=2\,160\\
			 &\Leftrightarrow& n^2+44n-540=0\\
			 &\Leftrightarrow& \hoac{&n=10\\&n=-54\quad(\text{loại}).}
		\end{eqnarray*}
		Vậy sau $10$ năm tổng tiền lương của kĩ sư bằng $2\,160$  triệu đồng.
	 }
\end{ex}

% nguồn: Ok_TraLoiNgan_Lop11_Lan2
\begin{ex}%[1D2V2-7]%[ma:0f98e4c7]
	Trong dịp nghỉ lễ $02/9$ gia đình anh An cần thuê một xe Taxi để di chuyển từ thủ đô Hà Nội về thăm quê tại TP Bắc Giang. Biết giá của kilômét đầu tiên là $10\,000$ đồng, kể từ kilômét thứ $2$ giá của mỗi kilômét tăng thêm $500$ đồng so với giá của kilômét trước đó. Biết quãng đường Taxi di chuyển từ thủ đô Hà Nội về TP Bắc Giang là $80$ km. Hỏi  gia đình anh An phải trả bao nhiêu tiền cho chuyến Taxi đó? (đơn vị nghìn đồng)
	 \shortans{2,38}
	\loigiai{ 
		Số tiền  ở từng Kilomet lập thành một cấp số cộng.\\
		Số hạng đầu tiên  là $u_1=10\,000$, công sai $d=500$.\\
		Cấp số cộng  có $80$ số hạng.\\
		Tổng số tiền là $S_{80}=8\,010\,000+\dfrac{79\cdot  80\cdot 500}{2}=2\,380\,000$ đồng.
	 }
\end{ex}

% nguồn: Ok_TraLoiNgan_Lop11_Lan2
\begin{ex}%[1D2V2-7]%[ma:8c35a000]
	Chiều cao (đơn vị: centimét) của một đứa  trẻ $n$ tuổi phát triển bình thường được cho bởi công thức $x_n=75+5(n-1).$
	\begin{center}
		\begin{tikzpicture}[rotate=90,x=1cm,y=0.4cm,scale=0.8]
				\draw[thick]  (0,0) -- (10,0);
				 \draw[thick,fill=pink!30,draw=pink!60] (-0.2,-0.3) rectangle (11,4);
				 \foreach \x in {0,0.5,...,10}{
					\draw[thick,red] (\x,0) --  (\x,1.5);
				 }
				\foreach  \x in {0,0.1,...,10}{
					 \draw[thin,red] (\x,0) -- (\x,0.55);
				 }
				\node[left,above]  at (10.2,0.8) {cm};
			\end{tikzpicture}
	\end{center}
	Một đứa trẻ có chiều  cao $120$ cm thì bao nhiêu tuổi?
	\shortans{10} 
	\loigiai{ 
		Ta  có
		\allowdisplaybreaks 
		\begin{eqnarray*}
			 & &x_{n}=75+5(n-1).\\
			&  \Leftrightarrow & 120=75+5(n-1).\\
			& \Leftrightarrow &  5n=50.\\
			& \Leftrightarrow  & n=10.
		\end{eqnarray*}
		 Vậy đứa bé $10$ tuổi.
	 }
\end{ex}


%%%%%%%%%%%%%%%% TỰ LUẬN %%%%%%%%%%%%%%%%

% nguồn: Ok_TuanLuan_Toan11_Lan2
\begin{ex}%[1D2H2-7]%[ma:2daa1b17]
	Khi một vận động viên nhảy dù nhảy ra  khỏi máy bay, giả sử quãng đường người ấy rơi tự do (tính theo feet) trong mỗi giây liên tiếp theo thứ tự trước khi bung dù lần lượt là $16; 48; 80; 112; 144;\ldots $ (các quãng đường này tạo thành cấp số cộng). Tính tổng chiều dài quãng đường rơi tự do của người đó trong $10$ giây đầu tiên.
	\loigiai{ 
		Công sai của cấp số cộng trên là $d =  32$.\\
		Chiều dài quãng đường rơi tự do là  $S_{10} = \dfrac{10 \cdot [2 \cdot 16 + (10-1) \cdot 32]}{2} = 1600$.\\
		Vậy tổng chiều dài quãng đường rơi tự do của người đó trong $10$ giây đầu tiên là $1\,600$ feet. 
	} 
\end{ex}

% nguồn: Ok_TuanLuan_Toan11_Lan2
\begin{ex}%[1D2V2-7]%[ma:8bc093d9]
	Người ta trồng $465$ cây trong một khu vườn hình tam giác như sau: Hàng thứ nhất có $1$ cây, hàng thứ hai có $2$ cây, hàng thứ ba có $3$ cây, \ldots . Số hàng cây trong khu vườn là bao  nhiêu?
	\loigiai{Cách trồng $465$ cây trong một khu vườn hình tam giác như trên lập thành một cấp số cộng $\left(u_{n}\right)$ với $u_{n}$ là số cây ở hàng thứ $n$  và $u_{1}=1$; công sai $d=1$.\\
		Tổng số cây trồng được  là
		\begin{eqnarray*}
			& &S_{n}=465 \\ 
			&\Leftrightarrow  &\dfrac{n(n+1)}{2}=465\\
			 &\Leftrightarrow &\dfrac{n(n+1)}{2}=465\\
			&\Leftrightarrow &\hoac{&n=30 \text{  (nhận)}\\
			&n=-31\text{ (loại)}.} 
		\end{eqnarray*}
	Vậy số hàng cây trong khu vườn là $30$.} 
\end{ex}

% nguồn: Ok_TuanLuan_Toan11_Lan2
\begin{ex}%[1D2V2-7]%[ma:a8669d8f]
	Cho tam giác $ABC$ có ba góc $A,\, B,\, C$ theo thứ tự lập thành một cấp  số cộng. Hãy chứng minh rằng:
	 $ AC^2+AB\cdot BC=AB^2+BC^2$.
	 \loigiai{
		 Theo định lý cosin, ta có
		\begin{align*}
			\cos  A&=\dfrac{AB^2+AC^2-BC^2}{2\cdot AB\cdot AC}, \\
			 \cos B&=\dfrac{AB^2+BC^2-AC^2}{2\cdot AB\cdot BC}, \\
			\cos C&=\dfrac{BC^2+AC^2-AB^2}{2\cdot BC\cdot  AC}.
		\end{align*}
		 Theo định lý sin
		\begin{align*}
			\sin A=\dfrac{BC}{2R}, \quad 
			\sin  B=\dfrac{AC}{2R}, \quad
			\sin C=\dfrac{AB}{2R}. 
		\end{align*}
		Vì ba góc  $A, B, C$ theo thứ tự lập thành một cấp số cộng nên
		$  A+C=2B $.\\
		 Do đó
		\begin{align*}
			&\sin (A+C) = \sin  (2B) \\
			\Leftrightarrow  &\sin A \cdot\cos C + \sin C \cdot\cos A = 2 \sin B \cdot\cos B \\
			 \Leftrightarrow & \sin (A+C) = \sin (2B) \\
			\Rightarrow &\sin A \cdot\cos C + \sin C  \cdot\cos A = 2 \sin B \cdot\cos B\\
			\Leftrightarrow & \dfrac{BC}{2R} \cdot\dfrac{B{{C}^{2}}+A{{C}^{2}}-A{{B}^{2}}}{2\cdot BC \cdot AC}+\dfrac{AB}{2R} \cdot \dfrac{A{{B}^{2}}+A{{C}^{2}}-B{{C}^{2}}}{2\cdot AB\cdot AC}=2\cdot\dfrac{AC}{2R}\cdot\dfrac{A{{B}^{2}}+B{{C}^{2}}-A{{C}^{2}}}{2\cdot AB\cdot  BC}\\
			\Leftrightarrow &  1=\dfrac{A{{B}^{2}}+B{{C}^{2}}-A{{C}^{2}}}{AB\cdot BC}\\
			\Leftrightarrow  & A{{C}^{2}}+AB\cdot BC=A{{B}^{2}}+B{{C}^{2}}.
		\end{align*}
	 }
\end{ex}

% nguồn: Ok_TuanLuan_Toan11_Lan2
\begin{ex}%[1D2V2-7]%[ma:5e7ea658]
	Khi kí kết hợp đồng lao động với người lao động, một doanh nghiệp đề xuất  hai phương án trả lương như sau:\\
	\textbf{Phương án 1:} Năm thứ nhất, tiền lương là $120$ triệu. Kể từ năm thứ hai trở đi, mỗi năm tiền lương  được tăng $18$ triệu.\\
	\textbf{Phương án 2:} Quý  thứ nhất, tiền lương là $24$ triệu. Kể từ quý thứ hai trở đi, mỗi quý tiền lương được tăng $1{,}8$ triệu.\\
	Nếu là người được tuyển dụng vào  doanh nghiệp trên, em sẽ chọn phương án nào khi.
	\begin{enumerate}
		\item Kí hợp đồng lao động $3$  năm?
		\item Kí hợp đồng lao động  $10$ năm?
	\end{enumerate}
	\loigiai{ 
		\textbf{Theo phương án 1:} Gọi $(u_n)$ là dãy số tiền lương của người lao động  theo phương án $1$ qua mỗi năm.\\
		Dãy số $(u_n)$ lập thành một cấp số cộng có số hạng đầu $u_1 = 120$ và công sai $d  = 18$.\\
		Khi đó số hạng tổng quát của cấp số nhân là $u_n =  120 + (n-1)\cdot 18$.\\
		\textbf{Theo phương án 2:} Gọi $(v_n)$ là dãy số tiền  lương của người lao động theo phương án $2$ qua từng quý.\\
		Dãy số $(v_n)$ lập thành một cấp số cộng có số hạng đầu $v_1 = 24$ và công sai $d  = 1{,}8$.\\
		Khi đó số hạng tổng quát  của cấp số nhân là $v_n = 24 + (n-1)\cdot 1{,}8$.\\
		a) Khi kí hợp đồng $3$ năm tương đương với  $12$ quý ta có\\
		\textbf{Theo phương  án 1:} $u_3 = 120 + (3 - 1)\cdot 18 = 156$ (triệu đồng).\\
		Tổng số tiền lương nhận được  sau $3$ năm là $S_3 = \dfrac{3\cdot (120+156)}{2} = 414$ (triệu đồng).\\
		\textbf{Theo phương  án 2:} $v_{12} = 24 + (12-1)\cdot 1{,}8 = 43{,}8$.\\
		Tổng số tiền lương  nhận được sau $3$ năm tương ứng với $12$ quý là $S_{12} = \dfrac{12\cdot (24+43{,}8)}{2} = 406{,}8$ (triệu đồng).\\
		Vậy nếu được tuyển dụng vào doanh nghiệp và kí  hợp đồng lao động $3$ năm thì nên theo phương án $1$.\\
		b) Khi kí hợp đồng $10$ năm tương  đương với $40$ quý ta có\\
		\textbf{Theo phương án 1:} $u_{10} = 120 + (10 - 1)\cdot 18 = 282$ (triệu  đồng).\\
		Tổng số tiền lương nhận được sau 10 năm là $S_{10} = \dfrac{10\cdot (120+282)}{2}  = 2010$ (triệu đồng).\\
		\textbf{Theo phương án 2:}  $v_{40} = 24 + (40-1)\cdot 1{,}8 = 94{,}2$ (triệu đồng).\\
		Tổng  số tiền lương nhận được sau $10$ năm tương ứng với $40$ quý là $S_{40} = \dfrac{40\cdot (24+94{,}2)}{2} = 2364$ (triệu đồng).\\
		Vậy nếu được tuyển dụng vào doanh nghiệp và kí hợp đồng lao động $10$ năm thì nên theo phương án $2$. 
	 }
\end{ex}

% nguồn: Ok_TuanLuan_Toan11_Lan2
\begin{ex}%[1D2V2-7]%[ma:511013df]
	Trong hội chợ tết Mậu Tuất 2018, một công ty sữa muốn  xếp $900$ hộp sữa theo số lượng $1$, $3$, $5$,\ldots từ trên xuống dưới (số hộp sữa trên mỗi hàng xếp từ trên xuống là các số lẻ liên tiếp). Hỏi hàng dưới cùng có bao nhiêu hộp sữa?
	\loigiai{\textbf{Cách 1:} Áp dụng công thức tính tổng $n$ số hạng liên tiếp của  CSC ta có
		\begin{align*}
			S_{n}=\dfrac{n}{2}\left[2 u_{1}+(n-1) d\right]  \Leftrightarrow 900=\dfrac{n}{2}[2 \cdot 1+(n-1) \cdot 2] \Leftrightarrow n^{2}=900 \Rightarrow n=30.
		\end{align*}
		Vậy $u_{30}=1+29 \cdot 2=59$.\\ 
		\textbf{Cách 2:} Áp dụng công thức $1+3+5+\cdots+(2 n-1)=n^{2}$, suy ra $n=30$.\\ 
	 Vậy $2 n-1=59$.}
\end{ex}

% nguồn: Ok_TuanLuan_Toan11_Lan2
\begin{ex}%[1D2V2-7]%[ma:e46cb8ff]
	Một gia đình cần khoan một cái giếng để lấy nước. Họ thuê một đội khoan giếng nước. Biết giá của mét khoan đầu tiên là $80\,000$  đồng, kể từ mét khoan thứ hai giá của mỗi mét khoan tăng thêm $5\,000$ đồng so với giá của mét khoan trước đó. Biết cần phải khoan sâu xuống $50$ m mới có nước. Hỏi phải trả bao nhiêu tiền để khoan cái giếng đó?
	\loigiai{Áp dụng công thức tính  tổng của $n$ số hạng đầu của cấp số cộng có số hạng đầu $u_{1}=80\,000$, công sai $d=5\,000$ ta được số tiền phải trả khi khoan đến mét thứ $n$ là
		\begin{align*}
			S_{n}=\dfrac{n\left(u_{1}+u_{n}\right)}{2}=\dfrac{n\left[2 u_{1}+(n-1) d\right]}{2}. 
		\end{align*}
		Khi  khoan đến mét thứ $50$, số tiền phải trả là
		\begin{align*}
			S_{50}=\dfrac{50[2\cdot 80\,000+(50-1) \cdot 5\,000]}{2}=10\,125\,000 \text { đồng.} 
	\end{align*}}
\end{ex}

% nguồn: Ok_TuanLuan_Toan11_Lan2
\begin{ex}%[1D2V2-7]%[ma:334d79ad]
	Một gia đình cần khoan  một cái giếng để lấy nước. Họ thuê một đội khoan giếng nước. Biết giá của mét khoan đầu tiên là $80.000$ đồng, kể từ mét khoan thứ hai giá của mỗi mét khoan tăng thêm $5.000$ đồng so với giá của mét khoan trước đó. Biết cần phải khoan sâu xuống $50$ m mới có nước. Hỏi phải trả bao nhiêu tiền để khoan cái giếng đó?
	 \loigiai{
		Áp dụng công thức tính tổng của $n$ số hạng đầu của cấp số nhân có số hạng đầu $u_1 = 80.000$, công sai $d = 5.000$ ta được số tiền phải trả khi khoan đến mét thứ $n$ là 
		$$S_n = \dfrac{n\cdot (u_1+u_n)}{2}  = \dfrac{n\cdot [2u_1+(n-1)\cdot d]}{2}$$
		Khi khoan đến mét thứ $50$ số tiền phải trả là 
		$$S_{50} = \dfrac{50\cdot[2\cdot 80000+(50-1)\cdot 5000]}{2} = 10.125.000 \text{ đồng}.$$ 
	 }
\end{ex}

% nguồn: Ok_TuanLuan_Toan11_Lan2
\begin{ex}%[1D2V2-7]%[ma:70218635]
	Chu vi một  đa giác là $158$ cm, số đo các cạnh của nó lập thành một cấp số cộng với công sai $d=3$ cm. Biết cạnh lớn nhất là $44$ cm. Số cạnh của đa giác đó là?
	\loigiai{Giả sử đa giác đã cho có $n$ cạnh thì chu vi của đa giác là $S_{n}=\dfrac{\left(u_{1}+u_{n}\right) n}{2}$ với $u_{1}$ là  cạnh nhỏ nhất.\\
		Suy ra $158=\dfrac{\left(u_{1}+44\right) n}{2} \Leftrightarrow 316=\left(u_{1}+44\right) n \Leftrightarrow 2^{2} \cdot  79=\left(u_{1}+44\right) n$.\\
		Do đó $u_{1}+44$ là ước nguyên dương của $316=2^{2}\cdot 79$  và đa giác có ít nhất ba cạnh nên
		$\dfrac{316}{3}>u_{1}+44>44$. Suy ra $u_{1}+44=79 \Leftrightarrow  u_{1}=35$.\\
	Số cạnh của đa giác đã cho là $\dfrac{44-35}{3}+1=4$  (cạnh).}
\end{ex}
